\chapter{第五次“围剿”与反“围剿”的战争}

\section{运动与攻坚}

\subsection{黎川失守}


第四次反“围剿”胜利后，中央苏区地域进一步扩大。1933年7月21日，苏维埃中央决定划中央苏区为4省40余县，即江西省的瑞金、于都、兴国、永丰、新淦、宜黄、乐安、崇仁、广昌、南丰、石城、万泰、胜利、杨殷、公略、洛口、赤水、龙冈、长胜、太雷、博生，\footnote{万泰在今万安、泰和交界处，胜利在今兴国、于都交界处，杨殷在今兴国、赣县交界处，公略在今吉安、吉水交界处，洛口、赤水在今广昌境内，龙冈、长胜、太雷分别在今永丰、瑞金、石城境内，博生为今宁都县。}福建省的长汀、上杭、武平、永定、连城、龙岩、新泉、宁化、清流、汀东、兆征、彭湃、代英，\footnote{代英、兆征均在今长汀境内，彭湃在今宁化境内，代英在今上杭、永定交界处。}粤赣省的会昌、寻乌、安远、西江、门岭、信康，\footnote{西江、门岭分别在今于都、会昌境内，信康在今信丰、南康交界处。}闽赣省的黎川、金溪、资溪、建宁、泰宁、光泽。以赣闽边界为中心，中央苏区由东向西形成分别向南平、漳州、粤北、赣州、吉安、樟树、临川辐射的广大区域。


中央苏区的壮大，尤其是控制区域的北移，使其与中央苏区控制区域继续北移，与闽浙赣苏区渐成连接之势。而黎川作为闽赣省委、省政府、省军区所在地，突出于中央苏区东北部，是中央苏区与闽浙赣苏区、赣东南与闽西北连接的主要孔道，具有重要的战略意义。早在1933年3月第四次反“围剿”时，中共中央就指示：“黎川、泰宁、建宁、广昌……这个区域是中央苏区战略的锁钥，是永远不能放弃的，这些城市仍须大大的巩固起来，特别应该注意广昌。”\footnote{《中央对目前作战计划与任务的指示》，《中共中央文件选集》第9册，第105页。}1933年9月，当国民党军由赣江自西南向东北基本完成对中央苏区的堡垒封锁后，其试探性进攻首先从这里发起。


当国民党军进攻黎川时，红军的防御兵力异常薄弱。1933年6月，根据时在上海的共产国际军事总顾问弗雷德的意见，红军决定展开分离作战，所谓“两个拳头打人”，红一军团组成中央军，部署在宜黄、乐安沿抚河一带，防守苏区北大门，第三军团及第十九师等部组成东方军，入闽作战，向闽西北方向进攻。在初期的作战计划中，红军指挥部对黎川的防御给予了相当重视，要求：“一、五军团依计划在北面地带，积极活动，隐蔽我军企图时应派独立第一团领导永兴桥地带部队前出新丰街南广部队，牵制敌第六师，使该敌于我十九师东出泰将后，不敢侵入黎川。”\footnote{《苏区中央局关于执行中央作战计划指示的布置（1933年6月13日）》，《中共中央文件选集》第9册，第232～233页。}随着福建方面战事的展开，东方军作战并不顺利，红军不断向福建增兵，红五军团十三师及原留置赣东北兵力陆续调往东方军协同作战，黎川一带兵力空虚，成为国民党军发动进攻的一个理想突破口。中革军委指出：“蒋介石的主要目的，是由南城向东南突击，以便割断江西的东北部与基本苏区的联系，并完成吉水、永丰、乐安、南丰，以达到邵武地区的坚固阵地的封锁线。”\footnote{《中革军委关于十月中战役问题致师以上首长及司令部的一封信（1933年11月20日）》，《中央苏区第五次反“围剿”》（上），第107页。}


在与红军多次交手后，蒋介石策动新一轮“围剿”时力图从战略上争取主动。他在庐山训练时就谈道：

\begin{quote}

现在土匪的主力是向闽北移动了，我们现在如果也跟到福建去找他来打，这样就是我们跟了土匪走，而陷于被动了！然则要如何才是立于主动呢？比方土匪现在在福建攻延平，攻了半个月还攻不下来，当然他向东不能发展了。如果我们现在有一个部队收复了建宁，或者是到广昌，那么，在福建的土匪，就断绝了后路，失掉根据地，不得不回来找我们打了！\footnote{蒋介石：《主动的精义与方法》，《庐山训练集》，第193～194页。}
\end{quote}


9月初，他又致电福建方面：“围攻延平之匪，必趁北路军未进剿时，先乘隙打破我在闽之主力，然后回师西向，再应付北路军，其计甚狡。但其在闽北延平一带，亦不能徘徊过久，此时只要延平各城能固守半月，一待北路军发动，则匪必西退，回顾老巢。故此时我闽军主力，应先待其向西撤回时而邀击之，不必正面急急进援，免遭暗算。”\footnote{《蒋中正电蒋光鼐蔡廷锴我闽军主力应待其西撤时腰击之不必正面进援（1933年9月7日）》，蒋中正文物档案020200019089。}这显然是希望闽方能将红军东方军拖住，以便其在赣东北从容部署。后来国民党将领也认为，第五次“围剿”中，国民党方面战术主要内容就是：“主动地选择苏区有战略意义而又便于用兵的地区为一战役阶段。集中强大的兵力，作周到的准备，进行有限目标的攻击；攻占后即利用有利地形构筑碉堡封锁线，对苏区严加戒备阻绝封锁……战略上取攻势，战术上取守势，自由主宰战局，不冒风险。”\footnote{杨伯涛：《蒋军对中央苏区第五次围攻纪要》，《文史资料选辑》第45辑，第185页。}蒋介石在日记中将这一战术概括为：“凡持重部队皆令其作前方之守备。机动者，集中要点，专备出击之用。先以一部突出，为匪之目标，以求决战。其余则分为二、三部队，在两翼后作沉机观变之备。”\footnote{《蒋介石日记》，1933年8月18日。}进攻黎川，和这一作战总思路是吻合的。同时，由于当时闽方内部反蒋活动暗潮汹涌，蒋介石对此已有知悉，进攻与福建相邻的黎川，还不无对此预作准备之意。


9月1日，蒋介石电令熊式辉等：“南浔路与九江部队可以尽量减少，并密令周浑元军第六师在内准备九月底占领黎川城为要。”\footnote{《蒋中正电示熊式辉、贺国光江西省金溪九江南昌等地军事部署（1933年9月19日）》，蒋中正文物档案002010200093004。}中旬，国民党军向黎川集结。19日，蒋介石致电熊式辉等，要求将第八十五师拨归周浑元指挥，“确实占领黎川”。\footnote{《蒋中正电熊式辉贺国光第八十五师归周浑元指挥周纵队确实占领黎川（1933年9月19日）》，蒋中正文物档案020200019093。}9月24日，国民党军周浑元纵队（下辖第五、六、九十三3师）在距黎川约20公里的硝石完成集结，次日向黎川发动进攻。当时，红军在黎川的力量薄弱。关于黎川留守部队的人数，各方面有不同说法。肖劲光回忆当时“只有一个七十人的教导队和一些地方游击队”。\footnote{《肖劲光回忆录》，解放军出版社，1987，第133页。}彭德怀回忆：“保卫黎川的是一个五六百人的独立团”。\footnote{《彭德怀自述》，第185页。}国民党方面战史记载，国民党军占领黎川后，“分窜黎川境内之匪约两千左右”。\footnote{《赣粤闽湘鄂北路剿匪军第三路军五次进剿战史》（上），第二章，第7页。}这三个数字不完全相同，考虑到黎川当时以地方部队为主的实况，大致可判断黎川中共方面部队在千人左右。后来追究肖劲光责任时，提到肖指挥下“有一师兵力足以拒止至少可迟阻敌人于硝资之线”，\footnote{周恩来：《反对红军中以肖劲光为代表的罗明路线》，《斗争》第38期，1933年12月12日。}由于曾驻赣东北一带部队除第三军团全部入闽外，第十三、十九、二十师全部及二十一师大部也开往福建，黎川守军薄弱是不争的事实，“一师兵力”的说法显然是夸大之辞。鉴于双方力量悬殊，刚刚由福建前线奉命赶回黎川的闽赣军区政治委员肖劲光为避免无谓损失，率领教导队和游击队于27日撤离黎川。据肖劲光回忆：“从我接受任务回黎川到撤出黎川这八天的时间里，军委没有给我下达坚守黎川的指示，也没有明确的撤离黎川的指示。”\footnote{《肖劲光回忆录》，第134页。}对战略要地疏于防守，中革军委一开始就犯了难以原谅的错误。


其实，对赣东北的严峻局面，中革军委及红一方面军总司令部的朱德、周恩来事先是有所预料的。早在9月3日，朱德、周恩来在注意到国民党军有南下企图后，就致电中共中央局，部署让东方军“准备回师”。\footnote{《周恩来年谱1898～1949》，第251页。}因为十九路军根据蒋介石指示，坚守不出，红军难觅有利战机，继续屯兵福建整体战局将陷不利。12日，项英致电朱、周，仍强调应“速将顺、将攻下”，同时指示：“你们必须注意蒋之行动，如赣敌有配合闽敌之企图，应有准备的以便迅速转移，消灭北线主要敌人。”\footnote{项英：《东方军必须注意蒋之行动（1933年9月12日）》，《项英军事文选》，第165页。}13日，周恩来再电博古、项英等，要求东方军结束当前战斗“迅速北上”。\footnote{《周恩来军事活动纪事》（上），中央文献出版社，2000，第195页。}中革军委复电虽同意回兵北上，但又判断：“蒋贼仍极力构造永、乐方面之封锁线，刻未东移，容我东方军迟于二十日若干时间再行北上……这样的时间和空间的条件，还容我们采用各个击破的手段，先打闽敌，以便将来独立对赣敌作战。”\footnote{项英：《各个击破，先打闽敌，尔后对赣作战（1933年9月14日）》，《项英军事文选》，第168页。}中旬末梢，中革军委实际已在做回师准备，要求红一方面军各军团向赣东方向“迅速派出有力的侦察队”，\footnote{项英：《保证机动迅速派出侦察队（1933年9月18日）》，《项英军事文选》，第173页。}准备大部队移动。同时拟调一军团往赣东，集中红军主力“消灭蒋贼东行的五个师”。\footnote{项英：《红一军团适时参加抚河以东的会战（1933年9月22日）》，《项英军事文选》，第178页。}博古在给共产国际电报中也提出：“占领顺昌和将乐后，我们将向东前进，迎击蒋介石军队。”\footnote{《秦邦宪给共产国际执行委员会远东局的电报（1933年9月23日）》，《共产国际、联共（布）与中国革命档案资料丛书》第13册，第500页。}但是，由于中革军委对东方军在福建的战果心有不甘，回师问题迟迟不能决断，而将红一军团东调红军高层也存在不同意见。项英曾致电周恩来，专门解释红一军团东移问题，指出：

\begin{quote}

此时一军团在现地挺进活动，破坏敌封锁线，固然抑留蒋贼主力，但须估量该敌在其开始大举进攻和发现我东方军北上时，可能由其远后方移兵东向。我一军团……秘密速到宜、南的机动地位，东则可以参加主力作战，北则可扼敌增援队之左背，还望加考虑处置。这一问题的本身，就在决定和运用钳制方面的兵力，而非如你所说将又成只手打人的问题。以次要方向积极的作战和伪装，原可以极小兵力争取决战所需短时钳制作用，特别紧要，也是在敌我兵力悬殊情形下抽出优势兵力决战之法。当然，不可说平分兵力于突击与钳制方面，才算是不是只手打人。\footnote{项英：《对红一军团作战行动的说明（1933年9月25日）》，《项英军事文选》，第181页。}
\end{quote}


24、25日，鉴于赣东北形势已十分紧张，周恩来连电项英，要求同意东方军“赶早北上”。\footnote{《恩来关于中央军东移、东方军北上的请示（1933年9月25日）》，《中央苏区第五次反“围剿”》（上），第70页。}25日，中革军委作出决断：“第一方面军应即结束东方战线，集中泰宁出其西北地带，消灭进逼黎川之赣敌。”\footnote{《军委关于消灭进逼黎川之敌的作战部署》，《中共中央文件选集》第10册，第359页。}从这一决策过程看，中革军委一段时间内对久攻不下的顺昌、将乐弃之不甘，冀望能有所得。其实，此时闽西对中共已成鸡肋，即使拿下顺、将，对全局也无影响，而黎川的重要性显然要超过上述两地，在已承认北上必要的大前提下，仍然抑留东方军于闽西北地区，丧失了宝贵的先机，犯下当断不断的错误。中革军委稍后自承“我们在原则上正确下了决心而犯了一个重大的错误（在决心本身上有之，而在移动兵力上则更大些），这错误就是没有正确地计上时间”。\footnote{《中革军委关于十月中战役问题致师以上首长及司令部的一封信（1933年11月20日）》，《中央苏区第五次反“围剿”》（上），第109页。}应该说，这一检讨部分道出了实情。


不过，中革军委此时不应有的迟延可能还和1933年年中围绕是否进军福建的争论有关，甚至和中共内部的政治、军事权力格局也不无关系。第四次反“围剿”胜利后，周恩来在中共军事、政治决策圈中的影响进一步增强，他不仅作为出色的政治活动家，同时作为优秀的军事指挥者被人们所认识。但是，周恩来作为所谓“老干部派”并不被中共正在执掌实权的年轻一代领导人所完全接纳。项英被安排为中革军委代主席，相当程度上就是希望借此某种程度上限制周恩来的影响。1933年6月上海方面提出组建东方军计划后，周恩来提出不同意见，而事态的发展又证明其所虑不为无因，这也使作出决策的中共年轻领导人可能不无尴尬，所以，力图在福建方面获得更大战果，证明东方军出击的正确，成为他们的一大心愿。从这个角度看中革军委明显有悖军事常识的决策，就不会显得那么突兀。


黎川易手，对此后战局影响甚大。国民党方面认为：“我军占领黎川之后，不惟伪东西两军为之隔断，而赣东北与赣南匪区之联系，亦成支解分离之象……故黎川之克服，实为我军战略上之一大收获，为五次围剿胜利之先声。”\footnote{《赣粤闽湘鄂北路剿匪军第三路军五次进剿战史》（上），第二章，第1页。}占领黎川后，国民党军还可趁机把赣东北一线一直未构成的封锁线完成，最终形成江西方面对中央苏区的完全包围。针对黎川失守后战略态势的变化，周恩来9月28日向项英和中央局连发三电，认为黎川之失，“将使赣敌得先筑据点以守”，使国民党军占得先机，在此敌情下，红军“必须以极大机动性处置当前战斗。正面迎敌或强攻黎川都处不利”。\footnote{《周恩来年谱1898～1949》，第252页。}要求主动出击，力图扳回已经失去的战略先机。


黎川之失，红军虽然出师不利，但远未到攸关大局的程度，在东方军北返，同时中央军也向黎川一带靠近后，国共间真正的交锋才刚刚开始。

\subsection{洵口、硝石之役}


黎川失守后，中革军委决定将东方军从福建撤回，中央军向黎川一线靠拢，准备在此对国民党军实施打击，相机夺回黎川。红一方面军指挥部根据这一计划发出命令：“首先消灭进逼黎川之敌，进而会合我抚西力量，全力与敌在抚河会战。”\footnote{《方面军关于歼灭黎川之敌后在抚河会战给各兵团的作战命令（1933年9月27日）》。}计划规定东方军（时辖红三、五、七军团）的基本任务是攻击国民党军左翼薛岳、周浑元两纵队，威胁南城。中央军（红一、九军团）先东移配合东方军作战，然后经康都西移棠荫、里塔等地，完成钳制抚河西岸吴奇伟纵队、粉碎抚河沿岸及抚河西敌军的基本任务。中革军委希望，红军以“消灭硝石、资溪桥、黎川地区敌人”为目的，在运动中“突击敌人之暴露翼侧”，以此打击敌有生力量，“造成对敌人中心根据地的威胁”。\footnote{《对作战部署的补充指示（1933年9月28日）》、《战地工委就近指挥突击敌人（1933年9月29日）》、《关于东方军等部基本任务和作战动作的决定（1933年9月30日）》，《项英军事文选》，第185、187、192页。}


从当时来看，中共方面总的思路是希望将部队顶到苏区外线作战，尤其要尽力打破国民党军封锁线，在战略上争取主动。这一思路其实和红军第四次反“围剿”作战思路是一脉相承的。朱德在总结第四次反“围剿”胜利原因时曾谈道：“此次战略的不同点是在择其主力，不待其合击，亦不许其深入苏区，而亦得到伟大胜利。”\footnote{中共中央文献研究室编《朱德年谱（新编本）》（上），中央文献出版社，2006，第331页。该文收入《朱德选集》、《朱德军事文选》，但均不见这一段。}面对和第四次反“围剿”有一定相似性的局面，中共继续采用这一战法应属正常。为达成顶出去打的目标，夺取南城与黎川之间的战略要点硝石具有重要意义，既可切断黎川国民党军与南城附近主力军的联系，使黎川成为孤城，从而一鼓而下，又可吸引南城一带国民党军主力出动，达到在此一带集中主力围点打援的目的。因此，中革军委一改黎川失陷之前的拖延态度，连电要求东方军迅速回师，指斥：“彭、滕又要围攻邵武，忽视上级命令或将延误军机。战机紧迫，对于命令不容任何迟疑或更改。”\footnote{项英：《对东方军中央军作战行动的指示（1933年10月2日）》，《项英军事文选》，第197页。}朱德、周恩来10月3日致电项英，明确谈道：“目前关键在中央军能以极大机动抓紧当前敌情变化，适时出现，东方军集结最大兵力以最大速度赶在抚西敌援未至硝、黎，工事未固之先，进击硝、黎敌人。”同时，对在硝石作战的风险也有充分估计：“估计敌知我军回师，如更知我中央军东移，其罗纵队有以两三师，或改由李九师、八十七师秘密移至南城，准备以硝石许师诱我，以便从南城、新丰、黎川三方面向我出击之极大可能。”\footnote{《朱、周关于东方军、中央军作战意图的电报（1933年10月3日）》，《闽赣苏区文件资料选编》，第108页；《朱德、周恩来致项英电（1933年10月3日）》，《中央苏区第五次反“围剿”》（上），第75页。}在硝石作战关键是要争取时间，尽可能早地控制硝石，否则就有落入对方包围圈的可能。事实上，朱、周的电报已经侧面提示了硝石作战的隐患，即进入对手后方作战，一方面应有强大的兵力作为后盾，另方面还要随时防止被断后路，而以当时两军兵力对比，红军在这两点上都难有成算。


在中共方面积极备战时，国民党军占领黎川后也四下出击，“清剿”分散活动的红军游击部队，巩固其对黎川的占领。10月3日，国民党军得到报告，洵口方向有中共军队千余人活动，周浑元判断其为红军独立团或游击队，于4日下令第六师周嵒部派出三团“前往扑击”。\footnote{《国民党军第6师周嵒部与红军在江西黎川洵口战斗详报》，《中华民国史档案资料汇编》第5辑第1编军事（4），第1页。}


当晚，国民党军第六师第十八旅旅长葛钟山率第三十一、三十四两团及第五师第二十七团向洵口方向进发，5日晨到达洵口后，并未发现红军大队踪影。此后，该部在回黎川还是继续搜索上与指挥部往复协商，一直到6日，侦知飞鸢方向有红军部队后，葛钟山决定“乘机袭扑该匪”，\footnote{《赣粤闽湘鄂北路剿匪军第三路军五次进剿战史》（上），第二章，第9页。}率部向飞鸢进击。


东方军9月底由福建回师，10月5日，奉命向硝石方向前进。6日，红三军团隐蔽到达黎川东北，突出于国民党军的葛钟山部成为红军向前挺进中的顺势攻击对象。6日下午，葛旅大部到达飞鸢后，红军主力突然向其发动攻击，将葛部压回洵口，同时资溪桥、湖坊一带红军自北、南两面向洵口出击，形成合围之势。7日凌晨，东方军下达攻击洵口命令，以第四师、第二十师、第五师对洵口发起进攻。国民党军被困的第二十七、三十一两团在赶来增援的部队接应下，突出包围，随增援部队撤回黎川，第三十四团被围困在洵口村内。据彭德怀回忆：被围国民党军“据守山顶土寨子，坡度很陡，不易爬上去，上面无水，再有一天半天时间，即可消灭”。\footnote{《彭德怀自述》，第185页。}但红军根据既定计划，主力前出进攻硝石，未全力消灭洵口被围部队，10日，国民党军援军到达洵口，将第三十四团残余部队接应回黎川。11月，蒋介石亲自出席仪式，为第三十四团颁授荣誉旗，称誉其在红军包围中，“危困至五日之久，卒能以寡击众，打退土匪，安全回到黎川，使我们全体剿匪军队的精神为之一振！”\footnote{蒋介石：《完成安内攘外的使命》，《先总统蒋公思想言论总集》第11卷，第600页。}是役，红军对敌3个团形成重击，缴获机枪29挺，长短枪1084支，无线电台1架。东方军伤亡700多人，国民党军阵亡458人、受伤810人、旅长葛钟山以下被俘1100余人，\footnote{《第36军洵口附近战役伤亡表》，《赣粤闽湘鄂北路剿匪军第三路军五次进剿战史》（上），第二章，第22页。}洵口死尸横陈，数日后仍“臭气大张”。\footnote{《陈伯钧日记（1933～1937）》，第107页。}洵口之役，“算是第五次反‘围剿’中一个意外的序战胜利”。\footnote{《彭德怀自述》，第185页。}不过，红军在是役中损失大批弹药，按中革军委事后保守的总结：“在我们的条件之下，战斗的胜利不是占领地方，而是消灭敌人的有生力量及夺取其器材。在这一次我们无论在哪一方面，也没有得到应有的程度。因此，以我们的损失与胜利来比较，那我们所付的代价是过份大了。”\footnote{《中革军委关于十月中战役问题致师以上首长及司令部的一封信（1933年11月20日）》，《中央苏区第五次反“围剿”》（上），第112页。}


洵口之战的胜利，对红军而言，只不过是接收了一个送上门来的礼物，根据中革军委的命令，东方军北上的目标是“消灭硝石、资溪桥以及黎川附近之敌”。\footnote{《东方军要消灭硝石等地敌人（1933年10月2日）》，《项英军事文选》，第196页。}10月7日，在得悉南城国民党军第十四、第九十四师企图增援硝石，硝石敌第二十四师将南援黎川后，决定以钳制部队“立即进攻黎川”，主力则迅速向北挺进，“占据硝石及由北侧击由南城前来之任何敌人纵队”。\footnote{《军委对我军作战部署的指示（1933年10月7日）》，《中央苏区第五次反“围剿”》（上），第87页。}虽然，这一带国民党军已构筑了野战工事，在兵力上也居优势地位，但中革军委仍然希望以迅速、突然的进攻态势抑制对方的攻击，尽力将战场突进到苏区之外。当然，在此期间，中革军委也强调，各部应遵循下列作战原则：“甲、不应攻击任何工事区域，不应向任何停滞的敌人作正面攻击。乙、东方军之行动，应仅限于侧击行进中之敌之纵队。为达此目的，须有大的运动机动、包围与迂回。”\footnote{《关于各兵团作战指导及战术的指示（1933年10月9日）》，《项英军事文选》，第214页。}同时，朱德、周恩来也向中革军委建议：“硝石东南为河所阻，恐亦不易强攻，东方军应以一部继续作有力佯攻，催促敌援主力，集结机动消灭援敌。”\footnote{《周恩来年谱1898～1949》，第253页。}双方都准备在硝石吸引对方主力打击之。


根据中革军委指示，彭德怀决定先消灭由硝石南援之敌，尔后乘胜收复黎川。10月8日令红二十师向硝石作有力佯攻；以红三、红五军团及红十九师集结于相埠、乌石地区，待机歼灭援敌。但援敌并未出现，蒋介石判断：“匪西区主力必藏伏在里塔附近，待我军向黎川增援时，其即伺机拊我侧背之狡计甚明，故此时应先严令黎川与硝石各部竭力固守，而我军主力必先设法击溃匪西区之主力，占领里塔官岭前枫林与百花亭一带，构成封锁碉堡线后，再向东进剿其彭匪主力。”\footnote{《手谕先击溃匪西区主力占领里塔一带再向东进剿（1933年10月9日）》，《陈诚先生书信集·与蒋介石先生往来函电》（上），第112～113页。}所以，国民党军暂时在此按兵不动。


围点打援方案不成，东方军遂就势将进攻重心转向硝石。拿下硝石，可截断黎川与国民党军基本区域的联络，使黎川国民党军陷于孤立。此时，国民党军防御硝石的是第二十四师许克祥部，根据红军当时的判断：该部“与红军周旋，多半迭遭损失，士兵对红军及苏维埃有相当影响，战斗力不强武器不精，非蒋介石主力军。我作战均系东方（军）全部，军力武器均比该敌占绝对优势，而又正当洵口战役胜利之后勇气更加百倍”，\footnote{《五军团第十三师硝石战斗经过详报（1933年10月9日至14日）》，《中央苏区第五次反“围剿”》（下），第1页。}对拿下此役相当乐观。但据国民党军将领回忆，硝石国民党守军“自第四次‘围攻’以来，陈就始终控置该师于总预备队中，迄末投入战斗，兵员装备比陈诚所部其它部队亦较为完整”，\footnote{陈振锐：《在陈诚所部参加第五次“围剿”见闻》，《福建文史资料选辑》第2辑，政协福建文史资料编辑室编印，1963，第30页。}并不像红军方面判断的那样薄弱。从战斗结果看，由于对此时国民党军整体战力提升的状况缺乏了解，红军的估计不无轻敌之嫌。


国民党军第二十四师进驻硝石后，立即赶筑碉堡、工事。当红军于10月9日向硝石发起进攻时，该部“依据硝石东北西北一带小高地构筑堡垒一连二十余个死守该地”，计划通过“固守硝石，吸引匪部，再抽调部队应援，聚匪于硝石附近而歼灭之”。\footnote{《五军团第十三师硝石战斗经过详报（1933年10月9日至14日）》，《中央苏区第五次反“围剿”》（下），第1页；《赣粤闽湘鄂北路剿匪军第三路军五次进剿战史》（上），第二章，第25页。}经过数天战斗，战局发展对红军明显不利，面对“工事坚强”\footnote{《赣粤闽湘鄂北路剿匪军第三路军五次进剿战史》（上），第二章，第27页。}的硝石守军，武器、地形处于劣势的红军攻坚战未能奏效，而国民党军的增援部队却大批开向硝石，“硝石西端沿南城河两岸灯火无数、上下不定，系敌之增援队已到达模样”，\footnote{《五军团第十三师硝石战斗经过详报（1933年10月9日至14日）》，《中央苏区第五次反“围剿”》（下），第5页。}红军主力有被围之虞，如彭德怀所说：“当时黎川驻敌三、四个师，南城、南丰各约三个师，硝石在这三点之间，各隔三、四十里，处在敌军堡垒群之中心。我转入敌堡垒群纵深之中，完全失去机动余地。”\footnote{《彭德怀自述》，第185页。}在此形势下，中革军委判断黎川之敌有与硝石据点国民党军内外攻击，相互策应切断东方军后路之可能。为避免陷敌包围，保存有生力量，中革军委决定立即结束硝石战斗，令东方军从现地南移。13日晚红军开始陆续退出战场，红五军团于资溪桥掩护部队转移，红十九、二十师转移至金溪地区，于敌后开展游击战争，红三军团继续在黎南地区待机。


硝石一战后，中革军委原先希望在局部地区采取攻势，以期打破国民党军包围圈的思路已难以继续。10月15日，中革军委发出指示，对战略方针作出调整，明确指出：

\begin{quote}

我们处于战略的防御地位，已成为不可争论的事实，所以我们的战术应根据这一事实来决定。我们东北方面主要的作战任务，曾经是现在还是制止蒋介石的大举进攻。惟有消灭他突击部署中一两个纵队的有生力量，才能达到这一任务。无论取得地方或获得局部的战术胜利，都不能代替这一任务。
\end{quote}


为此，他们规定了几项主要的作战原则：“尽量保存我们自己的有生力量”；“应该避免进攻要塞堡垒地域，甚至避免正面进攻停止的敌人。我们战术的特质就是要搜求运动中的敌人，特别是他的翼侧施行迂回，或因地形和时间的关系施行包围，以及迅速而猛烈地突击敌人纵队第二、第三梯队的翼侧”；“我们在少数的时机也进行临时的防御，但我们不应采取消极的，而应采取积极的和运动的防御。钳制敌人的兵力不得大于三分之一，而将三分之二或更多的兵力，用来机动突击敌人。突击队不是从配置钳制队的纵深，而是从其翼侧出来突击敌人”。\footnote{项英：《关于战术问题的指示（1933年10月15日）》，《项英军事文选》，第221～222页。}这一指示的核心是希望通过运动防御突击敌人一部，通过集中打击敌人有生力量取得战略优势，为此，中革军委要求红军继续在黎川一带寻找战机，利用国民党军新占黎川，各部调动频繁的缝隙觅得在运动中打击对方的机会。与此同时，已经熟知红军战法的国民党军一眼认清了红军意图，判断红军有“故作南北移动，引我注意，吸我调援，以遂其埋伏腰击之企图”，其应付方针一是继续强调“以我为主”，步步为营，二是在黎川与硝石间加紧部署：“对东取积极行动，完成黎硝交通，相机吸引伪东方军而扑灭之”。\footnote{赣粤闽湘鄂北路剿匪北路军第三路军作战计划，《赣粤闽湘鄂北路剿匪军第三路军五次进剿战史》（上），第三章，第12页；第四章，第2页。}


10月21日，中革军委鉴于国民党军周浑元纵队进入资溪桥地区，判断其刚刚进入，应有可乘之机，决定再次实行围点打援，命令：“方面军以在敌的援队未到前，先以有力兵团向周敌堡垒外与其间隙间的部队实施突然袭击”。\footnote{《军委对消灭进攻资溪桥敌之部署的指示（1933年10月21日）》。}22日，红军以第三、十三、十五师向资溪桥发动进攻，同时以红三军团主力集结于洵口、石峡地区，红一军团集中于湖坊地区，战役目标是希望通过攻击牵动对方，“以极大的机动迂回敌之翼侧后方及间隙中积极活动，以迫其与我作野战而全部消灭之”。\footnote{《苏元大元地域（资溪桥东北及西北）战斗详报》，《中央苏区第五次反“围剿”》（下），第7页。}时任第十三师师长陈伯钧日记记载：“敌周纵队在资溪桥、薛纵队在潭头市、硝石之线，我军有消灭该敌之目的。”\footnote{《陈伯钧日记（1933～1937）》，第109页。国民党方面根据缴获的红军作战命令也判断，红军作战目的在“先攻击资溪桥，吸引严和市、潭头市之我军前来增援，而予以夹击”。见《赣粤闽湘鄂北路剿匪军第三路军五次进剿战史》（上），第四章，第3页。}但是，国民党军在此线有10个师左右兵力，红军虽然集中了三、五军团主力部队及一军团部分部队，人数上仍然不占优势。周浑元纵队在资溪桥筑成坚固防御阵地稳守，驻守潭头市的薛岳纵队主力亦不轻动，“十分谨慎，出击的机会很少，即出击亦是小部队短距乘隙袭击，即集结两三个纵队多半取步步为营稳进办法”。\footnote{《彭杨关于我军首先解决信河流域赵部的建议（1933年10月23日）》，《中央苏区第五次反“围剿”》（上），第92页。}红军连攻4天，既不能奈何资溪桥，也没有调动敌人，反而暴露在国民党军堡垒之间，十分被动。


24日、25日，彭德怀、滕代远致电中革军委，提出：“切忌主力摆在敌垒周围，疲劳兵力，日间暴露，在敌机轰炸，晚上大多数露营消耗兵力特甚”；恳切建议：“望以远大眼光过细考虑这些至关重要的问题”。\footnote{《彭滕关于我军今后作战的意见（1933年10月24日）》，《中央苏区第五次反“围剿”》（上），第93页；《彭德怀、滕代远致中革军委电（1933年10月25日）》。}此时，由于赤白对立等问题的影响，红军前出到边区乃至国民党区域作战，面临着群众支持的问题，当时，红军就遇到“军队打仗群众旁观，请不到向导，弄不到担架，脱离群众不发动群众积极性的怪现象”，\footnote{《彭滕关于我军今后作战的意见（1933年10月24日）》，《中央苏区第五次反“围剿”》（上），第94页。}红五军团十三师报告：“每团都有拖枪逃跑的。”\footnote{《苏元大元地域（资溪桥东北及西北）战斗详报》，《中央苏区第五次反“围剿”》（下），第12页。薛岳报告，此役“投诚兵16名，俘匪九十五名”。见《薛岳电蒋中正陈诚》，蒋中正文物档案002090300066397。}这都是作出战略决策时不能不认真掂量的重要因素。中革军委对彭、滕强调的这些困难不以为然，认为：

\begin{quote}

敌人占领资溪桥一二天后，我东方军即到达资溪桥的东南，在这个时间内，敌人能构筑不可进攻的堡垒吗？这是不可能的。有支撑点的坚固堡垒要一个星期才能完成，在二三天之内，敌人只能构野战式的简易堡垒，并只有供给有限的部队（至多半师）。这种堡垒不是不可征服的障碍，如果原则上拒绝进攻这种堡垒，那便是拒绝战斗。\footnote{《中革军委关于十月中战役问题致师以上首长及司令部的一封信（1933年11月20日）》，《中央苏区第五次反“围剿”》（上），第117页。}
\end{quote}


但面对资溪桥连续攻击四日未取得进展、部队损失严重的现实，仍不得不同意退出战线，东方军回师江西后的初步攻击未获成效。红军将领回忆：“从表面上看，我们是在主动进攻，试图在运动战中消灭敌人，实际上由于红军兵力不集中，又是在敌人堡垒密布的白区作战，故我们一开始就处于被动的境地，可说是虎落平川。”\footnote{《莫文骅回忆录》，第229页。}陈诚则在写给妻子的家书中报称：“此间前日匪攻资溪桥（南城光泽间），匪伤亡颇巨（约三千）。我军仍照原定计划，先筑碉堡，实行封锁，再设法歼灭之。”\footnote{《匪攻资溪桥我军照原定计划先筑碉堡实行封锁再设法歼灭之（1933年10月25日）》，《陈诚先生书信集·家书》（上），第236页。}

\subsection{浒湾、团村之役}


硝石之战后，中革军委虽判断红军已经处于战略防御地位，但仍不愿立即退入苏区基本区作战，而是要求继续在苏区外围打击敌人，期望在运动中寻得歼灭国民党军有生力量的机会，规定三、五、七军团组成的东方军“要由目前期待机会的状态逐渐转入进攻敌人突击队翼侧动作”，一、九军团等组成的中央军“粉碎抚河西岸及抚河西的敌人”。\footnote{《军委决定各部队之部署任务（1933年10月29日）》，《中央苏区第五次反“围剿”》（上），第96页。}几次战斗不果，中革军委不免显得急躁，提出：

\begin{quote}

我们在运动战的条件中，如果在原则上也拒绝了进攻敌人的堡垒及野战阵地，这是我们战术的极大错误。我们应经常争取消灭敌人的有生力量于其移动时，然而，敌人新的战术的逐段跃进的移动，一到达指定地点就立刻构筑阵地，并且从其野战阵地出来进攻，一到退却时则缩入于堡垒里面就防御起来。因此，在敌人移动中消灭他的可能渐渐减少了。今后我们可能只消灭敌人的小部队，而其大部队差不多经常是逃脱了的。进行战役而限于这种战术的胜利是不够的，必须突击敌人的主力，追击他并且彻底消灭他才可了事。\footnote{《关于进攻坚固阵地战斗的指示（1933年11月7日）》，《项英军事文选》，第256页。}
\end{quote}


浒湾之战的展开混合着中革军委的上述目标和判断。


浒湾位于国民党军第三路军指挥部驻地南城以北，北路军指挥部驻地抚州（临川）以东，第八十五师驻地金溪以西，距中央苏区北沿约五六十公里，是国民党军在此构筑的切断中央苏区与闽浙赣苏区联系纵深堡垒线的重要支点：“浒湾当公路之要冲，水陆交通均称便利，以故商贾云集，贸易兴盛……金溪较大之地主阶级以及豪绅富户……咸集于是。”\footnote{贺明缨：《江西省田赋清查处实习报告书》，萧铮主编《民国二十年代中国大陆土地问题资料》第170种，第85083页。}从10月下旬起，国民党军第二十六师，第四师、第八十五师陆续进驻抚州、浒湾、金溪地区，其中以5个营的兵力驻守浒湾，并加紧构筑防御工事。由于浒湾位于金溪、临川、南城三点之间，三地都有国民党军重兵驻守，在此作战，对红军而言，存有相当风险。红军的战略仍是围点打援，即通过攻击浒湾，吸引正由金溪向浒湾移动的国民党援军，从而将其包围并消灭之。与此同时，蒋介石则希望利用红军主动出击之机，以优势兵力包围之，他判断红军出击的目的“一则向北在击留我主力，注重北区，不使我南下进剿其老巢；一则向南防我南进而牵制之。惟其两主力仍以黎川为中心，使乘机击我不意，恢复黎川，以打破我五次进剿计划”。\footnote{《蒋中正总统档案·事略稿本》第23册，台北，“国史馆”，2005，第364页。}针对此，蒋介石起初有些举棋不定，10月29日，他计划“以黎川为中心，向右转弯，使其急于突围南窜也”，重点放在逼红军退兵上。11月1日，则反复思量：“此时应分散其兵力乎？抑包围而歼灭之乎？”\footnote{《蒋介石日记》，1933年10月29日、11月1日。}6日，他在抚州确定“战略处置”数日后，他在日记中写下处置的具体内容：“剿匪先完成大包围与断绝之势”，\footnote{《蒋介石日记》，1933年11月11日。}下决心与红军在此展开大战。蒋要求：“我军惟一目的，乃在聚歼匪酋于黎川附近。其次则封锁其南丰与建宁之线，勿使其闯回老巢也。”\footnote{《蒋中正总统档案·事略稿本》第23册，第364页。}


11月11日，战斗开始。12日夜，金溪国民党援军第四师与红军主力在浒湾东北方向八里许的八角亭展开激战，红军集中了红三、五军团5个师兵力将国民党军第四师包围，负责指挥这一战役的彭雪枫颇具信心地认为：“几次都打得不痛快，这一次无论如何是十拿九稳的。”\footnote{彭雪枫：《八角亭战斗的教训》，《彭雪枫军事文选》，解放军出版社，1997，第1页。}由于当地四面都有国民党军重兵，红军围歼战必须迅速解决对方，三军团指挥部明确要求各部应于“拂晓前乘月色消灭被围之敌”。\footnote{《三军团浒湾八角亭战役详报（1933年11月12日）》，《中央苏区第五次反“围剿”》（下），第15页。}对红军的攻势，国民党方面报告：红军“真晨围攻浒湾，势极凶猛，真夜以三团兵力轮流向我冲锋”。\footnote{《谢彬电蒋中正顾祝同据周仲良等称十一日共军第十九师多部围攻浒湾（1933年11月15日）》，蒋中正文物档案002090300066458。}但红军在缺乏强有力攻坚武器的情况下，要实施迅速的歼灭战实在勉为其难。国民党军顶住进攻并在飞机掩护下尝试突围时，红军立即显出力不从心，红军战报写道：天明后，“敌机十二架行低空的袭击四处投弹，各方战斗完全沉寂，敌得于其空军掩护下以主力转向大山（仙）岭八角亭袭击，此时各方均受着敌机威胁而失去了战斗准备”。\footnote{《三军团浒湾八角亭战役详报（1933年11月12日）》，《中央苏区第五次反“围剿”》（下），第16页。}粟裕回忆战场情况时谈道：“十九师是红七军团的主力，战斗力强，擅长打野战，但没有见到过装甲车……部队一见到两个铁家伙打着机枪冲过来，就手足无措，一个师的阵地硬是被两辆装甲车冲垮。”\footnote{《粟裕战争回忆录》，第103～104页。}红军方面的说法可以从国民党方面得到证实，陈诚在家书中告诉妻子：“此役得力于飞机甚大。”\footnote{《昨匪与我第四师在浒湾激战初四师颇危险后转危为安（1933年11月14日）》，《陈诚先生书信集·家书》（上），第239页。}亲身参加浒湾之役的国民党军将领石觉在回忆中也提到，国民党军飞机到达后，猛炸红军阵地，使其“攻势停顿”。\footnote{《石觉先生访问记录》，台北中研院近代史研究所，1986，第49页。}


战斗进行到13日，红军包围圈被撕散，被迫撤出阵地，红十九师“沿着背后的抚河岸边撤了下去，以后才找到了军团部”。\footnote{《粟裕战争回忆录》，第104页。}是役，红军毙伤国民党军520多人，自身伤亡达1100余人。石觉认为红军在这一战斗中表现出的长处主要有：“1.能自主准备战场，布置陷阱，制人而不受制于人。在极优的环境和态势下从事作战。2.攻击精神旺盛，前仆后继，气势逼人。3.数线重叠配备，更番轮流连续攻击，突击力和持续力都很强大。4.守势方面，利用阵前埋伏，出我意料之外。”短处则为：“1.对我军行进路线判断错误。2.人海冲锋的战斗方式，在我狭小而坚强的阵地和炽盛的火力之下，死伤过重。”\footnote{《石觉先生访问记录》，第53页。}


浒湾之战时，11月中旬，红一、九军团奉命在宜黄见贤桥、棠阴地区突破国民党军碉堡封锁北上策应红三、七军团。14日，红一、九军团强行突破封锁线北上。国民党军5个师立即从东、南方向合围红军，“企图截我军于封锁线外”。\footnote{《一军团云盖山大雄关沙岗上等地附近战斗经过详报（1933年11月15日至19日）》，《中央苏区第五次反“围剿”》（下），第21页。}红军主力迅速南返，紧急退到云盖山一带。17日，与国民党军第七纵队3个师在云盖山展开激战。由于国民党军抢得有利地形，红军在激战十余小时后撤出战斗，退向神岗、党口地区。19日凌晨，鉴于国民党军第九师向党口一线挺进，目标较为暴露，朱德、周恩来要求“一军团及十四师应以迅速干脆的手段侧击该敌而消灭之”。林彪、聂荣臻随即根据这一指示下令所部“从敌侧面打击敌人运动的部队而消灭其一部”。\footnote{《一军团云盖山大雄关沙岗上等地附近战斗经过详报（1933年11月15日至19日）》，《中央苏区第五次反“围剿”》（下），第27页。}清晨，战斗打响，红军在大雄关一线与国民党军第九师李延年部激烈交锋。红军经过长途作战，师劳兵疲，“落伍者甚多”，“许多连上，还是前一天早晨吃过饭的”，\footnote{《一军团云盖山大雄关沙岗上等地附近战斗经过详报（1933年11月15日至19日）》，《中央苏区第五次反“围剿”》（下），第26页；罗瑞卿：《大雄关的战斗》，《红星》第19期，1933年12月9日。}但仍然发挥出很强的战斗力，国民党军主阵地967.5高地守军一度被红军压至最后阵地。在攻坚战中红军也付出重大伤亡，红一、二师师长负伤，二师政委胡阿林牺牲，总计伤亡600余人。由于红军始终未能击垮国民党军第九师，而敌第九十师又紧急向第九师靠拢，使“其左侧得到有力依托”，红军已“不宜在此持久战斗及与敌决战”，\footnote{《一军团云盖山大雄关沙岗上等地附近战斗经过详报（1933年11月15日至19日）》，《中央苏区第五次反“围剿”》（下），第28～31页。}19日黄昏，决定退出战斗继续南撤。对这一战，国民党方面战史认为：“我军已占领大雄关左右高地后，匪仍向我攻击，致受重伤，实匪不攻坚战术上之大错误。”同时，红军行动之迅捷和出其不意亦震慑了敌人：“是役上峰命令，匪似已窜至五都、东波（陂？）一带，竟至窜回大雄关与我激战，足见匪实为刁狡；嗣后行动，凡匪在二百里以内，应即准备应战，免受其愚，而遭大害。”\footnote{《赣粤闽湘鄂北路剿匪军第三路军五次进剿战史》（上），第四章，第20页。}蒋对红军能自其碉堡线进出也十分震怒：“窜入碉堡线主匪，仍由见贤桥原路退窜，愧愤之至，将师团长严惩。”\footnote{《蒋介石日记》，1933年11月18日。}


经过两个月的战斗，红军在国民党军的作战圈子里，始终未能实现调动、歼灭国民党军有生力量的目标，国民党军则加紧构筑碉堡、工事和封锁线，红军战略处境日益不利。关于这一阶段战事，后来的历史多称其为“进攻路线”。所谓进攻，当然是相对防御而言的。这一阶段中革军委对单纯的防御战十分担心，强调：“一切持久的防御，都是失败的开始。”“我们战术的基本原则，是要求以敏活的机动来实行进攻的战斗。对于占领的支撑点和阵地实行任何的防御，都是不适宜的。只有在极少的时机，如巩固某地点对于我们具有战略的意义时，则我们才留置小的队伍巩固起来，以求达我们的目的。”\footnote{《关于进攻坚固阵地战斗的指示（1933年11月7日）》，《项英军事文选》，第261页。}他们从第四次反“围剿”主动出击造成的机会中得到启发，希望延续第四次反“围剿”时的战绩。但是，由于国民党军部队密集，层层筑碉，且吸取第四次“围剿”中对红军主动出击战法的忽视，行动更加谨慎，这种战法并没有收到效果。而将初期战斗失利归结为不能攻坚，并企图在一定程度上开展攻坚战，更是错估形势，浒湾之战失败给了这种看法以迎头一击。事实上，这时中革军委已经进退两难，缺乏应对良策。


占领大雄关后，国民党军一方面加紧构筑碉堡、工事，另方面继续向南压迫。11月20日，福建发生反蒋事变，暂时给了红军喘息之机。事变爆发后，蒋介石不得不首先处理福建问题，决定以第二路军向福建出击，江西方面的第三路军则以攻为守，主动向红军发起攻击，抑留红军于其选择的阵地。11月24日发布的北路军作战计划规定该路“由黎川向东南德胜关泰宁方向进展，协同第二路军堵截伪三、七军团，冀歼灭其实力，并竭力掩护第二路军之推进”。\footnote{《赣粤闽湘鄂北路剿匪军第三路军五次进剿战史》（上），第五章，第5页。}12月11日，陈诚令第八纵队“进至团村，相机进占东山，与匪主力决战”。\footnote{《赣粤闽湘鄂北路剿匪军第三路军五次进剿战史》（上），第五章，第8页。}针对国民党军的进攻，红军集中第三、五、七、九4个军团准备阻敌，在团村地区设置埋伏，“预先布置一个师在敌正面钳制，分多组向敌佯动，引敌注意；另以三个师隐蔽在敌之第二梯队左侧后”。“准备取短促突击手段，消灭其前进之一部，迟滞和制止其前进之企图”。\footnote{《彭德怀自述》，第186页；《陈伯钧日记（1933～1937）》，第135页。}12日上午，当国民党军第九十六师、第六师两个师部队进至红军伏击圈后，红军开始发动进攻，双方激烈搏杀，时任红三军团第四师政委的黄克诚回忆：是役，“打得敌人狼奔犬突，消灭了大量敌人，只剩下一个寨子没有攻下来。敌人大部分兵力被我军击溃击散”。\footnote{黄克诚：《回忆在闽赣苏区的部分战斗》，《黎川文史资料》第1辑，政协黎川文史资料委员会编印，1989，第39页。}午后，国民党军退守三都一线，黎川第十一师也赶来增援。陈诚并策划集中7个师部队“全力与匪决战，务求歼灭该匪于团村附近”，双方“抗战异常激烈”。\footnote{《赣粤闽湘鄂北路剿匪军第三路军五次进剿战史》（上），第五章，第12页；《顾祝同电蒋中正近数日诸役报告（1933年12月18日）》，蒋中正文物档案002090300067095。}红军在兵力不占优势，又未能达到消灭敌军一部的不利形势下，13日下午主动撤出战斗，“避免与敌主力决战”。\footnote{《陈伯钧日记（1933～1937）》，第136页。}


团村战斗后，12月15日，国民党军集中3个师兵力以“占领德胜关求匪决战之目的”\footnote{《陈诚电蒋中正顾祝同该路部署明日攻占德胜关求匪决战（1933年12月14日）》，蒋中正文物档案002090300067083。}向德胜关展开进攻。德胜关“当黎川泰宁交通要冲，西南有盐隘草桥隘猴形沙矿各隘口，当江西福建的分界线，其余山岭屏障，林木森严，除此外无路可以通过，但德胜关两侧均系大山毗连，地形狭长，相当减少其防守价值”。\footnote{《五军团得胜关战斗详报（1933年12月15日至16日）》，《中央苏区第五次反“围剿”》（下），第39页。}红军在此有红十五师、三十四师、十三师3个师的防御部队，除十三师外，战斗力都不强，全力防御有些勉为其难。陈诚报告的战斗经过是：“今晨微雨大雾，左右两翼自九时起初各以一部分向德胜关南北侧高地攻击。至戌刻我傅黄两师确实占领德胜关。”\footnote{《陈诚电蒋中正顾祝同攻占德胜关经过（1933年12月14日）》，蒋中正文物档案002090300067087。}退出德胜关后，红军被压往泰宁方向。20日，国民党军又占领黄土关。至此，江西方面战事暂告一段落。


德胜关、黄土关相继攻陷后，国民党第三路军一方面构筑碉堡，巩固自身战线，伺机向苏区进一步深入；另一方面关注福建方面镇压十九路军的战事，准备配合其第二路军的进攻。同时，红军内部在作战指挥上发生争论，朱德、周恩来在团村战斗后致电中革军委，指出如一、三军团会合作战，战果必不如此，强调：“我突击兵团分割作战，在一般干部乃至战时将新战术运用，尚未了解与熟练条件下，常不能达到高级要求的胜利，且常付出过大代价，此点在目前特别重要。”\footnote{《周恩来年谱1898～1949》，第256页。}对中革军委的分兵策略及顶出去打的基本战略含蓄提出批评。16日，针对中革军委24小时内连续四次变更作战命令的做法，周恩来致电博古、项英，强调：“相当范围内职权似应给我们，否则亦请给相机处理之电给我们。”\footnote{周恩来：《请在相当范围内给予部署命令之全权》，《周恩来军事文选》第1册，第313页。}12月20日，中共中央局决定调整红军指挥机构，取消红一方面军总部，并入中革军委机关，中革军委直接指挥中央苏区各军团及独立师、团等。1934年1月起，中革军委代主席项英脱离军事指挥。2月3日，中革军委正式改组，朱德任主席，周恩来、王稼祥任副主席。就红军指挥系统而言，这样的改变应属必要。中革军委作为名义上的最高军事指挥机构，当中共中央有意抬高代主席项英的军事指挥权时，享有相当大的军事指挥权。与此同时，朱德作为一苏大后任命的中革军委主席兼任中国工农红军总司令、红一方面军总司令，周恩来任中国工农红军和红一方面军总政治委员，负有直接指挥作战之责。两个指挥系统难相上下，经常发生颉颃，相互指责不断，项英离开中革军委后，这种局面即被改变。


这一时期，值得强调的是李德在中共中央军事指挥系统中发挥越来越大的作用。李德1932年底来华，埃韦特向国际报告他的到来并称“可能稍后我们要把他派到苏区去”。\footnote{《埃韦特给皮亚特尼茨基的第二号报告（1932年12月初）》，《共产国际、联共（布）与中国革命档案资料丛书》第13册，第264页。}1933年9月26日，作为共产国际驻华军事顾问弗雷德的代表，李德到达瑞金。应该说，共产国际对红军作战指导的态度是谨慎的，共产国际执行委员会只在原则上提供指导，基本不会干预具体的作战方针，强调：“我们关于军事问题的建议不是具有约束力的指示，如何决定由中央和革命军事委员会负责，我们只是提出我们的想法供你们决定。”\footnote{《共产国际执行委员会政治书记处政治委员会给中共中央的电报（1933年9月29日）》，《共产国际、联共（布）与中国革命档案资料丛书》第13册，第509页。}远东局也表示：“你们在当地，应该根据你们的判断行事，并要考虑我们的建议。”\footnote{《共产国际执行委员会远东局给中央苏区的电报（1933年10月30日）》，《共产国际、联共（布）与中国革命档案资料丛书》第13册，第578页。}对中共内部关于军事问题的争执及中共和共产国际顾问之间的争执，埃韦特曾经作出报告：“中共中央和中央军事小组政治领导之间的分歧，以及就我们所建议的每一项重大措施同中央军事小组政治领导发生的持续不断的争论……所有这一切，都妨碍作出一些必要的和迅速的决定。”\footnote{《埃韦特给皮亚特尼茨基的第六号报告（1933年7月28日）》，《共产国际、联共（布）与中国革命档案资料丛书》第13册，第460页。}


当然，共产国际及其顾问的影响仍然非同一般，尤其是李德到达苏区直接参与指挥，在博古配合下，把共产国际的影响发挥到极限。李德在瑞金初期，共产国际军事顾问弗雷德经常会来电进行战争指导，在李德主导下，中共中央数次以现地为由否决了弗雷德的建议，由此导致1933年底至1934年初李德与弗雷德的冲突。而此时远东局书记埃韦特和弗雷德也在福建事变等问题上激烈争执，以致共产国际执行委员会国际联络部驻上海代表格伯特报告：“因为我很了解弗雷德、瓦格纳（指李德——引注）和代表的性格，所以我觉得，和解是不可能的。”\footnote{《格伯特给皮亚特尼茨基的信（1933年12月22日）》，《共产国际、联共（布）与中国革命档案资料丛书》第13册，第646页。}弗雷德和李德之间矛盾激化，甚至到了“越来越主张批评瓦格纳，进而取消他的代表资格”\footnote{《埃韦特给皮亚特尼茨基的信（1934年2月13日）》，《共产国际、联共（布）与中国革命档案资料丛书》第14册，第82页。}的地步，但是结果恰恰相反，1934年3～4月间，和李德及埃韦特两面作战的弗雷德离华，李德遂成为共产国际在华的唯一军事顾问。李德和弗雷德的争执可以看到的结果是以李德的胜利而结束，这大大提高了其发言权，加上不懂军事的中共中央领导人博古对他的依赖，使李德逐渐成为中共军事政策的决策者，根据李德本人后来的检讨：

\begin{quote}

实际上我掌握了红军行动的决定权……我对每个涉及红军的问题都提出了建议，并且直到红军进入贵州省之前我的所有建议均被采纳。除此之外，一些建议只由几个同志进行了讨论，主要是博古同志和周恩来同志，因此造成了对集体领导原则的部分违反。最终我直接干预了指挥部和司令部的工作，我自己起草了作战文件。\footnote{《布劳恩给共产国际执行委员会的书面报告（1939年9月22日）》，《共产国际、联共（布）与中国革命档案资料丛书》第15册，第344～345页。}
\end{quote}


毛泽民在共产国际的报告中直率地指责李德：“党和红军的所有重大事项，只有在取得他的同意之后才能贯彻执行。如果有什么事情没有取得他的同意或者没有按照他的愿望做了，那么他就会开始训人，不管谁都训。”\footnote{《毛泽民给共产国际执行委员会的书面报告（1939年8月26日）》，《共产国际、联共（布）与中国革命档案资料丛书》第15册，第329页。}

\section{福建事变爆发与国共的应对}

\subsection{福建事变中中共的应对}


在中共第五次反“围剿”过程中，1933年11月发生的福建事变是一个存在重要变数的事件。事变在蒋介石封锁线的东方打开了一个巨大的缺口，使江西、福建大块地区成为与蒋介石对立的整体，加上华南广东、广西与蒋介石实际上的离心状态，事变不仅仅对红军的“围剿”，对蒋介石的整个统治都构成了重大危机。但是，福建事变骤起旋灭，蒋介石突然遭遇严重危机，又几乎是兵不血刃轻松获胜。这一结果，和当时复杂的政治、军事背景密切相关。


对福建事变的应对，中共方面可谓一波三折。1927年大革命失败后，斯大林对中国革命形势的基本判断是，随着土地革命的深入，中国革命进入非资本主义发展道路的苏维埃革命阶段。在这一阶段中，资产阶级和小资产阶级出于“对日益增长的土地革命的恐惧”已经脱离革命。根据这一思路，共产国际对中共的指导中强调要划清革命阵营（工人、农民、城市贫民）与非革命阵营的界限。国际远东局代表1929年给中共中央的信中告诫：“你们不应同自称是我们的朋友或者是苏联的朋友的军阀进行任何交谈……如果他们只是把自己说成是自由的拥护者或者是蒋介石的反对者，或者是改组派的拥护者，那就要同他们这些中国劳苦群众的欺骗者进行无情的斗争。”\footnote{《雷利斯基给中共中央政治局的信》，《共产国际、联共（布）与中国革命档案资料丛书》第8册，中央文献出版社，2002，第178页。}在共产国际看来，不仅仅是汪精卫、冯玉祥这些人应该与之“进行认真的斗争”，孙中山的三民主义也应成为批判的对象：“以前我们承认孙逸仙主义的革命意义，这完全是因为民族资产阶级在转入反动阵营之前在革命中所起的作用。现在孙逸仙主义成了整个中国反革命势力的旗帜。”\footnote{《米夫和库丘莫夫给共产国际执行委员会远东局的信》，《共产国际、联共（布）与中国革命档案资料丛书》第8册，第160页。}正由于此，共产国际明确反对中共关于建立民族民主共和国的口号，强调工农民主专政的苏维埃政权的不可替代性。


在共产国际指导下，中共长期执行了关门方针，阻碍了团结更多更广泛同盟军的尝试。九一八事变后，出于对日本侵华后的日苏关系及法西斯在欧洲日益壮大的国际形势的新判断，共产国际在统一战线问题上的方针逐渐有所变化。1933年1月，中共宣言为反对日本帝国主义侵入华北，愿在立即停止进攻苏区、保证民主权利、武装民众三条件下与全国各军队共同抗日。对此，张闻天稍后曾解释道：“这一宣言也是对于所有国民党军阀们说的。在全国的民族危机前面，我们不但要号召工农民众武装起来参加民族革命战争，反对日本与一切帝国主义，而且也号召一切在反动营垒中真正‘爱国的’份子同我们在一起为中国民族的生存而战。”\footnote{洛甫：《关于苏维埃政府的宣言与反机会主义的斗争》，《斗争》第36期，1933年11月26日。}


事实上，虽然共产国际一再强调对国民党内外的各政治、军事势力不能抱有任何幻想，但中共在大革命前后激烈的政治、军事分化组合中形成的与各方间千丝万缕的复杂关系，远非共产国际所能了解，也不是有关原则阐述所能一概抹杀的，共产国际1929年对中共联络俞作柏的批评就证明了这一点。\footnote{1929年底，中共中央和共产国际远东局围绕着中国革命的一些问题发生过激烈争论，其中对俞作柏等的态度问题是这场争论的重要论题之一。远东局指责中共中央“对俞作柏有过幻想”，谈道：“他同你们耍花招，旨在为了改组派的利益来利用共产党日益增强的影响，并使党或某些党组织听其摆布。我们当即看出政治局有些动摇，甚至发展到李立三曾认真考虑能否接受他入党的地步。当广西特委要求公开同俞作柏结盟时，同志们，这没有任何夸张地、十分明确地反映了政治局一时不清楚怎么办和摇摆。”见《共产国际执行委员会远东局和中共中央政治局第三次联席会议记录》，《共产国际、联共（布）与中国革命档案资料丛书》第8册，第296页。}因此，共产国际态度的微妙变化，给了中共一定范围内扩大其活动目标的空间。1933年9月，当第十九路军蔡廷锴、蒋光鼐在沪与中共中央上海局接触始终不得要领，决然与苏区中央联系时，中共虽对其动机有所怀疑，猜测“此种行动极有可能系求得一时缓和，等待援兵之狡计”，\footnote{《苏区中央局致周恩来、朱德、彭德怀、滕代远电（1933年9月23日）》。}但仍对与蒋、蔡接触表现出积极态度。长期在中共掌管组织工作的周恩来指示：“蒋、蔡代表陈公培即吴明，此人为共党脱党者，常在各派中奔走倒蒋运动，并供给我们相当消息……可由国平前往西芹与吴明面谈，更可探知更多内容”。\footnote{《可派人与十九路军代表面谈——1933年9月22日致项、彭、滕电》，《周恩来军事文选》第1卷，第309页。}23日，彭德怀、滕代远、袁国平与蒋、蔡代表陈公培谈判，双方在停战、反蒋态度上基本达成一致，并商定进一步展开接触。停止内战。彭德怀回忆，谈判后，“请他们吃了饭，大脸盆猪肉和鸡子，都是打土豪来的。宿了一晚。我给蒋光鼐、蔡廷锴写了信，告以反蒋抗日大计，请他们派代表到瑞金，同我们中央进行谈判。把上述情况电告中央，中央当即回电，说我们对此事还不够重视，招待也不周，我想还是重视的。招待吧，我们就是用脸盆盛菜、盛饭，用脸盆洗脚、洗脸，一直沿袭到抗美援朝回国后，才改变了这种传统做法。”\footnote{《彭德怀自述》，第182页。}


25日，中央局明确指示“在反日反蒋方面：我们不仅应说不妨碍并予以便利，应声明在进扰福建区域时红军准备实力援助十九路军之作战，在反蒋战斗中，亦已与十九路军作军事之合作过”；强调：“应将谈判看成重要之政治举动，而非简单之玩把戏。”\footnote{《中央局对谈判之指示（1933年9月25日）》，《中央苏区第五次反“围剿”》（上），第132页。}


十九路军与中共主动联系，直接目标是解除身边的军事威胁，其实更重要的砝码还是押向苏联。正由于此，他们首先想到的是与共产国际接触。1933年6月，远东局报告：“19路军司令蔡廷锴建议，通过廖夫人与共产国际代表机构进行谈判。”对此，共产国际反应十分谨慎，强调：“不应当与第19路政府军司令进行任何谈判……您应当从中国同志们那里获得信息。如果他们与什么人进行谈判，那么他们只能以中国共产党的名义进行。”\footnote{《共产国际执行委员会远东局给皮亚特尼茨基的电报（1933年6月10日）》、《共产国际执行委员会政治书记处政治委员会给埃韦特的电报（1933年6月24日）》，《共产国际、联共（布）与中国革命档案资料丛书》第13册，第443、445页。}显然，出于对日本侵华后国际关系变化的认识，共产国际和苏俄对如何与南京政府及地方实力派打交道有自己的考虑，不想成为被利用的因素。这就是为什么他们几乎同时会在北方阻止冯玉祥同盟军行动的原由。\footnote{当1933年8月冯玉祥要组织同盟军时，共产国际致电远东局强调：“考虑到国际形势和冯玉祥想把蒙古人民共和国拖入挑衅行动，以便给日本人提供占领张家口和转向内蒙古的借口……请采取一切措施阻止起义。”见《共产国际执行委员会政治书记处政治委员会给埃韦特的电报（1933年8月9日）》，《共产国际、联共（布）与中国革命档案资料丛书》第13册，第473页。}不过，当十九路军直接与中共接触后，共产国际对此并不反对，鉴于中共面临的巨大威胁，从现实生存和需要考虑，远东局同意中共与福建方面达成协议，9月27日，远东局指示：“同19路军的谈判应尽快以令人满意的方式结束，特别是在与签订停战协定有关的军事问题上。”\footnote{《共产国际执行委员会远东局给中共中央的电报（1933年9月27日）》，《共产国际、联共（布）与中国革命档案资料丛书》第13册，第506页。}10月24日，远东局报告：“蔡告知，原则上他同意我们的建议。”\footnote{《共产国际执行委员会远东局给皮亚特尼茨基和共产国际执行委员会东方书记处的电报（1933年10月24日）》，《共产国际、联共（布）与中国革命档案资料丛书》第13册，第559页。}此中提到的建议即为双方签订停战协定。


10月，十九路军全权代表徐名鸿到瑞金与中共首脑会晤。关于谈判的情况，中共方面代表潘汉年1935年10月在共产国际有一个精彩的报告：

\begin{quote}

10月份他们的代表到来并向我们暗示，他们打算同我们进行认真的谈判时，苏维埃政府责成我作为苏维埃政府的全权代表到汀州同19路军的两位代表谈判。这两位代表表达了同我们进行认真谈判的十分真诚的愿望，甚至表示愿意前往苏维埃中国的首都会见我们中央执委会代表毛泽东同志。但我们中央的某些成员不想让这两位代表进入苏区，因为担心他们是特务。当我给毛泽东同志发去电报后，他不同意这种观点并建议：让他们来。这样，毛泽东同志便把他们请到苏区。

他们到达后，毛泽东同志为他们举行了正式宴会。宴会上，毛泽东同志发表了讲话，阐明了以前公布的统一战线的基本方针。两位代表听到毛泽东同志的讲话后很受感动，以至不知说什么好。让他们发言时他们说，“我们以为，毛是半土匪半游击队的头领，我们决没有想到，他竟是这样一位睿智的政治家”。\footnote{《潘汉年在共产国际执行委员会书记处非常会议上的发言（1935年10月15日）》，《共产国际、联共（布）与中国革命档案资料丛书》第15册，第48～49页。}
\end{quote}


26日，双方代表签订《反日反蒋的初步协定》，规定：双方立即停止军事行动；双方恢复输出输入之商品贸易；福建方面答应赞同福建境内革命的一切组织之活动；双方在上述条件完成后，应于最短期间，另定反日反蒋的具体作战协定。\footnote{《中华苏维埃共和国临时中央政府及工农红军与福建政府及十九路军反日反蒋的初步协定》，《红色中华》第149期，1934年2月14日。}中共与福建十九路军的议和，使中共的战略态势大为改观，不可否认，是中共统一战线政策的重要成果。


随着中共和十九路军关系的迅速缓和，中央苏区东方威胁大大减少，而且开始考虑通过福建方面打破封锁，获得物资援助。10月下旬，共产国际远东局报告，蔡廷锴“将有可能经欧洲然后经福州、延平发出大量的军事技术装备”。\footnote{《共产国际执行委员会远东局给皮亚特尼茨基和共产国际执行委员会东方书记处的电报（1933年10月24日）》，《共产国际、联共（布）与中国革命档案资料丛书》第13册，第559页。}11月，在闽方反蒋弯弓待发时，南京方面的情报称，中共和闽方达成协议：“由闽供给匪军盐卅万元，药品卅万元，兵工器材十万元。”\footnote{《吴醒亚电蒋中正闻李济琛抵闽（1933年11月16日）》，蒋中正文物档案002020200018001。}这应该不完全是猜测之词。对于中共方面希望经由蔡购买武器的设想，国际方面考虑得更周到一些，他们指示：“通过蔡廷锴购买需要的武器装备，我们认为不合适，因为收到定货后形势可能发生变化，他可以把我们的全部定货据为己有。我们建议立即从他那里购买重型火炮、飞机、防毒面具和药品。”\footnote{《皮亚特尼茨基给埃韦特的电报（1933年11月2日）》，《共产国际、联共（布）与中国革命档案资料丛书》第13册，第585页。}此时，蔡廷锴则对获得苏联支持抱有希望，当中共代表潘汉年见到蔡廷锴时，

\begin{quote}

他向我提出的第一个问题是：“您看，如果我们在福建成立政府，并向日本帝国主义宣战，苏联会不会承认我们的新政府？”……当然，我不能对他作出什么肯定的回答，但我想，如果我们真的成立强大的政府，那么苏联一定会支持它的。我补充说，我有一些朋友，可以同苏联人士取得联系。\footnote{《潘汉年在共产国际执行委员会书记处非常会议上的发言（1935年10月15日）》，《共产国际、联共（布）与中国革命档案资料丛书》第15册，第50页。}
\end{quote}


从这段对话可以看出蔡对苏联的期望之切以及与中共接触的良苦用心，而潘的回答也极具外交艺术，更有意思的是，当蔡在潘的提示下谈到其与美国的接触时，潘不仅没有表示反感，而且予以充分的理解，以致蔡高兴地表示：“你们有开阔的眼界。”此中体现的中共党人的灵活态度，绝非习称的所谓“左”倾教条可以概括。


1933年11月20日，以十九路军武力为基础，李济深、陈铭枢及蒋光鼐、蔡廷锴等领导的福建事变爆发。福建方面召集“中国人民临时代表大会”，成立“中华共和国人民革命政府”，改定国号，公然与南京国民政府全面对抗。虽然此前中共与十九路军早有接触，但对事变不悉底细，因此公开态度十分谨慎，代表中共中央态度的是张闻天发表的文章。文章一方面重申愿意在三条件下与任何人订立“作战的战斗协定”，同时申明共产国际的观点，即“整个国民党，不论是南京或是广东，都是投降帝国主义的。广东国民党政府的空喊‘抗日’，不过表明它是另一帝国主义，即英帝国主义的走狗”。文章强调，与靠拢共产党，甚至“抛弃国民党招牌，而采用一些‘动人的’新的名称”的国民党军阀签订反日反蒋作战协定，“这是一种妥协，也许是非常短促的妥协，但是我们并不拒绝这种妥协”。张闻天批评了否认这种妥协的“左”倾幼稚病，同时重点批评了主张信守协定的所谓右倾机会主义，强调：我们的任务“是利用一切条约上的可能去开展在他统治区域内的群众斗争，在最广泛的反日反帝反蒋的下层统一战线的基础上，最无情的揭破一切他们的动摇、不彻底与欺骗，来争取群众在我们的领导之下”。\footnote{洛甫：《关于苏维埃政府的宣言与反机会主义的斗争》，《斗争》第36期，1933年11月26日。}随后，中共中央发布的宣言则以教训的口吻告诫：

\begin{quote}

福建的人民革命政府，如若停止于目前的状态，而不在行动上去证明它真正把言论出版集会信仰示威罢工之自由给与人民，它真正采取了紧急的办法去改善了工农与贫民的生活，它真正准备集中一切武装力量，并且武装广大工农群众去进行反日反蒋的战争，那它的一切行动，将不过是一切过去反革命的国民党领袖们与政客们企图利用新的方法来欺骗民众的把戏。\footnote{《中国共产党中央委员会为福建事变告全国民众》，《红色中华》第133期，1933年12月8日。}
\end{quote}


作为十分重视宣传作用的政党，中共中央的公开表态和真实立场其实存有弹性，这一点，从共产国际的指示中就可看出端倪：“为了不致因同军阀的妥协而损害我们在群众面前的威信，苏维埃政府应该以红军斗争和改善苏区居民经济状况的需要来全面解释自己参加谈判的做法。”\footnote{《共产国际执行委员会政治书记处政治委员会给埃韦特的电报（1933年11月23日）》，《共产国际、联共（布）与中国革命档案资料丛书》第13册，第626页。}因此，与公开表态的尖锐批评不同，中共和共产国际实际对闽态度要复杂得多。陈云在阐述闽变后福建赤色工会任务时谈道：“不估计到一部分群众对于人民革命政府的影响，不估计群众今天觉悟的程度，在群众中简单的把反对国民党南京政府蒋介石，与反对人民革命政府并立起来，或者不在发动群众的运动中在群众面前实际证明人民革命政府不是革命的，而只是空叫反对人民革命政府的欺骗，那就非但不能组织真正群众的革命斗争，而且不能争取群众。”他进一步指出：

\begin{quote}

依照福建目前的具体情形，赤色工会与革命的反帝组织，必须向一切黄色工会与反动派别领导之下的反帝的工人的组织，与向他们的群众，提议建立反日反国民党与反动资本进攻的统一战线。这个行动的目的，是为了团聚一切群众的力量与可能联合，虽然是动摇的力量来反对日本帝国主义与国民党南京政府，是为了争取对于这些组织中的群众的领导。\footnote{陈云：《福建组织“人民革命政府”与赤色工会在福建的任务》，《苏区工人》第5期，1933年12月25日。}
\end{quote}


陈云提到的所谓“反动派别领导之下的反帝的工人的组织”，后人可能会觉得难以索解，但它却是中共在当时理论和现实环境下面对复杂形势的真实反映，字面的批判和实际的争取在这里被奇特地统一起来。1933年11月23日，共产国际针对广州方面寻求与苏区接触的电文明确指示：“如果广州人意在反对19路军，那我们就不该同他们谈判，以期不削弱19路军反日反蒋的立场。”此中所表现的对福建方面的支持和信任相当明显。而共产国际代表更直接就双方共同组建军队征询国际方面：“是否同意苏维埃政府建议蔡廷锴靠自己的力量组织不叫红军而叫人民革命军的工农武装？我们有人，但没有武器。这些力量应该同红军和福建军队实行合作来对付南京的进攻。”\footnote{《共产国际执行委员会政治书记处政治委员会给埃韦特的电报（1933年11月23日）》、《埃韦特给皮亚特尼茨基的电报（1933年11月24日）》，《共产国际、联共（布）与中国革命档案资料丛书》第13册，第626、628页。}


中共驻共产国际代表团对当时中共中央决策有着重要影响。1933年11月底、12月初，代表团两位主要领导王明、康生在共产国际第13次全会大会发言中都谈到了福建事变。王明认为：“个别的军阀派别，在和红军屡战屡败之后，在兵士群众和一部分下中级军官革命情绪压迫之下，不能不这样提出问题：或者继续和红军战，那么，毫无疑问地他们将要完全塌台；或者把自己枪头掉转过去反对日本及其走狗蒋介石，以便从绝路上另找出路。”\footnote{王明：《革命、战争和武装干涉与中国共产党底任务》，《王明言论选辑》，第332～333页。}显然，这一判断注意到了事变的积极意义，并不完全对其采取否定态度。康生则指出：

\begin{quote}

在对红军作战失败后，十九路军的部分指挥员认识到，若继续对红军发动进攻，那势必要招致更大的失败；但倘若不打红军，而打蒋介石和日本帝国主义，给老百姓以起码的民主、自由，那么，这支军队是能够使中国获得解放的。无论美国及其它帝国主义者是否想利用福建事变以达到它们反对日本的目的，无论福建的将领们是否会始终如一地实现自己的诺言，抗日运动结果，给福建劳动群众和十九路军士兵指出了一条道路，以便实现中华苏维埃政府和红军在1933年1月10日关于抗日武装斗争统一战线的宣言中所提出的要求。\footnote{《康生在共产国际执委会第13次全会第17次会议上的发言》，《共产国际有关中国革命的文献资料》第2册，中国社会科学出版社，1982，第251页。}
\end{quote}


康生这段讲话对事变的肯定更加正面，而且含蓄否定了有关十九路军是代表美国利益反对日本的说法，明确十九路军的行动有实现中共倡导的抗日武装斗争统一战线和使中国获得解放的可能，这样的表态不可忽视。由于中共代表团在中共中央决策体制中的特殊地位，他们的态度不可能不对中共中央发生影响。


事实上，中共地下党组织在福建的一系列行动也体现着中共对闽变欲迎还拒的真实态度。11月29日，中共厦门中心市委将原定在人民政府召开群众大会上的示威行动改成参加大会并宣传中共方面的主张。\footnote{参见《中共厦门中心市委给中央的报告（1933年12月4日）》，《福建革命历史文件汇集》甲11册，第211～213页。}福州中心市委在未接到中共中央指示信前曾正确判断“此次政变是有利于我们的，认为蒋介石才是我们的最大敌人，是冲破五次‘围剿’有利的条件。即决定发动群众参加反蒋大会”。但随后改变了这一看法，采取了与人民政府对立的一些做法，“犯了‘左’的关门主义的错误”。\footnote{《中共福州中心市委符镭关于福州工作情况给中央的报告（1933年12月4日）》、《中共福州中心市委关于工作情况给中央的报告（1933年12月16日）》，《福建革命历史文件汇集》甲13册，第177、202页。}12月中旬，在得到有关指示后，福州中心市委调整政策，决定将游击队改为人民抗日军，进一步开展反帝运动。罗明回忆，事变后，“中央局派我和谢小梅前往厦门、福州领导党组织推动地方群众工作。厦门、福州两个中心市委召开扩大会议，由我传达中央局指示，并决定开展群众抗日反蒋的民主运动，要求释放政治犯，加强党组织的发展和秘密工作”。\footnote{《罗明回忆录》，福建人民出版社，1991，第162页。}


作为政策转变的标志，中共这时把工作重心转到反蒋和中共自身的发展上，不再以十九路军为对抗对象。十九路军在福建连江开展“计口授田”工作时，正在此准备进行土地革命的中共方面当即表示：“如‘人民政府’实行分田”，他们“当让‘人民政府’去做。”\footnote{何公敢：《“福建人民政府”和“生产人民党”片段》，《文史资料选辑》第37辑，第96页。这一回忆可从事变基本失败后福建临时省委致连江县委信中得到侧面证实，信中批评连江县委没有充分进行士兵工作，“而只是与上层的官长来往”。见《中共福建临时省委致连江县委信（1934年1月30日）》，《福建革命历史文件汇集》甲7册，第214页。}同安地方政府召集大会时，为凑集人员求助于中共当地负责人，中共负责人“叫他们到各乡村公开打锣去号召，用汽车去载人，弄点心给农民吃”。\footnote{《中共厦门中心市委给中央的报告（1933年12月4日）》，《福建革命历史文件汇集》甲11册，第214页。}福州党组织慰劳队到前线慰劳十九路军，以“号召广大群众参加反帝反国民党的统一战线”。\footnote{《中共福州中心市委给中央的工作报告（1934年1月5日）》，《福建革命历史文件汇集》甲13册，第210页。}陈子枢甚至为厦门中心市委得知中共和福建人民政府签订协定后“表现出一种空洞的欣喜的情绪”感到担心，忧虑“这样必然会放弃了实际的艰苦的斗争工作，不可避免地由于这种错误的幻想，会有走上机会主义道路的危险”。\footnote{《陈子枢给中央的信（1933年12月16日）》，《福建革命历史文件汇集》甲13册，第205页。}中共在福建和十九路军的微妙关系，外界其实也多有注意：“共党宣传，表面上仍属反闽，实际上则共党密赴福建者甚多，对民众组织甚努力。”\footnote{《闽方内部意见纷歧》，1933年12月8日《申报》。}


由于中共与闽方政治上欲迎还拒而又相互有所企图的合作关系，当十九路军的反蒋军事展开后，中共与闽方的军事合作也呈现复杂的内容。蒋介石向闽方发动进攻之始，中共从自身战略利益出发，采取了有保留的协助措施。1933年11月24日，朱德、周恩来等要求赣东北和闽北中共武装，在国民党军向闽浙赣边境集中时广泛发展游击战争，扰乱其后方，红七军团主力则随时准备截击或尾追敌人。同时，周恩来致电中革军委，探询可否出动红三、五军团侧击国民党军入闽部队。\footnote{《周恩来年谱1898～1949》，第254页；《周恩来致中革军委电（1933年11月24日）》，《周恩来军事文选》第1卷，第315～316页。}中革军委对直接与十九路军配合作战抱有疑虑，强调“不应费去大的损失来与东北敌人新的一路军作战，而让十九路军替我们去打该敌”，但同时也指示红七军团及独立部队应“以游击战争的方式妨害敌人第一路军集中”，在“敌人第一路向闽北前进时阻滞并剥削之”。\footnote{《军委关于方面军动作的训令（1933年11月25日）》，《中央苏区第五次反“围剿”》（上），第123页。从国民党军提供材料看，其第五路军入闽时，红军在资溪至光泽一线确有节节抵抗的行动，但面对国民党军4个师的强厚兵力，红军的有限兵力事实上难以抵挡。见《国民党北路军顾祝同部与中央苏区红军作战情形报告书》，《中华民国史档案资料汇编》第5辑第1编军事（4），第20～21页。}12月中旬，中革军委要求红五、七军团及独立第六十一团组成的东方军，在建宁、泰宁、邵武、光泽、黎川一带展开游击战争，“侧击向资溪、光泽运动中敌人的中央纵队”，“在资溪务须进行顽强的防御”，“迟滞、钳制向光泽前进的敌人”。\footnote{《军委关于目前作战的决定（1933年12月13日）》，《中央苏区第五次反“围剿”》（上），第126页；项英：《关于钳制向光泽前进敌人之部署（1933年12月7日）》，《项英军事文选》，第310页。}12月中、下旬，朱、周又多次指示“广泛发展游击战争”、侧击敌“进剿部队之后尾”，“积极扩大并发展闽中游击战争，不断截击邵顺间敌人后方联络部队及进行一切破坏工作”。\footnote{《朱德、周恩来致刘畴西、聂洪钧并报项英电（1933年12月11日）》，《闽浙皖赣革命根据地》（上），第672页；《朱德、周恩来致寻淮洲、乐少华电（1933年12月28日）》，《周恩来军事文选》第1卷，第316页。据当时报载：“泰顺边境十九晚侵入匪千余，枪械不齐，经该处保安队中央军痛击激战，将二十一晨十时，至匪击退，现匪在寿宁坑底以北，仍企图进攻浙边。”见《大股赤匪犯泰顺边境被击退》，1933年12月22日《申报》。}12月底，项英明确指示红九军团第十九师“转移到将乐的地域，与十九路军的左翼部队取得直接连络”。\footnote{项英：《关于各军团部署及任务的指示（1933年12月29日）》，《项英军事文选》，第341页。}当时前方国民党军对红军行动的亲身感受是：“企图破坏道路阻我前进，无顽强抵抗力。”\footnote{《李默庵电蒋中正顾祝同卫立煌等（1933年12月17日）》，蒋中正文物档案002090300067096。}


中共的这一系列行动得到共产国际的支持，1934年1月2日，共产国际提出：“你们的部队应该避开蒋介石的进攻部队，在迅速重新部署后，从北面和南面同时打击进攻19路军的蒋介石部队的侧翼和后方，最好是在它们之间发生冲突的时候。”\footnote{《共产国际执行委员会政治书记处政治委员会给中共中央的电报（1934年1月2日）》，《共产国际、联共（布）与中国革命档案资料丛书》第14册，第8页。}就中共而言，维持十九路军这一反蒋力量的存在，是一种符合其利益的常识判断，罗明回忆，事变后他被派往厦门、福州“传达中央对闽变的指示，即军事上联合，帮助他”。\footnote{《罗明同志的回忆》，《党史参考资料》第4期，中共福州市委党史资料办公室编印，1982。}而蒋介石也报告：“最近赤匪因得闽方接济，并为牵制策应闽方及实行联成一片起见，突以伪第一军团向我崇仁、宜黄一带进犯，并以伪军第三、五、七各军团，集合全力，分向我黎川、金溪一带进犯，对我冀作中央突破及右翼迂回之企图。”\footnote{《蒋中正电汪兆铭叶楚伧等（1933年12月15日）》，蒋中正文物档案002090300068053。}国民党浙江省政府则注意到：“赣东赤匪乘中央军进迫闽垣之际，迭犯浙赣边境，企图牵制我方兵力，截断我后方联络。”\footnote{《浙江省政府呈行政院工作报告（1934年1月）》，《中央革命根据地革命与反革命斗争史料》（下），中国社会科学院近代史研究所藏。}


虽然红军受命骚扰入闽的国民党军，但要以实力代价为十九路军在闽北阻挡国民党军，也为中共中央所不取。在国民党军兵力居绝对优势情况下，红军要独力担起阻止国民党军入闽重任其实并不现实，何况，本在闽北的十九路军主力也在后撤。因此，虽然王明在事变后不久曾批评中共中央没有意识到：

\begin{quote}

问题的中心不在于我们红军愿意不愿意接受蒋介石这个打击，而问题的实质是在于：或者红军和十九路军一起来击退蒋介石的力量，或者是蒋介石先打败我们的同盟军——十九路军，然后再集中一切力量来打击我们的红军。\footnote{王明：《新条件与新策略》，莫斯科1935年印行，第64～65页。}
\end{quote}


但这究竟是事后诸葛亮之言。事变当时，中共和十九路军对双方的军事合作都没有充分的准备，对军事形势的发展也难有全盘的计划，在国民党军十几个师压迫下，中共即使勉为其难，效果也难以臆测，何况，国民党方面还观察到：“伪一、三、五等军团被我痛剿，迭受叵剧，喘息未定。”\footnote{《国民党军卫立煌部镇压“闽变”战斗详报》，《中华民国史档案资料汇编》第5辑第1编军事（5），第805页。}这虽然不无自夸之嫌，但红军自第五次反“围剿”以来连遭失利、军力受创、亟需休整确也属实。\footnote{朱德在1937年口述的自传中也谈道：当时红军出动兵力过小，“只以一个七军团去打，力量少小，当然没有牵掣得着……我们当那时，却想休息疲乏，就没有进行”。见《朱德自传（1886～1937）》，转见金冲及主编《朱德传》，人民出版社，1993，第317页。}


更重要的，事变发动后，中共和十九路军间在联合作战问题上，始终没有达成协议。潘汉年回忆：“在福州，蔡廷锴向自己的指挥部示意，有苏维埃政府的全权代表在这里，但同时他却不让师长们直接同我交谈。”\footnote{《潘汉年在共产国际执行委员会书记处非常会议上的发言（1935年10月15日）》，《共产国际、联共（布）与中国革命档案资料丛书》第15册，第52页。}蔡廷锴既希望借助中共及其背后的苏联力量，又对中共不无防范担心的微妙心理，于此显露无遗。对于福建方面的联共之举，闽方许多中高级将领也缺乏心理准备。十九路军与红军曾多次交手，就在事变发动前几个月，红军东路军还在与之作战，官兵中有一定的敌视心理。在事变前蔡廷锴主持的动员会上，许多将领“对同红军和平相处，对反对南京政府之事一无表示”，部分将领则“在政治上不同意即行反蒋，也不同意与红军合作”。\footnote{蔡廷锴：《回忆十九路军在闽反蒋失败经过》，《文史资料选辑》第59辑，第92～93页。}一些团长明确提出：“十九路军历来是反共的，为什么要和共产党合作？”\footnote{麦朝枢：《福建人民革命政府回忆》，《文史资料选辑》第37辑，第88页。}这反映了相当部分中高级干部的看法。曾任福建省政府委员的林知渊评判：“他们即使和中共合作成功，也只是貌合神离，只求各保边境，互不侵犯，并没有想到要进一步联合行动，统一作战。”\footnote{林知渊：《政坛浮生录——林知渊自述》，《福建文史资料》第22辑，政协福建文史资料委员会编印，1989，第54页。}这一说法颇中肯綮。至于发动事变的陈铭枢则“过于不重实在军事，毫无打仗准备”。\footnote{《冯玉祥日记》第4册，江苏古籍出版社，1992，第269页。}在福建人民政府参与机要的麦朝枢回忆，12月底，中共电告闽方，蒋军两个师东向闽境推进，陈铭枢得知情况后表示：“江西境内有红军，当可以把蒋军击退，不必顾虑。”而在麦看来，“十九路军和红军合作的具体条件还没有订定，红军实在没有代我们挡击蒋军的义务，我们为什么不派兵警备呢？”\footnote{麦朝枢：《福建人民革命政府的回忆》，《文史资料选辑》第37辑，第89页。}麦朝枢的这一疑问其实相当可以说明问题，当十九路军自身已在选择从闽西北后撤时，和其并无约定的中共确无代十九路军牺牲之理。


虽然中共对全力支持十九路军心存疑虑，但是，完全放任国民党军进攻十九路军也不符合中共利益。中共虽未选择以主力在闽赣边境与国民党军作战，但中革军委却谋划了一个更大更全面的战略构想：将中央苏区红军编为东方军、中央军和西方军，准备分路作战。红五、七军团编为东方军，在福建建宁、泰宁、邵武、光泽一带展开游击战争，钳制东线国民党军，“侧击向资溪、光泽运动中敌人的中央纵队”；红九军团编为中央军，在东方军左翼活动，防止江西方面国民党军进入苏区纵深；主力部队红一、三军团编为西方军，挺进到永丰地区活动，“从左翼绕至蒋介石军之后方，就是说渡过赣江由西向九江南昌进攻，以协同第十九路军前后夹攻”。\footnote{《军委关于目前作战的决定（1933年12月13日）》，《中央苏区第五次反“围剿”》（上），第126页；李光：《中国新军队》，第219页。}在中革军委看来，该计划既避开了以红军主力在闽西北直接与国民党入闽军决战，从而为十九路军火中取栗，成为其掩护部队的结果，在北线“敌人最弱的地位”主动出击，又不无围魏救赵之意，客观上帮助了十九路军；同时还可抓住国民党军东移机会，使中央苏区“打通与基本区域的联系”即湘赣苏区的联系，进一步在北线打开缺口，北出昌、九，全面打乱国民党军部署，争取战略主动，壮大自身。\footnote{《军委关于转移突击方向和组织三个军及各军动作的指示（1933年12月20日）》、《军委关于目前作战的决定（1933年12月13日）》，《中央苏区第五次反“围剿”》（上），第129、125～126页。}不失为一石三鸟之计。


为实施这一计划，湘赣的红十七师受命由湘赣苏区北上，出击南浔路，“和中央红军配合十九路军行动”；\footnote{萧克：《红十七师北上行动的回顾》，《湘赣革命根据地》（下），第1042页。}湘鄂赣的红十六师也根据中央电令，“向高安、万载附近行动”，\footnote{傅秋涛：《关于湘鄂赣边区内战中期后期历史情形报告》，《湘鄂赣边区史料》，第57页。}威胁南昌。李德回忆，这一决策是由共产国际执行委员会的军事代表（弗雷德）建议，“朱德、周恩来、博古和我讨论了这个建议，并对此基本上表示同意”。\footnote{〔德〕奥托·布劳恩：《中国纪事1932～1939》，第63页。}红十六师和十七师的出动确曾给国民党军造成一定压力，蒋介石致汪精卫电中谈道：“赤匪一部最近窜至南浔路附近地区，确是事实。”\footnote{《蒋介石致汪精卫等电（1934年2月25日）》，《中央革命根据地革命与反革命斗争史料》（下），中国社会科学院近代史研究所藏。}蒋更曾致电抚州方面的陈诚，告以：“清江万寿宫被匪占领，省防可虑，请抽调第五师或九十六师中之一师星夜集中南昌，用车运输为要。”\footnote{《手谕省防可虑请抽调部队星夜集中南昌（1933年12月27日）》，《陈诚先生书信集·与蒋中正先生往来函电》（上），第121页。}红军对国民党方面形成的压力可见一斑。杨永泰直承，当红十六师逼近南昌时，“南昌夙未留存预备部队……泰等在此唱了两夜空城计，仅能用飞机轰炸以威胁之”。\footnote{《杨永泰致汪精卫等电（1933年12月29日）》，《中央革命根据地革命与反革命斗争史料》（下），中国社会科学院近代史研究所藏。}国民党“围剿”军西路军总司令何键也谈道：

\begin{quote}

最近江西方面的匪情，不是北路之匪想渡江西窜，就是西路之匪想向东窜，所以西路军现在所负任务，一方面要消灭萧孔各匪，一方面要不许萧匪等东窜，同时又一方面要堵截北路之匪不能渡江西窜，因此每三方面同时发生战事。\footnote{《总司令在扩大纪念周中报告中央军在闽赣作战剿匪之胜利情况》，《西路军公报》第8期，1934年2月15日。}
\end{quote}


应该承认，中共中央这一计划和当时彭德怀建议并在日后得到毛泽东肯定的出江、浙进扰国民党军后方的设想可谓异曲同工。


不过，红军真要在南昌一线与国民党军大规模作战，也难有成算。这一带平原坦荡，堡垒众多，是国民党统治力量较强地区；虽然闽变吸引了国民党军一部分兵力，但其在北线的绝对优势并未改变。红一军团北出不久，即在丁毛山战斗遭到较大伤亡，北进计划难以实施。事实上，由于福建迅速出现不利形势，计划中的将红三军团西移设想并未实施，当蒋介石基本完成其对福建的部署但大规模战斗尚未打响时，红军已作出直接援助十九路军作战的决策。1月2日，中革军委命令红三军团向福建“沙县地域移动”。\footnote{项英：《关于第三第七第九军团的动作问题（1934年1月2日）》，《项英军事文选》，第345页。}滕代远记载：“中央军委当即令红军第七军团由太宁往将乐、顺昌地区以协助十九路军向浙边进展。同时命令已开抵广昌地区的红军第一军团、第三军团，星夜取捷径向顺昌开拔。并以红军第九军团位于建宁归化地区，以第五军团在黎川地区支持北面蒋军的进攻。”\footnote{李光：《中国新军队》，第220页。}但是，由于十九路军防线轻易瓦解，蒋介石又未给其任何喘息之机，红军来不及实现对十九路军的增援。所以中共后来说：“我们红军为了配合他们反蒋的战争，曾经在闽北积极行动，从占领沙县直下尤溪，但是这对那些表现丧魂失魄的狐群狗党依然是无用的帮助。”\footnote{《中华苏维埃共和国中央政府为福建事变宣言》，《中央苏区第五次反“围剿”》（上），第151页。}这固不无推脱责任之意，却也不全为空穴来风。


另外，还应指出的是，1934年1月下旬举行的第二次全国苏维埃代表大会上，毛泽东曾谈道：

\begin{quote}

有一个同志对于福建的所谓人民政府，说他带有多少革命性不是完全的反革命，这种意见也是不对的。我在报告中已经指出：“人民革命政府”的出现，是反动统治阶级的一部分，为着挽救自己将死命运而起的一个欺骗民众的新花样，他们感觉苏维埃是他们的死敌，而国民党这块招牌太烂了，所以弄个什么“人民革命政府”，以第三条道路为号召，这样来欺骗民众，没有真正革命意义，现在事实已经证明了。\footnote{毛泽东：《关于中央执行委员会与人民委员会报告的结论》，《苏维埃中国》，第305页。}
\end{quote}


毛泽东对福建事变的这种评价，征诸前引潘汉年回忆，当然并非如李德所称毛泽东当时坚持“不应该马上直接支持十九路军和‘人民革命政府’”，\footnote{〔德〕奥托·布劳恩：《中国纪事1932～1939》，第85页。}而是反映着事变后中共中央的态度。福建事变失败后，整个中共领导层一改事变中的谨慎态度而展开谴责、批判，在此背景下，毛的表态自也不能出此框架，它倒是提醒我们，正像不能简单用毛泽东事变后的发言等同于事变中毛泽东的态度一样，事变后中共中央对闽方的定性也不完全代表其事变中的真实态度。虽然，中共中央可以被批评缺乏有效措施将事变引向自己有利的方向，但如果看看蒋介石当年的部署，就可能发现，事情远不是那么简单。

\subsection{南京政府镇压福建事变}


福建事变爆发，对中共是机会，对蒋介石，却也未必全出意料。


福建的十九路军是淞沪抗战后进驻福建的。1932年蒋介石重掌南京中央后，作为粤陈（铭枢）势力的十九路军戍守南京势难继续，粤桂方面提出将十九路军调驻福建。对此，蒋介石内心并不情愿，在给何应钦的电报中谈道：

\begin{quote}

近日李黄陈诸兄急欲派十九路赴闽，其势似不可阻止。汪院长亦已赞成，其事必实现。如此恐伯南调赣南部队回粤，又碍中央剿匪计划。故中迟迟未肯下令也。前日罗师长回赣时托其面述此情，并派其往见余幄奇，最好留余部在赣南完成剿匪使命。但其直辖于伯南，如我方往留，则于公于私皆有为难。\footnote{《蒋中正电何应钦近日陈济棠等欲派十九路赴闽听其自决枪械暂勿送去（1932年5月13日）》，蒋中正文物档案002010200066033。}
\end{quote}


虽然心有戚戚，但蒋介石当时没有其他安置十九路军的办法。十九路军的抗日英名，蒋介石自身刚刚复职的脆弱，使其最终不得不同意十九路军赴闽。


十九路军到闽后，迅速控制福建局势，并与粤方谨慎接触，双方关系若即若离。蒋介石对十九路军以拉拢为主，但也不无搞垮十九路军，将福建收为己有的心思。1933年2月，陈立夫向蒋介石报告福建和粤桂形势，透露出南京中央图闽的隐秘动机：

\begin{quote}

蔡对闽省客军极仇视，而于卢兴邦部为尤甚，常欲伺机解决之。闻中央接济卢部机枪五十挺迫击炮八尊之讯为蔡所悉极度不安。2.刘珍年之调驻浙东与闽北，配之以卢兴邦与刘和鼎诸部在蔡视之为中央对十九路军之包围。3.中央此次调六十师赴赣剿匪在蔡视之为有分散其兵力。

西南中心系于陈济棠之一身，陈如效忠中央，则西南风云可以消，盖无广东则西南活动将无经济基础。陈氏乃解决西南问题之锁。\footnote{《陈立夫电蒋中正缕陈西南问题（1933年2月7日）》，蒋中正文物档案002090300009103。}
\end{quote}


1933年5月，十九路军的老上司陈铭枢游欧回国，开始积极筹划反蒋。参与陈铭枢策划的梅龚彬回忆，陈提出上中下三种方案：“第一种方案（上策）是联合粤桂反蒋；如果陈济棠不肯参加的话，就执行第二种方案（中策），先搞闽桂联合倒陈，再发动反蒋；如果陈济棠和李宗仁都不肯干，那只有采取第三种方案（下策），争取与红军合作反蒋。”\footnote{《梅龚彬回忆录》，团结出版社，1994，第82页。}陈铭枢的活动，南京方面迅速得到讯息，与粤方有千丝万缕联系的汪精卫早在6月15日就电蒋报告广州方面的异动，提醒蒋“恐将有军事行动”；17日，更直接点出浙江有被犯之虞：“浙省空虚，不肖生心，乘虚冒进固愚妄所为，但天下乱事往往由愚妄之人所造成。不如益兵为备，使之知难而退，弟固确有所闻故力言之，并非欲轻启兵衅也。”\footnote{《汪兆铭电蒋中正西南分离运动已成熟旬日内必实现其自立政府（1933年6月15日）》、《汪兆铭电蒋中正浙江省应益兵为备（1933年6月17日）》，蒋中正文物档案002080200097032、002080200097100。}汪精卫在此故作玄虚，并未点明具体的犯浙者，但衡诸时、地二势，有可能对浙江构成威胁的，必为福建无疑。


其实，陈铭枢的活动，蒋介石多有掌握。在陈铭枢接触陈济棠、胡汉民未取得进展时，蒋介石在日记中记有：“陈铭枢等联合反动，似告失败，则西南渐稳。”\footnote{《蒋介石日记》，1933年6月17日。}17日，他致电吴铁城时表示：“西南一切酝酿，一切误解，应恳切劝导，设法消弭，必尽其在我。如仍逞私见，害大局，吾人职责所在，固不容瞻顾畏缩也。”\footnote{《蒋中正总统档案·事略稿本》第20册，第592页。}虽然对广州方面和陈的活动有所警戒，但蒋的判断还是偏于乐观，认为其一时难成气候。8月，吴铁城也致电蒋介石报告：“粤闽军事联络恐难实现。”\footnote{《吴铁城电蒋中正顷接港讯现陈铭枢拟扩充十九路粤闽军事联络恐难实现（1933年8月23日）》，蒋中正文物档案002080200115183。}


作为在长期内部混战中摸爬滚打出来的实力派人物，蒋介石当然不会对陈铭枢的活动完全掉以轻心。虽然按照蒋介石惯常的坐观其变处事方式，他未对福建方面和陈铭枢采取积极行动，但并不意味着对此无所作为。准备第五次“围剿”时，蒋在浙赣闽边区部署警备部队5个师另4个保安团，这样的重兵配置，极具防备闽方的意味。尤其是1933年9月，蒋介石令国民党军进攻位于闽赣边境的黎川，应为一石二鸟之举，既防范赣南和赣东北红军的联系，对其后来的进兵福建也大有裨益。


1933年10月，陈铭枢活动益繁，陈济棠曾电蒋介石，请其适当安置陈铭枢、李济深，以免引起异动，但未得到蒋的积极回应。\footnote{《陈济棠电蒋中正李济琛陈铭枢恳给以名义（1933年10月3日）》，蒋中正文物档案002020200018001。}稍后，蒋介石又接到戴笠的报告：“陈铭枢前来闽用意在与蒋蔡密商联络桂系倒蒋，以求西南切实联合，反抗中央。”\footnote{《江汉清电蒋中正转报陈铭枢赴闽联络蒋光鼐蔡廷锴意图反抗中央（1933年10月18日）》，蒋中正文物档案002090300009011。江汉清为戴笠化名。}对此，蒋介石仍然没有明确反应。11月9日、10日，朱培德连电蒋介石，告以福建陈铭枢等“谋不轨”的消息，建议其速劝时在福州的国民政府主席林森“回京坐镇”。\footnote{《朱培德电蒋中正接黄实来电陈铭枢等谋不轨（1933年11月9日）》、《朱培德电蒋中正转报陈铭枢等谋不轨闽将先发》，蒋中正文物档案002020200018002、002090300009012。}此时，坐以观变的蒋介石方才出手。11日，蒋介石致电林森，望林“即日回京”并代劝陈铭枢“回中央襄助一切”。\footnote{《蒋中正电林森可否劝陈铭枢回中央襄助一切》，蒋中正文物档案002020200018005。}12日，蒋介石在日记中自我安慰：“陈铭枢入闽作乱，消息渐紧，但无妨耳。”15日，得到福建将有事变的确实消息，蒋当夜“几不成寐”；次日仍“对闽事，思虑入神，不觉疲乏”。\footnote{《蒋介石日记》，1933年11月12、16日。}16日，他做最后的努力，拿出惯常的封官许愿招数，致电蒋光鼐：“许陈军事总监或参谋总长，内政部长亦可。”\footnote{《蒋中正电嘱蒋光鼐探询陈铭枢任职意愿及其反抗中央之传闻（1933年11月16日）》，蒋中正文物档案002010200097057。}但这样的表态，未免失之太晚。


事变既起，在判断其将局限于福建范围内后，蒋介石迅速确定军事解决闽变的方针。十九路军原辖3个师，1933年6月扩充2师，总共有5个师10个旅，每师4000～4500人，加上直属部队，实际兵力5万人以上。\footnote{参见《闽方逆军新编部队番号及各级逆首姓名调查表》，《中华民国史档案资料汇编》第5辑第1编军事（5），第806～807页。}事变之初，戴笠即向蒋介石报告，闽方“新兵多，逃亡众，能作战者不上三万五千人”。\footnote{《江汉清电蒋中正十九路军逃兵众多请速发兵戡乱》，蒋中正文物档案002090300009326。}据此，蒋介石致电汪精卫表示：“总计逆军号称六军十二万人，实际能作战者最多三、四万人。”\footnote{《蒋介石致汪精卫等电（1933年12月13日）》，《福建人民政府与共产党合作反蒋史料》，中国社会科学院近代史研究所藏。}对于蒋介石而言，这样的实力并不足以构成致命危险。何况十九路军“此次师出无名，其军心必动摇，干部钱多，必不如前之肯牺牲”。\footnote{《电呈预防西南异动及应付闽变之刍见（1933年12月24日）》，《陈诚先生书信集·与蒋中正先生往来函电》（上），第120页。}1933年12月5日，在给驻日公使蒋作宾的电报中，蒋介石乐观判断：“闽变必可速平，饶有把握。”\footnote{《蒋中正总统档案·事略稿本》第24册，第21页。}而陈诚早在12月中旬对事变的趋势也作出了准确预测：

\begin{quote}

闽变当不难解决，报载军事行动多不确。现我军早至建阳，且建瓯尚有刘和鼎所部，闽军决不能北进。以现在情形观之，彼只能守延平附近。将来在延平或有一场恶战，此一战之后胜负即决定。再进一步，即闽省善后问题耳。所可虑者，或因此引起他方之变动，及日帝国主义者之再侵扰，而共匪亦得苟延也。\footnote{《德胜关工事已完成此后匪在赣南与赣东北完全隔绝此举实其致命伤（1933年12月19日）》，《陈诚先生书信集·家书》（上），第248页。}
\end{quote}


1933年11月24日，军事委员会南昌行营制定北路军作战计划：“入闽军应以较匪优势之有力部队集中赣东，以主力猛烈压迫匪第三、第七军团，乘机推进闽北，以迅速之行动，向南进展。”\footnote{《赣粤闽湘鄂北路剿匪军第三路军五次进剿战史》（上），第五章，第4页。}12月初，进一步确定攻闽方针为：“以有力之国军一部编成数个纵队，由赣、浙边区分道入闽，先击破逆军之主力，并将其余逆部，由南北两方夹击，一举歼灭之。”\footnote{《中国现代历次重要战役之研究——剿匪战役述评》，台北，“国防部史政编译局”编印，1983，第127页。}具体攻击部署是：以第二路军两师从浙赣边界的上饶、广丰入闽，第四路军两师从浙西入闽，加上总预备队两师于12月15日前集中闽北浦城，准备分由建瓯、屏南攻击延平、水口；第五路军四个师加上总预备队一师由金溪、资溪入闽，于12月20日左右集中光泽附近，负责掩护攻击部队侧翼，并由邵武、顺昌拊十九路军之背；第三路军主动向德胜关方向出击，牵制中共部队，掩护第五路军入闽并配合其确保攻闽军右侧背安全；海军陆战队准备进攻福州、厦门。这一部署将进攻重点放在闽北方向，欲乘十九路军主力“未集中以前，迅速击破其现驻闽北之部队”，\footnote{《蒋介石致蒋鼎文电（1933年12月11日）》，《蒋中正总统档案·事略稿本》第24册，第76页。}而在闽西北由于顾忌到红军的威胁，以保持警戒状态为主。


应对红军威胁是蒋介石平定闽变不能不考虑的一个重要因素，虽然蒋判断“赤匪未必急助伪闽……必在闽北赣乐间地区，以阻止我对闽行动，而以消极助逆”，但在抽调10个师左右兵力入闽时，仍然不敢大意，自我警醒曰：“匪主力既在黎光之间，我军动作极应慎重也。”\footnote{《蒋介石日记》，1933年11月29日、12月12日。}国民党军在江西保持了强厚的兵力，留在江西及赣浙边境的第三路军辖7个师的进攻部队及6个师的守备队，加上赣西第一路军部队，兵力仍数倍于红军。对浙江后方地区，也多有部署，事变发动当天，蒋就致电浙江省主席鲁涤平，提醒其“闽乱既起，浙防应从速准备”；后又电示浙省：“在龙泉、庆元、泰顺、平阳各县对福建之松溪、政和、寿宁、福鼎之各要隘，从速派员负责，专员修筑闭锁堡并囤积粮秣，以防万一。”\footnote{《蒋中正电鲁涤平俞济时闽乱既起浙防应从速准备（1933年11月20日）》、《蒋中正电鲁涤平在龙泉等各县对福建之福县等各要隘筑堡屯粮等蒋介石致俞济时电（1933年11月26日）》，蒋中正文物档案002020200018024、002010200098057。}由于江浙一带为其基本区域，实力坚强，他甚至十分期望闽方攻浙。11月26日，蒋计划研究“如何使闽逆来攻浙”，晚间有闽方攻浙消息传来时，他更“不禁转忧为乐”。\footnote{《蒋介石日记》，1933年11月26日。}为使自身在宣传上居于主动，蒋电告陈布雷等：“自即日起即宣传闽逆进攻浙边庆元、泰顺之消息，逐渐发布使国人注意闽逆开衅之罪恶。惟宣传方法应须有系统与计划，不可使人知为虚构也。”\footnote{《蒋中正电陈布雷俞济时廿日起发布闽逆开衅消息惟宣传应有系统计划（1933年12月1日）》，蒋中正文物档案002010200099076。}深悉内幕的陈诚在家书中说得很明白：“闽逆军事行动与报纸所载完全不同，现已处被动，我军已过建瓯、邵武，即可知报载所谓犯浙，不过以祸首予闽逆耳。”\footnote{《此次胜利确足以慰劳非但为剿匪之关系实开讨闽逆胜利之基础也（1933年12月25日）》，《陈诚先生书信集·家书》（上），第251页。}12月10日，蒋介石亲向第五路军入闽先头部队训话，强调：

\begin{quote}

你们第十四军这两师人此次同走一路出发，力量非常雄厚，而这一路兵又是土匪和叛逆所料不到的，敌人一定想不到我们能够有这样一个实力雄厚的部队，由我们所决定的这个路线出去，你们这两师人的目的是要占取此后战争的中心要道，这一点对于剿匪讨逆战争最后的胜利，实有最大的关系。\footnote{蒋介石：《为闽变对讨逆军训话——说明讨逆剿匪致胜的要诀》，《先总统蒋公思想言论总集》第11卷，第625页。}
\end{quote}


随后，为防止由赣东入闽时遭遇红军阻拦，蒋先发制人，要求第三路军“由黎川向东南德胜关泰宁方向进展，协同第二路军堵截伪三、七军团，冀歼灭其实力，并竭力掩护第二路军之推进”。\footnote{《赣粤闽湘鄂北路剿匪军第三路军五次进剿战史》（上），第5章，第5页。闽变期间，国民党军序列调整较大，这里所谓第二路军，实际即由金溪—光泽一线入闽的第五路军。}希望通过攻击赣闽边境红军，将其逼向建（宁）、泰（宁）地区，敞开入闽通道，使入闽军进展顺利。根据这一计划，12月11日，黎川一带国民党军奉命沿团村向闽赣边境的德胜关地区进攻，寻找红军主力决战。16日，国民党军进占德胜关，红军被压往泰宁方向。17日，蒋介石又指示：“为防匪由泰宁绕道北窜，扰我第五路军后方起见，第五纵队应即占领金坑。”\footnote{《蒋介石1933年12月17日电》，《赣粤闽湘鄂北路剿匪军第三路军五次进剿战史》（上），第5章，第17页。}该部随即向东北方向的熊村、黄土关、金坑一线推进，截断建、泰红军往光泽一带的去路；同时加紧构筑碉堡，打通、巩固至光泽方向联络，此陈诚所谓“决先完成黎川至德胜关，及黎川至金坑、东山至熊村之封锁线”。\footnote{《电呈构筑封锁线计划（1933年12月21日）》，《陈诚先生书信集·与蒋中正先生往来函电》（上），第119页。}这样，在由金溪、资溪入闽通道东南方向，国民党第三路军构筑了一条环形防御带，确保其入闽通道安全。由于浙赣边界和浙西国民党军入闽部队本身就受红军威胁甚小，因此，当闽赣边境入闽通道基本被打通后，国民党军入闽事实上有了相当的安全保证。对于德胜关的占据，陈诚在家书中更揭示出另一层意义：“德胜关工事已完成，此后匪在赣南与赣东北完全隔绝，此举实其致命伤。尤以黎川附近之丰富资源，现被我掌握，对匪之物资补充更感困难，实可致其死命也。”\footnote{《德胜关工事已完成此后匪在赣南与赣东北完全隔绝此举实其致命伤（1933年12月19日）》，《陈诚先生书信集·家书》（上），第248页。}


随着第三路军的顺利进展及第五路军的入闽，蒋介石对红军威胁的提防逐渐减小，原计划主要用于警戒的自光泽一带的入闽部队除留一部分继续执行警戒任务外，有3个师部队投入前线，南京政府军兵力使用更为充裕。25日，蒋介石抵闽北浦城就近指挥作战，“虽一团一旅之众，亦亲临训话，砥砺士气”；\footnote{《闽省残局收束中》，《国闻周报》第11卷第6期，1934年12月29日。}同时确定攻击计划，以延平、古田、水口作为首期主攻对象。


延平、古田、水口互为犄角，是控扼福州重要外围据点，直接关系福州乃至整个闽东南地区安危。十九路军在此却只是布置了新编的谭启秀第五军，分由该军第六师守延平，第五师两个团守古田，另一团及军直属部队守水口。而蒋介石布置的围攻部队是：第四、第三十六、第五十六3个师攻延平，第八十七、八十八两师攻古田，第九师及第三、第十师各一部攻水口，仅从编制而言，就均为闽军3倍，至于实际兵力和战斗力更远远超出。因为“抽调入闽的兵力，全系蒋介石的嫡系部队，并集中了海空军及炮兵的优势力量”，\footnote{宋希濂：《我参加“讨伐”十九路军战役的回忆》，《文史资料选辑》第37辑，第109页。}而谭启秀部是十九路军原补充旅（1933年6月改为补充师）基础上成立的新军，战斗力和战斗意志都有限，这样的接战态势使闽方一开始就处于十分不利的境地。1934年1月3日，蒋介石在日记中写下其攻击计划：“微日攻击延平城，八日攻击水口，十日占领闽清。十三日占永泰。十六日占莆田。廿日占泉州、漳州。”


1月5日，战事刚一爆发，延平守军就告不支，南京政府军第三十六师、第四师、第五十六师分从城南、城东、城东北展开攻击，守军退路也被切断，被迫于次日缴械，延平易手。7日，包围水口、古田的南京政府军发起总攻，当天即占领水口。此时，位于三城犄角顶端的古田已成孤城，蒋介石对古田引而不发，欲以古田作为诱饵，围点打援。攻克水口当天，他致电前方：“蔡逆决率其主力来援古田，并言十日可达古田附近，逆军出巢来犯，正我军求之不得者，现决对古田城逆暂取包围监视之姿态，不必猛攻。”\footnote{《国民党陆军第八十八师古田围城之役战斗详报》，《中华民国史档案资料汇编》第5辑第1编军事（5），第764页。张治中在回忆录中一力说明是他冒着违背蒋的命令风险坚决主张缓攻古田以争取守军投降（《张治中回忆录》（上），文史资料出版社，1985，第92～95页），证之上述电报，似不可靠。}次日，蒋再次强调：“古田城逆，只可包围，昼夜佯攻，一面严密监视，不许其逃遁，亦不必留缺口，但不可攻破，务使蔡逆主力仍来增援古田。”\footnote{《国民党陆军第八十八师古田围城之役战斗详报》，《中华民国史档案资料汇编》第5辑第1编军事（5），第764页。}9日，当十九路军一部前出准备北上往援古田守军时，蒋介石更信心满满地指示卫立煌：“逆军已倾全力来犯古田水口之线，刻已进至白沙以西地区，望兄迅速准备在白沙洪山桥间地区，选择多数之渡河点，设法渡河，袭击逆军侧背，整个包围而歼灭之。”\footnote{《国民党军卫立煌部镇压“闽变”战斗详报》，《中华民国史档案资料汇编》第5辑第1编军事（5），第833页。}其一心期望以古田为诱饵，吸引十九路军主力于古田、水口地区实施歼灭。但是，十九路军并没有足够勇气与蒋在闽北对垒，1月12日，北上往援的沈光汉部与南京政府军第三师稍有接触，虽然政府军根据蒋的指示后撤诱敌，但沈部并未乘势前进，反而见其“不战而退，更致狐疑”，\footnote{《国民党陆军第三师参加镇压“闽变”战斗详报》，《中华民国史档案资料汇编》第5辑第1编军事（5），第747页。}当晚即向白沙方向退却。其实，根据蔡廷锴的回忆，1月9日，蔡与陈铭枢、蒋光鼐等放弃福州向闽南撤退，这时的应援行动更多的只是一种姿态。12日，孤处敌后的古田守军投降。


上述电文、战报体现出的是蒋介石一意诱敌而十九路军不敢应战的过程，此中的蒋介石可谓运筹帷幄、信心满满。然而，如果对照蒋介石日记，却会发现档案、电文中无法反映的另一面，看到作为一个人的蒋介石心态的复杂变幻。战事爆发后，蒋在日记中对福建方面是否将主力出福州，在闽西主动出击一直高度关注。就军事常识言，单纯防御福州几无可能，水口、古田为福州防御必守之地。因此，蒋在全面攻击即将展开时，密切注意：“我军攻击水口时，逆部主力由省来袭否？”\footnote{《蒋介石日记》，1934年1月14日。}1月6日攻克延平后，他全力注视闽方军事动向，猜测“福州逆军，其或反守为攻乎？”\footnote{《蒋介石日记》，1934年1月6日。}当时，由于担心日方借事变有所动作，蒋对进攻福州没有信心，多次在日记中写道：“对倭只有避战，如不得已，则不攻福州，以延平为省会，成立政府”；“对福州叛逆，如果集中负隅，则以封锁之法处之”。\footnote{《蒋介石日记》，1933年12月29日、1934年1月1日。}如果闽方出福州在闽西决战，对蒋而言，不失为一个一举解决闽省的机会，前文中说到的蒋的诱敌之计盖出于此。但是，十九路军的战斗力毕竟不可小觑，两军正面交锋，成败也未可料。所以，1月7日，当蒋得到事后证明并不确切的消息，报告“蔡逆果率主力来援古田”时，当时的反应却并非“吾计已售”的得意，而是令览史者感慨万端的“喜惧交集”。\footnote{《蒋介石日记》，1934年1月7日。}喜的自然是闽方出击，其计可售；惧的则是决战结果，事前难有绝对把握。


古田不守，福州外围防御据点尽失，南京政府军开始向福州推进。与此同时，南京政府海军早在12月下旬就先后占领长门、马尾两要塞，时时威胁福州安全。1月9日，海军在厦门市长黄强配合下接收厦门，威胁漳州地区，对十九路军后方形成巨大威胁。在不利形势下，闽方向南京提出三项停战条件：“海军守中立”、“中央军在距福州十英里之线停止”、“福州治安交由海军陆战队接收”，\footnote{《汪兆铭电蒋中正闽逆托日方提出停战三条件（1934年1月12日）》，蒋中正文物档案002020200018117。}欲以此退让换取蒋介石息兵。对此，汪精卫认为“逆军能如限撤退，如能做到仍为有利”，\footnote{《汪兆铭电蒋中正闽逆托日方提出停战三条件（1934年1月12日）》，蒋中正文物档案002020200018117。}但信心满满的蒋介石根本不为之所动，反而加紧对十九路军的攻击。四面楚歌声中，十九路军撤出福州，向闽南退却。16日南京政府军进占福州。


在进攻福州外围据点同时，蒋介石已经开始部署从闽西北插向十九路军后方。1月4日，战事尚未打响，他在日记中标列的注意事项就有“进取闽南利害之研究”。\footnote{《蒋介石日记》，1934年1月4日。}7日，鉴于水口已下，蒋介石考虑：“卫第五纵队挺进闽南计划是否实施，当注意之。”\footnote{《蒋中正总统档案·事略稿本》第24册，第174页。}次日，电卫立煌令其分兵南下永泰，“但须隐秘中央军兵力队号为要”。\footnote{《国民党军卫立煌部镇压“闽变”战斗详报》，《中华民国史档案资料汇编》第5辑第1编军事（5），第830页。}10日，再电卫立煌，令其渡河南下，行动须守秘密，“不可使逆军发觉我有渡河企图”。\footnote{《蒋中正电卫立煌所部在闽清须极端隐秘（1934年1月10日）》，蒋中正文物档案002020200018113。}永泰地处福州西南部，由此前进可扼住十九路军退路，蒋的一系列动作旨在于此。因此，蒋介石日记明确记有：“逆军如向闽南撤退，则第五路仍照原计划向永春、漳州急进。”\footnote{《蒋介石日记》，1934年1月10日。}12日夜，在确知十九路军将全线后撤时，蒋介石命令“主力明日速向永泰急进。除留一旅守永泰外，其余主力再向仙游沙溪急进，以行截击”。\footnote{《国民党军卫立煌部镇压“闽变”战斗详报》，《中华民国史档案资料汇编》第5辑第1编军事（5），第835页。}如计占领永泰后，蒋介石大感得意，在日记中写道：“本日我军已占永泰，此心为之大慰，从此必可如计截击，在莆田海滨歼敌，使之片甲不返也。”\footnote{《蒋介石日记》，1934年1月14日。}


由于蒋在准备围点打援、诱敌实施歼灭战同时，已有展开追击战的腹案，因此，当十九路军沿着沿海公路南撤时，南京政府军从侧翼对十九路军展开所谓“行动之艰苦与神速，俱达极点”\footnote{《国民党军卫立煌部镇压“闽变”战斗详报》，《中华民国史档案资料汇编》第5辑第1编军事（5），第862页。}的超越追击，东路军总司令蒋鼎文指挥四路大军以莆田、仙游、安溪、同安、漳州等为目标，直插十九路军后方。蒋介石要求前方将领：“我军只要正面稳固，尽可多抽部队，到达惠安以东或以西地区，分组截断公路，节节袭击，横断其退路，总须达成一网打尽之目的，以为我战史创例也。”\footnote{《国民党军卫立煌部镇压“闽变”战斗详报》，《中华民国史档案资料汇编》第5辑第1编军事（5），第846页。}在南京政府军快速推进下，全线溃退的十九路军不断遭到追击部队的堵击，狼狈不堪。17日，南京政府军第八十三师已进至仙游，次日，第九、第十师到达。南京政府追兵和夺路而逃的十九路军在仙游、涂岭一带激烈交锋，虽然南京政府军未能在此完全堵截十九路军并予以消灭，但十九路军“蒙受巨创，士气沮丧，致入于不堪再战之境地”。\footnote{《国民党军卫立煌部镇压“闽变”战斗详报》，《中华民国史档案资料汇编》第5辑第1编军事（5），第863页。}20日，莆田被南京政府军占领，十九路军大部纷纷向泉州退却。21日，南京政府军第三师由厦门嵩屿登陆，对泉州一带的十九路军形成南北夹击态势。蔡廷锴见大势已去，被迫离开部队，所部随即向蒋介石请降，轰轰烈烈的闽变从大规模交战开始到失败不过半个月时间即告瓦解，确如军事发动前蒋介石所言：“闽乱不逾一月，必可敉平。”\footnote{《蒋介石致汪精卫电（1933年12月11日）》，《蒋中正总统档案·事略稿本》第24册，第80页。}


迅速镇压福建事变后，蒋介石顺利将福建纳入自己手中，对其“剿共”军事的继续展开，大有裨益。

\section{中央苏区核心区域的争夺}

\subsection{蒋介石与陈诚关于进攻路线的讨论}


蒋介石发起第五次“围剿”，初期进攻重点放在江西方面，以广昌、石城、瑞金这一直通苏区中心区的进攻路线作为主攻方向。福建事变后，随着国民党军第二路军大批进入福建及红军主力一部分向福建进军，国民党方面在主攻方向上有过动摇。


1934年1月初，鉴于福建事变后国民党军在福建进展顺利，红军北出永丰地区的行动又不顺利，中革军委改变将红一、三军团组建为西方军的设想，决定加强福建方面兵力，牵制国民党军在福建的军事行动，命令红三军团进兵福建。红三军团入闽后，主动出击，很快于11日包围闽西沙县县城。同时，红七军团也从泰宁南下配合行动。红军首先打败了赶来增援的国民党军第四师，取得第二次入闽的初战胜利。25日，攻占沙县，毙伤敌军700余人，俘虏1300余人。红军的行动，引起蒋的很大不安，1月中旬，蒋要求第三路军以主力进军福建建宁，强调：“福州不难收复逆军必可在兴化海滨为我全部解决，我北路第三路军主力应即设法迅占建宁或泰宁，并限本月杪达成任务。”\footnote{《蒋中正电顾祝同等福州不难收复逆军必可在兴化海滨为我全部解决（1934年1月16日）》，蒋中正文物档案002020200018119。}同时以一部留守现地，再将“第三十六军速向西移，以第四军、第三十六军、第九十二师、第九十三师编为一路，向沙溪、龙岗进展”，\footnote{《蒋中正电顾祝同陈诚伪第五军团在沙县附近速率第三路向泰宁兼程前进（1934年1月16日）》，蒋中正文物档案0202000200002，台北“国史馆”藏。}对兴国方向保持警戒并相机发动进攻。由于第三十六军辖第五、第六、第九十六3个精锐师，主力东开，三十六军西移，将使第三路军形成东、中、西三个兵力点，造成严重的分兵局面。对此，陈诚提出意见：“第三十六军不必西移”，认为该军“一时难以抽调，俟第三路军主力占领建宁后，该军可推进至康都，协同完成南丰康都建宁封锁线，并相机策应主力，向广昌进展”，坚持以广昌为主要进攻方向。同时，陈不赞成由龙岗一带向兴国进兵，表示部队“只能推进至富田，如向兴国进展，由泰和经沙村为妥”。\footnote{《赣粤闽湘鄂北路剿匪军第三路军五次进剿战史》（上），第六章，第2～3页。}对此，蒋经过慎重考虑后同意了陈的建议，指示顾祝同：“预定三十六军及四十六师、十三师将来之进展计划均照所拟办理。”\footnote{《顾祝同1934年1月26日致陈诚转告蒋介石命令电》，《中华民国史档案资料汇编》第5辑第1编军事（4），第26页。}改变了原来的分兵计划。


1月底，鉴于福建境内对十九路军的战事已基本结束，红军虽有攻占沙县的动作，但以入闽部队对付之，应该可保无虞，陈诚致电蒋介石，建议第三路主力部队留在江西境内，巩固黎川东南战线：

\begin{quote}

樟横村为黎川通建宁之咽喉，西城桥为南丰入建宁之孔道，两者均为伪赣闽省苏之中心区域，实匪必争之地……军为巩固后方并节约兵力起见，拟先进占西成桥，构成石沟圩西成桥樟村间之封锁线，将钟贤东坪西坪一带守兵推进接守，届时取建宁或入广昌均较为易。\footnote{《陈诚电蒋中正顾祝同拟先进占西城桥构成封锁线俾取建宁或入广昌（1934年1月29日）》，蒋中正文物档案002090300069260。}
\end{quote}


2月1日，陈诚再电蒋介石，提出：“樟横村附近粮食被匪搜刮，给养困难万分。待封锁线完成，给养略事补充，即向建宁前进。”\footnote{《陈诚电蒋中正待黎川经樟横村等至西成桥各封锁线完成补充给养后即向建宁前进（1934年2月1日）》，蒋中正文物档案002080200146025。}2日，蒋介石回电仍然坚持令陈诚“设法早占建宁或泰宁”。\footnote{《赣粤闽湘鄂北路剿匪军第三路军五次进剿战史》（上），第六章，第17页。}同日，陈诚致电蒋介石，以将在外而君命有所不受的姿态坚决主张：“此间决先解决当面之匪，再用交互前进法，向建宁进展”，\footnote{《陈诚电蒋中正（1934年2月2日）》，蒋中正文物档案002090300069328。}坚持主力留在黎川东南作战。从战略态势看，由于福建十九路军已被击垮，红军有限兵力不足以影响福建全局，此时离开主攻方向向福建增兵颇有不得要领之感，而黎川作为战略要地，是双方进退的要点，一旦黎川丢失，则其已经获得的先机将一朝尽失。因此，蒋介石坚持进占建宁的主张颇有昧于大局、轻重倒置之嫌，而陈诚的主张显然更具主动意义。所以陈诚强调：“此后我如得一地一城，即不再为匪所有，保障民众不为匪用，虽觉迟缓，实万全之策。”\footnote{《陈诚电蒋中正（1934年2月2日）》，蒋中正文物档案002090300069328。}面对陈诚的抗命举动，蒋介石不以为忤，权衡再三后，次日复电陈诚，肯定其“所见甚是”，同意“照办”。\footnote{《赣粤闽湘鄂北路剿匪军第三路军五次进剿战史》（上），第六章，第17页。}


但是，蒋介石并不甘于自己的计划被搁置，2月11日，蒋致函陈诚，提出：“中前日回赣审察匪情，其目的仍欲先求闽北方面之决战得胜，然后再回攻我赣边之师，故此时我赣边部队，对其建泰方面暂勿进取，以其工事已成且有备也，不易垂手而得。不如先占康都与广昌，使匪不注重我主力之东向，以懈其戒心，俾其主力转移于沙县方面，专力对我第四、第八十师时，然后我进取广昌、康都部队之主力，再向东进占建宁，必易得手……故弟路务于删日以前占领康都与广昌，最好同时占领，待至二十五日以前，再转移至建宁与康都间地区，俾与我闽北部队夹击残匪，至转移兵力时，广昌最多留驻三团或竟放弃亦可。”\footnote{《手谕务于十五日前占领康都广昌（1934年2月11日）》，《陈诚先生书信集·与蒋中正先生往来函电》（上），第125页。}蒋此函表面同意陈诚部队留在江西作战，但其重心仍放在福建方面，而且限几日内攻占广昌也让人不知所云，和国民党军一直奉行的稳扎稳打方针明显抵触。正因如此，此函实际并未发出，陈诚2月12日家书中提到：“昨蒋先生给我一函，因计划稍有变更，仍取回未送来。惟不知此后进剿计划如何耳？”\footnote{《我已拟一计划寄蒋先生如能照此进剿当万无一失也（1934年2月12日）》，《陈诚先生书信集·家书》（上），第270页。}当是指的此函。蒋介石计划改变，应和其过于冒进有关。2月12日，蒋在日记中表示：“进剿不能过急”，实际是对陈诚意见的首肯，同时详细写下其“进剿”计划：“预定三月十日前第二期作战开始，三月底到达连城、宁化、广昌、龙岗之线；四月底或五月十五日到达汀州、宁都、兴国之线；六月底占领瑞金、于都，肃清残匪。”\footnote{《蒋介石日记》，1934年2月12日。}几乎与蒋写下这段日记同时，陈诚也起草了给蒋介石的报告，详论其“围剿”计划：

\begin{quote}

我军为始终立于主动地位，而达收复一地一城即不再落匪手，保障民众今后永不为匪所用之目的，应本钧座战略取攻势，战术取守势之原则，以罗军长封锁围进，乘机突进之方法，第一步本路军先完成西城桥至石沟圩封锁线后，即进占宁都，构筑由西城桥向西延伸之碉堡线，截断匪南丰与建宁之联络，并策应我东路军收复建泰。待东路军确实占领建泰后，掩护本路军取广昌，同时各部仍赶筑碉堡，俟完成后，以刘和鼎部分守将乐泰宁，以卢兴邦部守沙县，以陈明仁部守延平顺昌一带，以周志群守邵武光泽，以黄子咸守黎川建宁，同时令薛岳部相机进占龙岗，李云杰部出富田。第二步东路军一部出归化，掩护一部主力由永安出，本路军出宁都。第三步东路军一部取石城，主力出长汀，而本路军相机占领兴国及瑞金。最后会师，寻匪主力而围歼之。如此筑碉堡交互推进，我碉堡完成之日，即剿匪成功之时。虽似觉缓慢，实万全之策。\footnote{《函呈进剿闽赣赤匪意见（1934年2月12日）》，《陈诚先生书信集·与蒋中正先生往来函电》（上），第126页。}
\end{quote}


如果单将这一报告和蒋介石同日日记联系看，很容易得出蒋被陈诚影响改变作战计划的结论，但陈诚在写给妻子的家书中透露，这一报告当天并未送出，陈在12日写道：“我已拟一计划寄蒋先生作其参考，拟明日请派飞机来勾去。如能照此进剿当万无一失也。”\footnote{《我已拟一计划寄蒋先生如能照此进剿当万无一失也（1934年2月12日）》，《陈诚先生书信集·家书》（上），第270页。}可见，蒋作出决定时，并未读到陈的报告。蒋介石和陈诚殊途同归，几乎同时作出在江西缓步推进的决定，这在蒋陈二人而言，应为洞悉大局的明智之举。


2月15日，陈诚致电蒋介石，详论今后作战方针：“本路军以出广昌、宁都，将匪区截成两段为有利，盖我军如能占领广昌、宁都，既可防匪西窜，并可策应我西区，且可吸引匪之主力，而使我东路军进占建泰容易也。我军无论出建宁或广昌，第一步均须先将当面之伪一、五、九军团，及伪三军团之一部，驱逐或击溃之。”\footnote{《电呈进剿方略并督饬所部达成任务（1934年2月15日）》，《陈诚先生书信集·与蒋中正先生往来函电》（上），第128页。}陈诚的意见，得到了蒋介石的支持，此后，国民党军的主攻方向始终保持在广昌、石城一线，即所谓“剿匪战略重西轻东，左急右缓”。\footnote{《蒋介石日记》，1934年3月6日。}福建方面则由东路军“由北向南”，\footnote{《蒋中正电卫立煌闽省进剿计划决先由北向南（1934年2月24日）》，蒋中正文物档案0202000200011。}顺势推进。在国民党军进攻压力下，2月中旬，红三军团又奉命返回江西广昌一带，国民党军再次抢得先机。陈诚在国民党军战略制定上发挥着重要作用，不过，对作战方针的坚持也使陈诚承受了极大的压力，“时陈诚对当前形势及所担负任务，愁虑重重。据他随从副官透露，陈常深夜不眠，在办公室内徘徊，忧心忡忡”。\footnote{陈振锐：《在陈诚所部参加第五次“围剿”见闻》，《福建文史资料选辑》第2辑，第31页。}


在国民党方面计划下一步作战方针时，中共方面也在评估数月来反“围剿”战争的结果。虽然没有实现歼灭国民党军有生力量的目标，但国民党军在苏区外围的徘徊也颇让中共领导人引以为慰。当时《红星》发表的社论写道：红军“在黎川之东北，进行了三个月胜利的战斗，伸展到浒湾附近，击溃敌人的第四师，在硝石、洪门、新丰街一带，发展了广大的游击战争，取得了很多的胜利。在北面战线，我们的大兵团，可以自由的出现于敌人堡垒线的后方，许多独立兵团，在敌人的后方出没与发展”。社论把国民党军加强构筑堡垒的举动归因于红军的主动攻击，强调正是由于红军对国民党军的打击，使之“不得不改变其计划，相隔二里路甚至四百米，就要构筑一堡垒。而且他的堡垒要向着南昌来构筑。这种堡垒政策的结果，又使其攻击精神的减弱”。\footnote{《五次战役中我们的胜利》，《红星》第25期，1934年1月21日。}作为公开的舆论宣传，这种说法有其进行正面阐释的可理解一面，但仍显得过于避重就轻，缺乏准确的反思和评估。


关于对第五次反“围剿”进程如何看待的问题，在二苏大上也有反映。朱德在二苏大报告中谈道：“在最近，我们的英勇红军，已经把帝国主义国民党的六次‘围剿’粉碎了一半，给了敌人严厉的打击，阻止敌人向苏区前进。”\footnote{朱德：《向大会致粉碎六次“围剿”的敬礼》，《苏维埃中国》，第240页。}毛泽东则针对大会上认为“围剿”已经粉碎和仅在准备粉碎中的两种说法指出：

\begin{quote}

照前一说，是过分估计了自己的胜利，把苏维埃最后粉碎“围剿”的严重任务轻轻取消了，而实际上，蒋介石正在集中一切力量最后向我们大举进攻，所以这种估计是不正确的，并且是非常危险的。照后一说，是看不到几个月来红军从艰苦战争中已经给了敌人以相当严重的打击，已经取得了第一步胜利，这种胜利，同粉碎五次“围剿”的伟大胜利合起来，就成为我们粉碎六次“围剿”的坚固的基础。对于自己成绩估计不足，同样是很危险的。\footnote{毛泽东：《关于中央执行委员会与人民委员会报告的结论》，《苏维埃中国》，第304～305页。}
\end{quote}

\subsection{平寮、凤翔峰战斗}


1934年2月20日，为部署推进第五次“围剿”后期计划，蒋介石在南昌召集顾祝同、陈诚、熊式辉、陈调元及西、南两路将领举行军事会议。21日，南昌行营重新调整兵力部署，将入闽军队改编为东路军，委任蒋鼎文为东路军总司令，率第二路军与第五路军及预备队共16个师又1旅2团，向中央苏区东面的建宁、泰宁、龙岩、连城等地推进，目标是夺取广昌及中央苏区中心地长汀和瑞金，协同已组成的北、西、南三路军，形成对中央苏区四面合围之势。


此前，1934年1月下旬，国民党军第三路军已经开始向黎川以南发动进攻。初期根据蒋介石的指示，以建宁为攻击方向，于25日攻占樟横村。樟村、横村是黎川南50里左右的大村，“两处地方极大，约有三、四千户，如不遭匪患，实非常之富”。\footnote{《此次出发占领建宁大概不成问题此或匪对我生日所赠礼物也》，《陈诚先生书信集·与蒋中正先生往来函电》（上），第259页。}随后，国民党军向由江西进出福建的要隘邱家隘、寨头隘进犯。邱家隘为闽赣交界的分水岭，“两翼系高峰，中为一谷地，道路蜿蜒曲折其间，加之构筑一些工事，更为壮严之天险”；“冬天积雪未溶，山头竹林茂密，真是易守难攻之地”。\footnote{《陈伯钧日记（1933～1937）》，第161页；曹嘉忠：《邱家隘战斗》，《建宁文史资料》第1辑，政协建宁文史资料组编印，1982，第41页。}红五军团两个师部队在此凭借天险，面对对方优势兵力“顽强抵抗”，\footnote{《赣粤闽湘鄂北路剿匪军第三路军五次进剿战史》（上），第六章，第6页。}国民党军不得不集中七倍于红军的兵力，经过反复十数次冲锋，付出重大伤亡后才将红军迫出战场。随之，又占领寨头隘阵地，并乘势进占无人驻守的平寮。


2月1日，红军向平寮一带国民党军发动反攻，计划“消灭进占坪了附近之敌并相机占领里岭下向樟村攻击，消灭敌之增援部队”。\footnote{《坪了附近战斗经过情形报告（1934年2月1日）》，《中央苏区第五次反“围剿”》（下），第53页。}红一军团第一、第二两师主攻平寮，“来势殊为涌猛”，\footnote{《国民党北路军顾祝同部与中央苏区红军作战情形报告书》，《中华民国史档案资料汇编》第5辑第1编军事（4），第44页。}五、九两军团在邱家隘、寨头隘一带牵制对方部队。平寮之战，据国民党方面战史记载“战事激烈前所罕见”。其所记战斗经过为：“匪惯以小部队牵制各处，主力集结于一地，用密集队形重叠冲锋，如此平寮战役第一二两日，均以小部队牵制我寨头隘阵地，主力集结于平寮前方……其攻击部署，犹是率由旧章，但我工事坚固，士气旺盛，匪亦无所施其技矣。”\footnote{《赣粤闽湘鄂北路剿匪军第三路军五次进剿战史》（上），第六章，第14～15页。}经过两天攻击后，红军虽予国民党守军以重创，歼其七十九师四七○团大部，但国民党军在坚固工事和飞机轰炸配合下，利用占据“坪了至里岭下一带高地”\footnote{《坪了附近战斗经过情形报告（1934年2月1日）》，《中央苏区第五次反“围剿”》（下），第53页。}的有利地形顽强抵抗，红军也遭遇较大伤亡，陈诚家书载：“据俘匪供称，匪之第二师生还者不及十分之一”，\footnote{《匪第一军团攻我樊崧甫师第九军团攻霍揆彰师均被我击溃》，《陈诚先生书信集·与蒋中正先生往来函电》（上），第262页。}此虽或有夸张，但红军损失重大应属事实。因未能达到击垮、歼灭国民党军有生力量的目标，红军被迫于2日午后撤出战斗。


平寮之战后，根据蒋介石的想法，第三路军应继续向福建挺进，但陈诚认为福建战事已基本结束，此后进攻方向应直接指向苏区的中心区域广昌，这一进攻方向很快得到蒋介石的认同。2月初，陈诚下达作战命令，令所部在樟横村一带集结，构筑工事，主力向位于黎川西南地区的西成桥方向推进，揭开第三路军向广昌进攻的序幕。


为破坏国民党军继续南进的意图，红军在西成桥西南、西北方向骚扰国民党军，阻止其修筑碉堡、工事，破坏其占领意图，双方在鸡公山、司令岩等高地展开拉锯战。9日，国民党军十四师及九十四师占领鸡公山地区，红军退往“康都五通桥”方向。\footnote{《本师西成桥附近之役战斗详报》，《汗血月刊》第1卷第4期，1934年4月20日。}10日，红一军团第一师第一团二营增援鸡公山途中，在三甲嶂顽强抗击国民党军绝对优势兵力的进攻，同时第一师其他部队根据军团指挥部命令“从西而东的向敌之后路打出去”，击溃其数个师的进攻，创造了以弱抗强的范例，被誉为“战斗中最光荣的模范”。\footnote{《鸡公山附近战斗经过详报（1934年2月11日）》，《中央苏区第五次反“围剿”》（下），第61～62页。}聂荣臻撰文提出：“这应当写在我们的战史上，成为我们第一团三甲嶂上光荣战斗的一页。”\footnote{聂荣臻：《把第一团顽强抗战的精神继续发扬光大下去》，《红星》第31期，1934年3月4日。}但这样的胜利不足以影响大局。


15日，国民党军第七十九师进至乾昌桥一带，准备在此筑碉。红一军团根据中革军委命令，决定在俯瞰乾昌桥的凤翔峰左、右两翼埋伏红一、二两师部队，红四师在敌侧后“负责钳制79师之后续部队及敌之11师策应樊师之部队”；同时以二师一个团“兵力固守凤翔峰阵地并吸引敌人主力”，待敌仰攻凤翔峰而抵前沿时，左、右两翼“突击截断敌之归路，以迅速最短的时间歼灭”。\footnote{《一军团凤翔峰战斗详报》，《中央苏区第五次反“围剿”》（下），第66～67页。}凤翔峰地当要冲，为国民党军之所必攻。15日上午，国民党军第七十九师以两个营部队向凤翔峰展开进攻，当其接近阵地时，红军展开突击，但因相互间协调不够，未能歼灭对手。下午16时左右，国民党军增调第六十七、第八师等部向红军发起更加猛烈的攻击，同时出动飞机对红军阵地展开轰炸，而其6个团兵力绕向红军后方，“有抄袭我归路模样”。\footnote{《一军团凤翔峰战斗详报》，《中央苏区第五次反“围剿”》（下），第69页。}红军全盘计划落空，被迫撤出战斗。凤翔峰战斗次日，林彪、聂荣臻曾就战斗情况作出总结：

\begin{quote}

凤翔峰战斗又重复了三甲嶂的血训，敌又是用飞机轰炸及堡垒炮兵火力向我阵地轰击，尤其是敌人的飞机所起的作用太大，使我一举一动甚受妨碍。敌步兵在其空、地火力掩护之下，向我阵地推进，我军以反突击将其击溃后，敌缩回堡垒附近从容整顿。我军因被敌方堡垒所阻，而不能行较长距离的进击，行反突击亦只能做到向敌之侧面搏击及正面冲锋，而无法做到压迫到敌之后面……我们时刻感觉得敌人堡垒外的近距离或到堡垒间隙中去求运动战，结果仍变成堡垒战。以大部队在这种场合想行短促而突然的袭击，结果打响之后仍然不易拉（摆）脱而成了对峙局面。这种战斗办法最好的成绩也只能消灭敌之一部，否则仅仅将敌击溃而不能消灭。但这种战斗，我们他们的伤亡和弹药的消耗都很大。我们在五都在丁毛平寮在三甲嶂和昨日的战斗，性质和结果都差不多，平寮三甲嶂及昨日的战斗，我们伤亡在约一百人以上，因此觉得以后须尽量的避免这样的战斗。在距敌人半天或一天行程处隐蔽，应于敌前进而于敌之运动中或初到时消灭之……我们只要敌离开其（堡垒）八里到十里路以外，我们便能完全集中的消灭敌人。虽敌数倍我军的集团兵力，只要敌后尾离开其后面堡垒在十里以外，我们是有把握消灭他的。\footnote{《凤翔峰战斗经过及经验教训（1934年2月16日）》，《中央苏区第五次反“围剿”》（下），第63～64页。}
\end{quote}


凤翔峰战役中，还出现一个对红军相当不利的现象，即红军行动不再像先前那样可以充分隐蔽，相反，资料显示，该地“封建关系浓厚，故仍有仇视反对或破坏我们现象与事实。特别是南丰敌派遣其走狗或利用当地流氓密布侦探网，对我军行动消息往往容易漏网”。\footnote{《一军团凤翔峰战斗详报（1934年2月20日）》，《中央苏区第五次反“围剿”》（下），第66页。}这样的状况，对红军是一个巨大的隐忧。


在江西方面作战取得进展后，国民党军在东路福建也展开进攻战。2月22日，东路军第十纵队汤恩伯部夺回沙县。3月7日，又攻占将乐，进逼泰宁。3月19日、23日，泰宁、归化相继失守，中央苏区闽赣间进出的门户建宁受到严重威胁。中革军委紧急调红一、三军团主力入闽，在大洋嶂、将军殿、梅口、江家店等地与进犯敌军展开激战，大洋嶂一战，成功阻截国民党军第八十九师的进攻，消灭其一个团兵力，尸横遍野，以致“将近半个月的时间，臭气难闻，路上无人往来”。\footnote{梁印浩：《大洋嶂阻击战侧闻》，《泰宁文史资料》第1～3辑合订本，政协泰宁文史资料委员会，1990，第38页。}时任红三军团政治部主任的袁国平称誉：

\begin{quote}

以两连兵力迎拒十数倍于我的优势故人，激战终日，越战越勇，营长、连长、指导员及代营长连长、代指导员继续伤亡二次，仍无退却动摇，坚忍卓绝的与敌人血战到底。敌旅长亲率大刀会队督战，并不惜以敌兵死尸及垂死重伤敌兵堆成进攻工事，拼死命地向我猛进，但是终被我英勇的五、六连以顽强拒战与神勇的反冲锋打得落花流水的溃败下去。\footnote{袁国平：《大洋嶂战斗》，《泰宁文史资料》第1～3辑合订本，第177页。}
\end{quote}


红军的英勇奋战延缓了国民党军在福建的进攻步伐。不过，福建本非国民党军主攻方向，这里的战局不足以对全局形成影响。

\subsection{广昌外围战斗}


西成桥战斗后，国民党第三路军摆出欲由西成桥经康都斜向进攻广昌的态势，吸引红军在此防御，主力则逐渐集中到南丰西南地区，准备经白舍南下直插广昌。2月25日，蒋介石令第十八军“主力先向杨林渡白舍罗坊伸展”。\footnote{《蒋中正电顾祝同等第十八军占领荷田冈后主力先向杨林渡白舍罗坊伸展（1934年2月25日）》，蒋中正文物档案0202000200013。}3月初，国民党军第三路军在南丰西南地区集中11个师部队，在沧浪—杨林渡一线构筑碉堡线，准备以此为基地，向南推进。


国民党军完成在南丰西南地区的集结后，红军逐渐察觉其意图。为阻止其对广昌的进攻，红军集中红一、三、九军团北进，向南丰西南地区出击，“有在白舍、三坑、三溪之线与敌决战之决心”。\footnote{《陈伯钧日记（1933～1937）》，第184页。}此时国民党军已经完成在此地区的战略展开，6个师的主攻部队集中在南丰以南方圆20公里左右区域内，声气相通，配合便捷。陈诚确定应战计划：

\begin{quote}

伪三军团全部在我五都寨东华山石股山阵地前，伪一军团由早坳向我侧翼迂回。明真日部署：霍李傅三师固守原阵地，夏黄两师集结立壁岭牛形岭茅坪间地区，樊师主力集结磨刀渡附近，伺机以击破之。而令孔师主力星夜向神冈党口前进，威胁匪之侧背。\footnote{《陈诚电蒋中正顾祝同（1934年3月10日）》，蒋中正文物档案002090300071188。电文中提到的霍、李、傅、夏、黄、樊、孔分别为霍揆彰、李树森、傅仲芳、夏楚中、黄维、樊崧甫、孔令恂。}
\end{quote}


3月11日，红军向国民党军发起全线攻击，占领五都寨、东华山阵地，并继续向北攻击。主力则由西南方向插向国民党军右后方，准备侧击其阵地右翼，断敌后方。但是，国民党军在此厚集兵力，红军初期进攻虽有得手，总体形势并不乐观。国民党军在红军全线攻击下，一方面顽强固守，另方面积极准备利用兵力集中优势，全线发动反扑，蒋介石致电前方提醒“伪一军团踪迹不明，希注意详侦，防其暗算”同时，指示各部主动出击，“相机猛击其侧背……设法歼灭之”。\footnote{《蒋介石文午行战电（1934年3月12日）》，《赣粤闽湘鄂北路剿匪军第三路军五次进剿战史》（上），第七章，第10页。}陈诚也亲赴前线督战，通电全军，“亟望各级官兵，自勉勉人，贾我余勇，不顾一切，急起直追”，\footnote{陈诚：《勖勉前线将士》，《赣粤闽湘鄂北路剿匪军第三路军五次进剿战史》（上），第七章，第14页。}激励士气，向红军发动反攻。


3月13日，红军向杨梅寨国民党军后方部队第九十八师发起进攻。黄昏，国民党军各部全线反击，“一时战事骤转激烈，各阵地炮火之猛为剿匪以来所未有”。\footnote{《国民党北路军顾祝同部与中央苏区红军作战情形报告书》，《中华民国史档案资料汇编》第5辑第1编军事（4），第46页。}红军由于部队深入，处在国民党军环形阵线之中，在国民党军并不多见的反击下，准备不足，虽努力奋战，仍遭失利。彭德怀鉴于战场形势不利，建议中革军委“应有舍心”，\footnote{《彭德怀关于作战等问题给军委的信（1934年4月1日）》，《中央苏区第五次反“围剿”》（上），第154页。}主动撤出战斗。22时，第一军团第一师接到命令：“东华山五都寨均到敌，我军决转移地区，一师担任最后掩护，候彭杨部队（第四师）及我军完全通过后，再节节掩护到宝石朱坊附近宿营。”\footnote{《五都寨附近战斗详报（1934年3月10日至14日）》，《中央苏区第五次反“围剿”》（下），第80页。}红一师奉命掩护部队撤退，到次日晨5时，红军主力大部撤出战场。此役，红军歼灭对手1200多人，自身损失达1700多人。\footnote{参见《杨尚昆回忆录》，中央文献出版社，2001，第89页。国民党方面对这次战役伤亡情况的描述是：“匪伤亡约在七、八千以上，阵前死尸约有三千余具……我方亦伤亡近千。”见《南昌行营第二厅致汪精卫等电》，《中央革命根据地革命与反革命斗争史料》（下），中国社会科学院近代史研究所藏。}国民党方面战史后来总结此役时说：

\begin{quote}

匪军常集中其大部，对我一点，施行猛烈攻击，以求其突破或包围之成功，此种时机，匪军弱点，即在处处不能对我主攻，故处处均感虚薄。此次匪以伪一三军团主力，猛攻我杨梅寨夏师阵地时，我全线不顾一切，果敢出击，致使匪虚弱部分，均被我击破，其主攻部分，亦被我截断，全线遂不能不总崩溃。\footnote{《赣粤闽湘鄂北路剿匪军第三路军五次进剿战史》（上），第七章，第24页。}
\end{quote}


陈诚此次战斗中采用的反击和跃进战术给林彪留下了深刻印象，他当时数次说起，国民党军“前进的距离和速度，是随着群众条件，随着对我主力行踪的了解，随着他本身兵力的大小，随着当时红军打击他轻重的程度，以及地形情形等条件而变动的……陈路军向广昌前进时，他虽然发现红军在其附近，然因他已集结了充分的优势的兵力，同时在地形上也容许他大兵力的使用，所以他依然前进。”\footnote{林彪：《短促突击论》，《革命与战争》第6期，1934年7月。}彭德怀则具体谈道：

\begin{quote}

是役敌是纵深的，紧靠着南丰数年来所构成的坚固堡垒，这种情形无条件的突击是徒然剥削了自己的力量，得不到任何代价。东华山、五都寨的教训不是不了解企图纵深的突击，而是把敌人看得太呆笨。我觉得东华山、五都寨取得后而不能乘胜完成和展开其胜利，主力不应再留恋，这种留恋在任何方面会达到无谓的牺牲。\footnote{《彭德怀关于作战等问题给军委的信（1934年4月1日）》，《中央苏区第五次反“围剿”》（上），第154页。}
\end{quote}


此战成功，蒋介石大为开怀，开出赏单，参战各师每团发赏金1000元。\footnote{《蒋中正电熊式辉贺国光参加南丰战役各师准每团给赏洋一千元（1934年3月15日）》，蒋中正文物档案02020002000025。}同时，他也提醒陈诚：“推测匪情，其在枫林、三溪现在抵抗一线为其之生命线，亦为我军进入赣南匪区胜利之第一步，故匪必在现地死守顽抗也。我中路军主力不必求急进，只要固守现地，作成持久之局，以求薛、汤两路之发展，则匪经此战必崩溃更速，不必心急也。”\footnote{《枫林三溪一线为我军进入赣南之第一步并示作战要领（1934年3月13日）》，《陈诚先生书信集·与蒋中正先生往来函电》（上），第131页。}


红军后撤后，国民党军乘胜继续进击，逐渐进入苏区基本区域，逼近广昌。早在1933年底，红军指挥部就指示在广昌地区构筑防御工事。因此，之前“从未发现红军占领整然一线的防御阵地，构筑工事，采取有形的防御战斗”的国民党军将领，此时“沿白舍东西之线，突然发现红军构筑有整然的防御阵地，重要山头筑有两层射击设备的坚固碉堡”。不过，红军在这一带仍没有采取死守方针：“当面守备的红军不是主力部队，是新近扩军所编成的队伍；装备很差，携带的多是破旧步枪，堪用的不多，而且弹药极少；战士素质不佳，老弱参半”。\footnote{杨伯涛：《蒋军对中央苏区第五次围攻纪要》，《文史资料选辑》第45辑，第190、191页。}国民党军没有经过大的战斗，很快于3月15日占领白舍，随后向甘坊地区展开进攻。16日，双方在甘坊地区战斗稍烈，国民党军得到“空军轰炸”\footnote{《国民党北路军顾祝同部与中央苏区红军作战情形报告书》，《中华民国史档案资料汇编》第5辑第1编军事（4），第48页。}支持，付出一定代价后于17日控制甘坊。至此，国民党军打开了通往广昌的大门。


1934年4月1日，蒋介石在南昌召开军事会议，推出北、东两路军行动计划，决定北路军第三路军陈诚部4月进取广昌，5月推进宁都；第六路军薛岳部4月底进占龙岗，5月底攻取古龙岗；其余各部在陈、薛部后跟进。东路军第十纵队汤恩伯部和第八纵队周浑元部5月进取闽北建宁、宁化得手后推进石城；第四纵队李延年部和第九纵队刘和鼎部进取闽西连城、永安，续进长汀。该计划以广昌、建宁作为攻取重点。


根据全盘作战计划的要求，1934年4月初，国民党军第三路军开始向广昌推进。该路军制定的进攻计划是：“沿盱河两岸，逐步筑碉，向甘竹、广昌进展，完成南广公路，并诱匪主力决战而歼灭之。如情况许可，则一举进占广昌”\footnote{《赣粤闽湘鄂北路剿匪军第三路军五次进剿战史》（上），第八章，第3页。“盱河”即“盱江”。}具体部署是：右翼之第九十六师、四十三师、九十七师向红三军团第四师扼守的甘坊石下寨、池埠阵地进攻；左翼第六十七师附山炮—连，向扼守瑶陂的红五军团第十三师进攻；第九十四师附山炮一连，攻击扼守白舍附近的红九军团。其攻击重点放在盱河西岸。循着运动战的思路，红军在广昌外围没有采取节节防御的作战方针。根据第五次反“围剿”以来的基本用兵思路，红军主力部队一、三军团开至广昌附近后，被置于机动位置，前线防御主要由新编成的红九军团及地方独立部队担任，广昌外围第一道防御线甘竹以北地区只布置了象征性的防御。\footnote{对一、三军团的使用，早在第五次反“围剿”的初始阶段，当时的中革军委代主席项英就指示了一个基本原则：“一、三军团……要隐蔽控制，以便突击。”（项英：《对敌情的分析和作战意见》，《项英军事文选》，第176页）此后，一直到广昌战役，一、三军团战略上始终处于机动状态。}4月初战役开始，国民党军进展顺利，几乎没有经历大的战斗，很快占领甘竹以北地区。即使如此，国民党军也并不急于伸展，而是按部就班步步推进。10日，国民党军首先进占罗坊。13日，双方在甘竹外围展开争夺战，国民党军占领罗家堡、李家堡等地，随后控制甘竹。参加战斗的红一军团二师四团团长耿飚回忆：“当时我们在甘竹‘守备’。敌人前进半里多一点，便开始修乌龟壳。”红军在与其对阵中，“敌人有碉堡依托，火力又猛，我们的掩护部队由于弹药匮乏，根本无法对射。等我们冲锋部队冲到双方中间地带时，敌人的大炮便实施集火射击。由于敌人实现早已设计好战斗层次，炮火很准。我们一次又一次地被炮火压回来，除了增加一批又一批伤亡之外，一无所获”。\footnote{《耿飚回忆录》（上），第148页。}


面对国民党军的进攻，在广昌应否防御问题上，红军高层看法基本一致。李德回忆：“党的领导人把这个本来不很重要的县城，视为必须保住的战略要地，因为他卡住了通向苏区心脏地带的道路。此外他们认为，将广昌不战而弃，政治上无法承担责任。”\footnote{〔德〕奥托·布劳恩：《中国纪事1932～1939》，第92页。}这一说法可从周恩来当时发表的文章中得到证实：“每个同志都要认识，敌人这次占领广昌的企图，与以前四次战役更有着不同意义的形势。敌人在持久战略与堡垒主义的战术下，进占广昌是其战略上重要的步骤，是深入中区，实行总进攻的主要关键。我们要为保卫广昌而战！战斗胜利了，将造成敌人更大的困难与惨败的条件，将造成我们彻底粉碎五次‘围剿’的更有力的基础。”\footnote{周恩来：《工农红军和全苏区群众一致动员起来，为保卫广昌而战！！！》，《红星》第33期，1934年4月22日。}随着苏区的巩固和发展，中共在战略抉择上受政治、经济背景制约，选择空间反而受到影响。初期反“围剿”作战中，红军规模相对较小，资源供给也较多倚赖打土豪的收入，大规模后退和前进的运动作战游刃有余。而随着红军的扩大，苏区周围土豪被打尽，对资源吸取的正规化（如税收制度的建立），红军再要流动作战已不像初期那样较少顾忌。湘鄂赣就反映：“如果红军在有时候未打的（得）胜仗，说红军是吃饭的，所有慰劳品，都不送去了……前次敌人进攻万载，红军没有与之抵抗，就说红军是吃饭的，更是说把红军吃，很（肯）把狗吃。”\footnote{《××同志关于湘鄂赣青年团工作给中央的报告（1932年9月3日）》，《湘鄂赣革命根据地文献资料》第2辑，第441页。}共产国际远东局也提到如果实行诱敌深入，“当地老百姓就会对我们失望，我们就会丧失补充红军队伍的可靠来源”。\footnote{《共产国际执行委员会远东局给共产国际执行委员会政治书记处政治委员会的电报（1933年4月3日）》，《共产国际、联共（布）与中国革命档案资料丛书》第13册，第374页。}这样的群众反应事实上代表了红军在新形势下所应承担的义务，这是中共制定全盘战略时不得不考虑的一个因素。


与此同时，红军内部也不无压力，国民党方面收集的有关资料谈道：

\begin{quote}

现匪军之所谓战斗员，苏区农民，几占十分之七八，彼等皆被伪政府所欺骗利诱，即每人或分有田地，或惑于所谓“红军眷属优待条例”，故在匪军中较为坚决可靠，唯其眷属及所分得之田地，均在苏区，若为国军所占，则向之藉以维系彼等者，自失效用。我军占领广昌后，该县籍之匪兵，日久势将渐渐离异逃逸，影响匪军本身之战斗力，故有不得不死守广昌之苦衷。\footnote{《赣粤闽湘鄂北路剿匪军第三路军五次进剿战史》（上），第八章，第57页。}
\end{quote}


另外，国民党方面还注意到，红军“时时提出夺取中心城市之口号，但是日久终未见诸事实。加以受我军物质封锁，生活一天比一天困苦，自五次围剿后，每次战役，均打败仗，伤亡残废者，触目皆是，因以人人自危，个个灰心。即所谓坚决分子之党员团亦常提出疑难质问，并有逃逸事情发生。上级指挥员，每感难于应付”。\footnote{《国军五次围剿赣匪崩溃近况》，《军政旬刊》第19、20期合刊，1934年4月30日。}


为展开广昌保卫战，苏区中央决定调集红军主力9个师，尽力抵挡国民党军的进犯。在前方成立野战司令部，朱德为司令员，博古为政治委员，直接指挥前方战事。根据中革军委当时的解释，其主要作战方针是：“1.集中红军主力打击和消灭敌之主要进攻。2.以必要的兵力尽力钳制其它方面。3.派遣得力的地方独立部队，挺出敌人近的与远的后方，发展游击战争，创造新苏区，以钳制和调动敌。”\footnote{《朱德周恩来王稼祥致闽浙赣电（1934年4月2日）》，《闽浙皖赣革命根据地》（上），第702页。}可以看出，中共当时虽然迫不得已在广昌实行防御，但还是希望不要把战争打成消极防御。中共中央局机关刊物《斗争》发表社论强调：“在敌人的堡垒政策面前，发展游击战争，可以使敌人力量很大的分散与削弱，使主力红军的战斗得到更便利的条件。”“建筑支撑点，制造和使用地雷、弩箭等防御武器来打击敌人（这方面赣东北有很好的模范）。但必须反对把中心力量完全放在这个工作的防御路线，并反对乱筑防御工事。”\footnote{《五次战役的第二步的决战关头和我们的任务》，《斗争》第58期，1934年5月5日。}张闻天也指出：“分兵把口，与堡垒主义，是紧密联系着的，这是单纯防御的机会主义倾向的又一种具体表现。这种倾向，实际上不但不能保卫苏区，而且正便利于敌人的各个击破。在这里，我们应该清楚指出，积极的发展游击战争，把我们的基本游击队深入到敌人远后方与侧翼去活动，是我们保卫苏区的最好办法。”\footnote{张闻天：《我们无论如何要胜利》，《红色中华》第183期，1934年5月17日。}只是，在总体战略属于被动防御的背景下，局部的小范围的运动战，面对国民党军优势兵力的稳步推进，效果实属有限。


其实，在具体的战争指导中，运动战思路也难以真正得到贯彻。根据对国民党军主攻方向的判断，14日，林彪、聂荣臻向朱德提出建议：“我军主力目前宜隐蔽于千善、石嘴以南诸地，而以一部伪装主力在现地诱敌，主力准备突击经河西前进之敌，和准备突击向大田市、溪口前进之敌。如周（浑元）纵队联合向南采取跃进时，我们更便于突击他。”\footnote{林彪、聂荣臻：《关于我对敌罗、樊、周纵队的部署建议（1934年4月14日）》。}16日，林、聂再次提出：“即令在敌人采取编成两个纵队同时架河而上的行动，我一三军团亦不应分开。”这一建议实际是主张将主力收缩，待国民党军充分展开后，再待机出击歼敌。所以他们强调：“如三军团在现地不动，不仅不便于对付敌人自由河西前进的情况，对令地方队，对敌经河东活动南进时亦成了在正面和距敌的短距离内阻敌。”\footnote{《林聂对陈敌将向广昌前进的估计及突击该敌的部署报告（1934年4月16日）》，《中央苏区第五次反“围剿”》（上），第158页。}而彭德怀、杨尚昆则判断国民党军将由盱江东岸南进，主张：“我主力应在芙蓉塅、大罗山地带与敌决战，以充实的一营固守延福庵，扼守制敌两个纵队不易联系……以第六师自延福庵至大罗山钳制樊纵队，以四师三师一军团十三师为突击兵团在芙蓉塅、里峰地域决战。”\footnote{《彭杨关于在芙蓉塅、大罗山地带与敌决战的部署报告（1934年4月16日）》，《中央苏区第五次反“围剿”》（上），第157页。}这种运用主力在芙蓉塅、大罗山山地地带与敌决战的设想，和中革军委一段时间来实行的防御思路具有相当的一致性。但是，虽然山地作战有利于防御一方尤其像红军这样火力较差的防御者，但这一战地事实也在国民党军预料之中，难以达到出奇制胜的效果。所以，林、聂等的设想可能包含着更多的制人而不制于人的争取主动思路。不过，彭德怀和中革军委也许可以辩解，根据五次“围剿”以来国民党军的一贯方针，他们“取逐步构成野战工事节节推进以求得火力掩护的可能极多……未完成前两翼暴露南进广昌的可能减少”。\footnote{《彭杨关于在芙蓉塅、大罗山地带与敌决战的部署报告（1934年4月16日）》，《中央苏区第五次反“围剿”》（上），第157页。}在此情况下，林、聂的设想固然不错，但不能排除落空的可能，而放弃在大罗山山地地带作战，广昌也将陷入无险可守的境地。所以，大罗山的阻击战仍然成为当时背景下的逻辑选择。


4月19日，国民党军在判断红军主力集中于盱江东岸地区后，转调部分兵力用于东岸，开始向该地区的延福嶂、大罗山一带红军主力发动进攻。上午10时半国民党军第六师开始猛攻大罗山，“未刻占领大罗山五二五七高地”。\footnote{《樊崧甫电蒋中正等（1934年4月19日）》，蒋中正文物档案002090300072421。}由于一军团未按预定时间赶到战场，红军对延福嶂、大罗山一带阵地并未取固守态势，据周恩来报告：“三军团主力七时半到马鞍寨、磜上，他们未依军委突击攻大罗山之敌，而拟待敌攻天井围、墓坑时再突击。”\footnote{《周关于敌占大罗山等地后我军突击敌部署给朱、博、李电（1934年4月19日）》，《中央苏区第五次反“围剿”》（上），第161～162页。}而据三军团的命令：“我军以于邓家庄、石源、浮竹、大罗山地带突击该敌于我防御地带之前而歼灭之为目的。”\footnote{《三军团歼敌命令（1934年4月19日）》，《中央苏区第五次反“围剿”》（上），第160页。}为此，红军集中6个师兵力分左、中、右三队向大罗山一线进发，左路为第十三、第六师，中路为第四、第五师，右路为第一、第二师，准备取三路包围之势，突击并消灭深入红军防线的敌人。


国民党军占领大罗山后，一度曾继续向纵深追击，第十八旅旅长向该师师长报告：“当面之匪击溃后，向大罗山东南溃窜，我已派队追击中。”但是，第六师师长周嵒在下午3时半下达的命令中根据国民党军第五次“围剿”以来一贯的稳扎稳打方针强调，该部应“迅即构筑守备公事限本夜完成”。同时，该部第三十六团九连攻至红军重兵集结的天井围附近时，请求炮兵予以火力支持，周嵒当即指示：“一、天井围过于突出，该团第九连应在炮火掩护下，即行撤回。二、大罗山至平山间工事，须迅速构筑。”\footnote{《赣粤闽湘鄂北路剿匪军第三路军五次进剿战史》（上），第八章，第23页。}国民党军的如上处理，使其能避免深入红军阵营。


当晚，红军主力完成集结，19时左右，向大罗山一线国民党军第六师、第七十九师发动猛烈反攻，准备歼灭突进之敌。国民党军凭险顽抗，据守大罗山的第六师十八旅三十六团团长李芳在率部向红军反击时被炮火炸死，由该团第三营营长接替指挥。是役，红军志在必得，集中了几乎所有能打硬仗的部队，反攻“炮火极为猛烈，双方死伤亦极奇重……战斗时间竟达十五小时之久，可谓作战以来仅有之剧烈斗争”。\footnote{《国民党北路军顾祝同部与中央苏区红军作战情形报告书》，《中华民国史档案资料汇编》第5辑第1编军事（4），第80页。}战至20日凌晨3时许，红军虽然集中了最精锐的主力部队，仍未实现歼灭敌军的目的，被迫撤出战斗。


大罗山反攻失利后，红军退至饶家堡一带，准备利用深山密林继续对来犯敌军实施打击，力争歼灭其突出部队，命令“三军团由墓坑及其以南山地和天井围向樊敌主力及肖师行猛攻干脆的突击”，“准备集结主力，下最大决心与敌六、七个师作较大的决战”。\footnote{《军委关于阻敌向甘竹、长生桥前进的作战指示（1934年4月20日）》，《中央苏区第五次反“围剿”》（上），第163页；《陈伯钧日记（1933～1937）》，第207页。}20日下午，红军向深入饶家堡地区的七十九师二三五旅部队发动突击，“三面围攻，其势汹汹，大有‘请君入瓮’之概”。\footnote{《赣粤闽湘鄂北路剿匪军第三路军五次进剿战史》（上），第八章，第28页。}同时，红军一部绕向二三五旅后方，准备截断其后路，但在前排遭遇国民党军七十九师主力，无功而返。在无法截断敌后路的情况下，红军加紧对突进部队二三五旅的打击，“这天晚上，阴雨绵绵，不便射击，红军与敌人进行白刃格斗，战斗异常激烈”；\footnote{《王平回忆录》，解放军出版社，1992，第55页。}“饶家堡西北高地及大坪咀山阵地，失而复得者，凡五六次”。\footnote{《赣粤闽湘鄂北路剿匪军第三路军五次进剿战史》（上），第八章，第29页。}中革军委对这场战役高度重视，与前方电文往来不绝，朱德不断就前方状况发出电报，并强调：“这不是命令而是给你们下决心的建议。”\footnote{《朱对我军攻击饶家堡敌之建议（1934年4月21日）》，《中央苏区第五次反“围剿”》（上），第166页。}21日凌晨4时，他还乐观指示：“饶家堡战斗得手后，应集中一、三军团炮兵与迫击炮作有组织之炮击。”\footnote{《朱关于集中炮兵、迫击炮和一、三军团行动的指示（1934年4月21日）》，《中央苏区第五次反“围剿”》（上），第166页。}但是，形势发展并不以人们意志为转移，由于红军久攻不下，到21日拂晓，对方援军第九十七师源源到来，红军无法实现歼灭敌军的目标，被迫退出战斗。随后，国民党军又进占云际寨、香炉峰、高洲瑕一线。红军主力退往广昌地区，盱江东岸战事告一段落。红军在大罗山、饶家堡的两次战斗，是前方指挥员在当时总体以被动防御为主的总战略下，尽力发挥红军运动战特长的两次尝试，但由于红军久战疲劳，国民党军兵力又过于厚集，红军歼灭敌人的目标难以实现。另据国民党方面战史载：

\begin{quote}

伪参谋长林义光供：“……匪在日间，畏我飞机之轰炸，枪火之猛烈，为避免损害计，采取夜战。其攻击部署，以少数兵力，用于正面佯攻，以重兵力用于两翼，如冲锋两次不成，即行撤退，并在日间，预行选定进攻路线地区及目标，但匪兵缺乏训练，且多新兵，常畏缩不前。”\footnote{《赣粤闽湘鄂北路剿匪军第三路军五次进剿战史》（上），第八章，第59页。}
\end{quote}


红军战记则报告，在此前不久的东华山战役中，“七团有些新战士不会打手榴弹，敌人冲来时，把手榴弹交给班长打”。\footnote{刘鹤孔：《东华山战斗》，《火线上的一年（1933～1934.8）》，红军总政治部红星社编印，1934，第93页。}来不及得到必要训练的新战士的大量增加对红军战斗力有着重大影响，这也是战役难以获胜的不可忽视的原因。大罗山、饶家堡两战役，尤其是大罗山战役，红军出动了几乎所有主力部队，面对国民党军不完整的两个师，仍然不能取得充分的战果，这和第五次反“围剿”以来，国共双方战斗力的消长无法分开，它预示着红军此后的战斗将更加艰难。


相对于盱江东岸，盱江西岸红军力量更为薄弱，仅有两个师番号实际不到一个师的部队，部队作战能力也相对较差。国民党军在此出动了第十一师和第九十八师两师部队，攻势发动后，很快占领长生桥、伞盖尖、火神岩等地。广昌已处于国民党军直接威胁之下。

\subsection{广昌保卫战}


外围战事连遭失利后，4月21日，林彪、聂荣臻以“万万火急”致电朱德、周恩来，提出：“如突击当前之敌无把握且广昌××（原文如此——引者注），三军团本晚须即由沙子岭以南渡河，与敌决战于广昌附近”，\footnote{《林聂关于与敌决战于广昌附近的建议（1934年4月21日）》，《中央苏区第五次反“围剿”》（上），第167页。}主张放弃在外围继续抵抗，直接在广昌城附近与敌决战。同日，中共中央负责人博古、中革军委主席朱德、代总政治部主任顾作霖发布命令，号召继续展开广昌保卫战，要求红军“应毫不动摇的在敌人炮火与空中轰炸之下支持着，以便用有纪律之火力射击及勇猛的反突击，消灭敌人的有生力量”。\footnote{《中央、军委、总政保卫广昌之政治命令（1934年4月21日）》，《中央苏区第五次反“围剿”》（上），第167页。}同时，面对前线不利形势，中共中央指出，由于国民党军战略的变更，“使我们红军消灭敌人的战斗，须在一些新的条件下来进行”，强调：

\begin{quote}

动员群众武装起来，参加革命战争，发展广大的游击战争，是战地党和苏维埃的第一等重要的任务。

要以更多的地方部队，发展广大的游击战争，在敌人左右前后，在敌人的封锁线外，在敌人的堡垒间隔之中，在敌人的远近后方，到处去寻找敌人作战，冲破封锁，钳制敌人，分散敌人，疲惫敌人，隔断敌人，瓦解敌人，这样来配合和掩护我主力红军，得以运用自如，实施突击，而最终的消灭敌人。\footnote{《中共中央委员会、中央人民委员会给战地党和苏维埃的指示信（1934年4月24日）》，《斗争》第58期，1934年5月5日。}
\end{quote}


虽然中共中央要求开展广泛的游击战争来改变当前的被动局面，但远水难解近渴，游击战事实上已难以担起改变战场形势的重任。相反随着广昌外围防线被步步压缩，红军活动空间愈来愈小，广昌战役已越来越向阵地遭遇战方向发展。22日，周恩来致电朱德、博古、李德，提出三项建议：“1.最紧急时须调二十三师主力加强广昌守备。2.一、三军团要能在一起突击敌。3.……拟令董朱二十四日西移二十五日可参加广昌战斗。”\footnote{《周为保持广昌战斗胜利的建议（1934年4月22日）》，《中央苏区第五次反“围剿”》（上），第169页。}这实际是要求把红军最精锐的一、三、五（董、朱部）军团全部投入保卫战，反映出中共高层对广昌防御的极端重视。同日，博、朱、李复电周恩来，未采纳将红五军团西调的建议，而仍指望通过以红一军团在盱江西岸诱敌，再由“三九军团包括十三师在内突击该敌”。\footnote{《博朱李关于广昌战斗部署情况致周电（1934年4月22日）》，《中央苏区第五次反“围剿”》（上），第170页。}不过，随着国民党军迅速向广昌逼近，这一计划也迅成泡影。26日，中革军委下令组成3个作战集团：东方集团，由红九军团及红十三师组成，红九军团军团长罗炳辉、政委蔡树藩负责指挥，任务是在盱江东岸钳制敌人；西方集团，由红一、三军团及红二十二师组成，由朱德直接指挥，任务是在盱江西岸广昌以西及西北地域消灭进犯之敌；守卫广昌部队，由红十四师等部组成，任务是坚守广昌工事。这一部署意味着中革军委已不顾双方实力对比，准备在广昌城外围进行大规模的兵团作战。不过，这时中革军委对战役前途其实已不乐观，《火线》发表社论强调：“保卫广昌战斗虽是五次战役中的一个重要战斗，但不能认为是五次战役唯一的决定最后胜负的一个战斗。五次战役决定最后胜负的战斗，主要的在于我们能否消灭敌人的有生力量。假如我们能消灭敌人的有生力量，我们不仅能恢复某些被敌人一时侵占的苏区，而且可以扩大更广大的新的苏区。”\footnote{《论保卫赤色广昌》，《火线》第130期，1934年4月25日。}显然已在军中为最后放弃广昌做舆论准备。


4月27日，国民党军经过短暂休整、准备并构筑碉堡、封锁线后，出动6个师兵力分左、右两路沿盱江两岸开始向广昌发动进攻，“河西三个纵队并进，河东一个纵队前进”。\footnote{《博朱关于广昌战况及决令二十八日晨放弃广昌致电周（1934年4月29日）》，《中央苏区第五次反“围剿”》（下），第82页。}广昌附近山地较少，多绵延起伏的丘陵，地势相对平坦，“虽田畴相望，小溪横流，然实以适合大军团运动天然地带”。\footnote{《陆军第十四师广昌长阴之役战斗详报》，《汗血月刊》第1卷第5期，1934年7月20日。}虽然红军事先做了一定准备，在广昌外围构筑工事，期望进行顽强防御，\footnote{时任红十四师师长、广昌警备司令的张宗逊回忆：“红十四师用了五个月时间，在广昌城和其以北到甘竹一线，共构筑了坚固的大小支撑点和碉堡十多个，在广昌城周围修了七个支撑点……每个支撑点驻守的兵力不等，有的为一个营，有的一个连，有的一个排。红十四师以四个营的兵力守支撑点，以五个营的兵力作为反冲锋部队。”（《张宗逊回忆录》，解放军出版社，1990，第120页）作为守备部队，红十四师在防御工事上做的这些工作不能说是偏于死守的，而其兵力使用也基本符合当时中革军委贯彻的“短促突击”战术原则，只是相对比一般部队要更重守御一些。}但红军工事在国民党军重武器攻击下，往往无法发挥作用，对国民党军的突击由于对方兵力厚集也难有效果。国民党军战史记载：“是日匪以一部守平面岭卖竹坪大仙山坚固匪碉，以伪三军团全部，及伪九军团一部，由卖竹坪附近，向我第十四第六十七两师正面猛烈反攻……午后，伪三军团犯我十四师正面，伪二师犯我六十七师左翼，战斗更为激烈。”\footnote{《赣粤闽湘鄂北路剿匪军第三路军五次进剿战史》（上），第八章，第42页。}经过一整天激战，红军虽表现英勇，向国民党军“迭次冲犯”，并“以密集部队往复冲锋，毫不混乱”，\footnote{《陈诚罗卓英电蒋中正顾祝同等（1934年4月27日）》，蒋中正文物档案002090300073052；《陆军第十四师广昌长阴之役战斗详报》，《汗血月刊》第1卷第5期，1934年7月20日。}予国民党军以重大杀伤，但自己也付出相当惨重的代价，对等的消耗，兵力上处于绝对劣势的红军显然难以承受。当晚，在广昌前线直接指挥作战的博古、朱德、李德联名致电留守瑞金的周恩来，提出：“广昌西北之战未能获得胜利，现只有直接在广昌支点地区作战之可能，但这不是有利的，提议放弃广昌而将我们的力量暂时撤至广昌之南。战斗经过另报，请立即以万万火急复。”\footnote{金冲及主编《朱德传》，第322页。}周恩来随即复电，表示在红军主力受到较大损失而在广昌直接作战又无把握的情况下，“原则上同意放弃广昌，但仍须以一部扼守广昌，迟敌诱敌，抽一军团秘密东移，突击汤（恩伯）敌”，并强调“最后决心由你们下”。\footnote{金冲及主编《周恩来传1898～1949》，人民出版社、中央文献出版社，1989，第274页。}在周恩来复电未予反对后，28日，朱德下达放弃广昌的命令。当天上午，盱江东岸国民党军第六十七师、七十九师进占广昌，红军全线后撤。红一军团东移珠市坪、尖峰地区，九军团后撤至扬家坪、新安地区，三军团南撤头陂等地，十三师在北华山、马坊寨、里丰一线留守原阵地抗敌一天，掩护主力撤退。广昌保卫战进行了18天，国民党军伤亡2600多人，红军伤亡5000余人，占参战总兵力的1/5，其中红三军团伤亡2700余人，达到军团总人数的近1/4。\footnote{国民党方面当时报告：“是役匪方前后伤亡在四千以上，获枪二千余枝。”见《南昌行营第二厅致汪精卫等电（1934年4月29日）》，《中央革命根据地革命与反革命斗争史料》（下），中国社会科学院近代史研究所藏。}红九军团的第十四师因伤亡过大，已难成建制，余部被并入其他部队。


广昌战役，红军打破不固守坚城的惯例，展开大部队参加的保卫战，被认为是苏区中央执行的被动防御军事政策的体现。彭德怀在回忆录中提到，他对李德实施的拼消耗战术当时就提出尖锐批评，彭还强调，广昌保卫战中，国民党军“每次六、七架飞机轮番轰炸。从上午八、九时开始至下午四时许，所谓永久工事被轰平了。激战一天，我军突击几次均未成功，伤亡近千人”。\footnote{《彭德怀自述》，第190页。}也许就是针对彭德怀当时的指责，博古、朱德在放弃广昌后给周恩来的报告中，特地强调在国民党军步兵向广昌攻击之先，“并未有炮兵与空中轰炸”，同时其对27日当天战役的发展状况描述是：

\begin{quote}

敌先攻我翼侧，即占我左翼第一线阵地……因地形系狭山，我全部力量正面不过十里，故四个师突击师以充分够用。敌主要纵队立即密集队形在西岸谷地前进，我们决今让他近一些，而以三军团突击敌人之后部队，一九突击敌人之先部队，而实际上三军团过早进入战斗，且系突击敌之先头部队，因此敌停止前进构筑工事与准备反突击，而我一军团则不能全部展开一师之利用，只能转移至右翼四、五师之间，且只在战斗最后阶段才进入战斗，形成敌我对峙，未获结果，决定脱离战斗。\footnote{《博朱关于广昌战况及决令二十八日晨放弃广昌致电周（1934年4月29日）》，《中央苏区第五次反“围剿”》（下），第82～83页。}
\end{quote}


博古、朱德的这一报告，颇值玩味，它道出了争论另一方的看法，而这在以往通常是被忽略的。当然，这并不意味着该报告陈述的都是事实。首先，报告所说国民党军进攻之前并未使用炮兵与空中轰炸，就事实本身言，应为可信。国共双方战史都提到，27日的广昌战斗在拂晓前就开始进行，战斗开始后，双方战线很快形成犬牙交错状态，空军轰炸事实上难以措手。包括彭德怀在内许多回忆录中提到的空中轰炸，应是将第五次反“围剿”中的一般状况作了误植。值得一提的是，自4月10日广昌外围战斗开始到28日广昌失陷，正是南方春雨连绵季节，19天内，阴雨天就占了15天。\footnote{参见《陈伯钧日记（1933～1937）》，第202～214页。上述统计还不包括27日当天傍晚的大雨。}以当时的技术能力，这样的天气飞机实际难以发挥作用，所以无论就整个广昌战役还是当天的广昌战斗言，国民党军的空中优势由于气候制约并未得到充分发挥。另外，当时技术条件下，炮兵作用发挥尚较有限，李德观察到：“敌人少有系统的预先进行空军轰炸及炮兵的火力准备。其进攻一开始，就出动步兵（步兵的火器及其突击队），炮兵及迫击炮的射击，主要的是为着直接援助步兵；而飞机则不断的进行战场上的观察，以妨碍我们在战场上的机动，飞机轰炸也和炮兵一样主要的是为着直接援助其步兵。”\footnote{华夫：《再论战术原则》，《革命与战争》第4期，1934年5月18日。}这一观察从国民党军第六路军所编战史中侧面得到证实：“各师所有之迫击炮，多属旧式，瞄准机构不甚健全。对广大目标之射击，勉可使用，若以之摧毁碉堡，则耗弹多而命中公算少。”\footnote{《剿匪纪实·赣南围剿》，第六路军总指挥部，1937，第9页。}不过，博、朱的报告当然不仅仅止于陈述上述事实，其潜台词应为通过对国民党军火力的压低，为广昌战斗的决策辩解。应该指出，虽然在27日进攻之前，国民党军没有采用炮火和空中集中轰炸，但其火力优势在战斗中仍然体现得至为明显，国民党军战史多次提到，其在激烈的攻防战中获胜的主要原因为“火力旺盛”，红军进攻“被我炮兵火力压倒，始未得逞”。\footnote{《赣粤闽湘鄂北路剿匪军第三路军五次进剿战史》（上），第八章，第42～44页。}试图以国民党军没有大规模使用炮火和空军来掩盖双方火力上的巨大差距，不会有充足的说服力。


其次，博、朱报告将战斗失败相当程度上归咎于三军团的过早突击，这很可能也就是当日彭德怀和李德发生激烈冲突的直接诱因。但是，在宽十里的战场上以4个师兵力欲对国民党军5个师实行突击，在红军已经屡遭损失，本身师建制就难以和国民党军相比，火力又远逊对手的情况下，这一决策本身就不现实。何况，战场形势瞬息万变，由于国民党军进展甚快，红军的反应往往是被动的，兵力又捉襟见肘，将战斗失败归因于前线指挥，而不检讨自身在战争指导上的失误，确实有失公见。


广昌战斗中的这种战场上的拼消耗战法是中共中央对敌我力量对比还缺乏深切了解的一种反映，广昌战役后，红军转而采取依托坚固工事，实行固守的战法，虽然这种战法在当时和日后均遭到批评，但起码可以比在广昌坚持更长时间。


在国民党军第三路军进攻广昌同时，为配合进攻广昌，4月中下旬，位于藤田、沙溪一带的国民党军第六路军向龙岗一线展开进攻，以“与正向广昌进展之第三路军齐头并进”，\footnote{《剿匪纪实·赣南围剿》，第17页。}牵制红军主力集中。21日、22日，其先头部队第九十师、第九十二师与红军在韶源、上固交火，红军在此方向兵力薄弱，尽力阻击后即撤出战场。薛岳报告，国民党军在此遭遇了红军的地雷：“地雷系用烟罐竹筒内实炸，故其爆炸力甚微。”\footnote{《薛岳电蒋中正梁师推进至上固构筑近碉楼在上固南方路上发现匪之地雷又韩师搜剿队在含下路上亦发现匪之地雷等（1934年4月22日）》，蒋中正文物档案002090300072360。}显然，由于技术的原因，红军此时的地雷尚未发挥出太大威力。30日，国民党军第九十九师向龙岗进兵，红二十三师及独立第一、二团4000余人与对手激战竟日，5月1日下午，龙岗北端红军碉堡在与国民党军“肉搏十余次”\footnote{《剿匪纪实·赣南围剿》，第19页。}后终告失陷，龙岗随之为国民党军占领。陈毅当时曾谈到这一战斗：“龙岗之支点，在构筑上不十分坚固，被敌人包围一天一夜，两翼的突击部队侧击又未得手，而其上级首长尚令该支点内守备队死守，请示上级，致失时机，而使守备队受到损失。”“在突击队失效，估计敌兵力强大不能固守，即应机断的给敌人以短促的突击，随即撤退，最好先保持退路不陷于敌人包围中。”\footnote{陈毅：《几个支点守备队的教训》，《革命与战争》第4期，1934年5月18日。}


4月中旬，南路国民党军陈济棠部在蒋介石一再催促下，向苏区南部粤赣省发动进攻。3月中旬，粤军第一军第一师占安远，第二师占信丰，第二军第四师进驻南康、赣州。随后，又以第三军李扬敬部为主编成第二纵队，辖4个部队，主攻南线。其作战计划为：“纵队先进攻筠门岭，俟巩固占领后，候机向会昌城进攻。以主力从寻乌、吉潭经澄江、盘古隘，向筠门岭正面攻击；另以一部从武平的岩前经武平所向筠门岭东南侧攻击，并掩护纵队主力右翼安全。”\footnote{曾其清：《陈济棠进攻筠门岭红军的经过》，《文史资料选辑》第45辑，第153页。}当时，红军在粤赣只有二十二师一个师的主力部队，在汶口、盘古隘、筠门岭设置三道防线。4月初，南路军由安远、寻乌和平远、武平两个方向推进筠门岭、会昌。陈济棠部发动进攻后，红军寡不敌众，汶口、盘古隘相继失守，退守筠门岭。筠门岭距粤赣省府所在地会昌仅55公里，位于赣粤闽三省交界处，是江西通往广东、福建的交通要塞。21日拂晓，陈部集中重兵猛攻筠门岭，战至当日午后，红军被迫放弃筠门岭，退守门岭县站塘地区。由于陈济棠并不真心想对苏区形成威胁，攻占筠门岭后即电蒋报捷，\footnote{《陈济棠电蒋中正顾祝同等（1934年4月22日）》，蒋中正文物档案002090300074460。}算是对蒋介石和舆论有所交代，随即与红军联络，止步不前，南路战事仍保持相对稳定。


4月下旬，东路国民党军第八、第十纵队对建宁发动进攻。红军以第五军团为主力，红一、七、九军团参加战斗，在建宁外围将军店、驻马寨等主要阵地节节抵抗，其中，红一军团是在广昌失利后，马不停蹄赶往建宁前线。5月7日，国民党军第八纵队向将军店发起进攻，同时其第三、第十纵队等部准备从侧后向建宁挺进。9日，红军未能阻挡住敌军的攻势，放弃将军店，退守离建宁城十多里的驻马寨一线。12日，国民党军向驻马寨发动进攻，红军且战且退，虽然予国民党军以重大杀伤，但无法阻挡其前进步伐。16日，红军退出驻马寨，午后，国民党军第八十八师占领建宁。到5月中旬，国民党军已控制西起龙岗，中经广昌，东至建宁、泰宁、归化一线的广大地区，苏区区域被大大压缩。

\section{红军战略转移的准备}

\subsection{中共中央计议突围}


反“围剿”的不利形势，提示在苏区内打破“围剿”的可能性日趋渺茫，虽然博古1934年5月为《红色中华》撰写的社论中仍然宣称：“我们要保卫土地、自由、苏维埃，直至最后一个人，最后一滴血，最后一口气！”\footnote{博古：《我的位置在那边，在前线上，站在战线的最前面！》，《秦邦宪（博古）文集》，第257页。}但这更多只是政治宣传。政治离不开宣传，但宣传并不就是政治。就在博古发出上述豪言壮语同时，中共中央内部关于突围问题的讨论已提上日程。1934年2月，共产国际驻华代表埃韦特报告：“江西和福建的形势很困难，我们近期的前景不妙。不带多余的悲观主义应该承认，包围圈越来越小，敌人兵力在接近向我地区突破的一些地方。”\footnote{《埃韦特给皮亚特尼茨基的信（1934年2月13日）》，《共产国际、联共（布）与中国革命档案资料丛书》第14册，第81页。}面对此一局面，共产国际和中共中央都在苦寻出路。李德回忆，1934年3月，他曾提出以主力在中央苏区的“西南部或东南部突围”\footnote{〔德〕奥托·布劳恩：《中国纪事1932～1939》，第112页。李德1939年对这一问题的记述是：“自1934年夏提出以后的行动问题时，革命军事委员会在我的影响下立即提出两个方针：一个是竭尽全力保卫苏区；另一个是疏散。”见《布劳恩给共产国际执行委员会的书面报告（1939年9月22日）》，《共产国际、联共（布）与中国革命档案资料丛书》第15册，第348页。}的设想。5月，中共中央书记处在瑞金召开会议，提出在敌人逼近中央苏区腹地、内线作战不利的情况下，将主力撤离中央苏区。埃韦特在给共产国际转呈中共中央报告时说明了中共中央关于下一步行动的两个建议：“留在中央苏区，转入游击战，将其作为我们斗争的最重要方法。”“否则我们只有保卫中央苏区到最后，同时准备将我们的主力撤到另一个战场。”埃韦特个人虽明显倾向撤离中央苏区，但他同时强调：“只有在实行保卫的各种可能性都用尽之后并且在保存着我们大部分有生力量的情况下才应使用。”\footnote{《埃韦特给皮亚特尼茨基的报告（1934年6月2日）》，《共产国际、联共（布）与中国革命档案资料丛书》第14册，第128、129页。}


中共中央的撤离计划和共产国际一贯思路其实是契合的。早在1931年初，共产国际就指示：“必须进行顽强的斗争，把赣南基本的最主要的根据地保持在我们手里。但考虑到军队的主要核心力量在敌人压迫下有暂时被迫撤退的可能性，我们认为，现在就采取措施在湘西南和黔桂交界地区筹建辅助区是适宜的。同时务必更加重视在鄂湘川交界地区建立第二个主要根据地。”\footnote{《共产国际执行委员会政治书记处政治委员会给共产国际执行委员会远东局和中共中央的电报稿（1931年1月11日）》，《共产国际、联共（布）与中国革命档案资料丛书》第10册，第23页。}1933年3月，共产国际又提出：“在保卫苏区时，对于中央苏区来说特别重要的是保持红军的能动性，不要以巨大损失的代价把红军束缚在领土上。应该事先制定好可以退却的路线，做好准备，在人烟罕至的地方建立有粮食保证的基地，红军可以在那里隐蔽和等待更好的时机。”同时强调：“要建立和具备几个新的根据地，使政府军难于同我们对抗。我们积极评价第4军主力向四川转移。我们认为，在四川、陕南，以及尽可能在新疆方向开辟苏维埃根据地具有很大意义。”\footnote{《共产国际执行委员会政治书记处给中共中央的电报（1933年3月19～22日）》，《共产国际、联共（布）与中国革命档案资料丛书》第13册，第353、354页。}在中央苏区面临大规模进攻时，共产国际强调建立新根据地及对红四方面军转移的肯定，意味深长，以致埃韦特对此一度颇有疑虑，为消除其可能带来的负面影响，特电中央苏区解释共产国际的指示：“在我们看来，该电报是出于这样的考虑：万一敌人取得重大的胜利，我们必须保存和加强我军力量。而我方取得胜利时，我们应消灭敌人，必须时刻保持高度谨慎，使我军主力不受威胁。”\footnote{《埃韦特给共产国际执行委员会的报告（1933年4月8日）》，《共产国际、联共（布）与中国革命档案资料丛书》第13册，第393页。}虽然埃韦特试图淡化共产国际对转移的肯定，但其潜台词不可能不对中共中央未来的行动规划形成影响。正由于有此思想基础，当军事不利后，考虑转移就顺理成章。


得到中共中央有关报告后，1934年6月16日，共产国际复电指出：

\begin{quote}

动员新的补充人员的过程证明，中央苏区的资源还没有枯竭。红军作战部队的抵抗能力、后方的情绪等，还没有引起人们的担心。如果说主力部队可能需要暂时撤离中央苏区，为其做准备是适宜的，那么这样做也只是为了撤出有生力量，使之免遭打击……目的首先是要保存有生力量和为其发展创造新的条件，以便在有利的时机对日本和其它帝国主义军队和国民党军队展开广泛的进攻。\footnote{《共产国际执行委员会政治书记处政治委员会给埃韦特和中共中央的电报（1934年6月16日）》，《共产国际、联共（布）与中国革命档案资料丛书》第14册，第144～145页。}
\end{quote}


复电基本同意了中共进行战略转移的计划。次日，共产国际再电中共中央，明确谈道：“我们建议发动福建战役，将其作为预防和吸引敌人，进而便于保存苏区或从那里撤离（如果不可避免这样做的话）。”\footnote{《共产国际执行委员会政治书记处政治委员会给中共中央的电报（1934年6月17日）》，《共产国际、联共（布）与中国革命档案资料丛书》第14册，第146页。}


中央红军突围已成必然，各方对此了然于胸，1934年6月初，共产国际远东局委员赖安认为：“如果中国其他地区的军政形势以及国际因素不会导致发生‘出人预料的’重大冲突，以后几个月在阶级力量对比和政治重新组合方面也不会导致发生重要变化的话，那么在最近的将来，可能是秋天，中央苏区红军的主要力量将不得不放弃江西、寻找出路和在湘川方向寻找发展苏维埃运动的新的地区。”赖安在信中还透露了埃韦特等人的想法：“在目前的情况下，预先决定了中央红军主力‘迟早’将不得不放弃中央苏区。换句话说，这些同志，特别是从去年12月起，一方面阐发了关于本阶段在中央苏区不可避免地要发生决战的失败主义‘理论’。而另一方面与此相联系，仍在考虑红军能否坚持到春季或夏季，而现在是到秋季的问题。”\footnote{《赖安给哈迪的信》，《共产国际、联共（布）与中国革命档案资料丛书》第14册，第132、134页。}可见，撤离中央苏区的想法早在福建事变前后已成共产国际代表的议题。


红军即将展开的突围行动，其方向几乎不言自明。北面是国民政府核心区，向北无异于自寻绝路。东面濒临大海，也无出路，南面陈济棠虽不愿与红军作战，但保境欲望强烈，向南难免陷入粤方和宁方夹击之中，以此看来，向西几乎是唯一出路。1934年6月2日，埃韦特向共产国际报告中共突围计划时已经对行动方向有了明确交代：“虽然敌人从4军团1932年的远征时起学会了许多东西，虽然他们有比1932年更加强大的、我们必须克服的防线（赣江、赣湘边界和湖南湘江一线），虽然敌人在采取行动之初会拥有比1932年多得多的军队来组织追击，但对于我们的进一步推进来说反正都一样，恐怕我们也没有别的选择。”\footnote{《埃韦特给皮亚特尼茨基的报告（1934年6月2日）》，《共产国际、联共（布）与中国革命档案资料丛书》第14册，第129页。}赣江、赣湘边界和湖南湘江一线，这是明显的西进四川路线，而这和共产国际前一年指出的西进川陕路线完全一致。8月，回到莫斯科的埃韦特向共产国际汇报湖南省的情况时也侧面提到中央红军的行动方向：“最近一个时期，特别是在中央苏区，为了与这个省建立联系和在那里组建党组织作出了很大努力。作出这种努力是出于打通湖南的迫切需要，并考虑到保住江西已不可能了。”\footnote{《共产国际执行委员会政治书记处政治委员会会议第395号记录（1934年8月15日）》，《共产国际、联共（布）与中国革命档案资料丛书》第14册，第188页。}


与中共中央决定向西突围几乎同时，国民党方面也预想到中共的可能动向。5月中旬，蒋介石指示：“赣南残匪，将必西窜，酃县、桂东、汝城、仁化、始兴一线碉堡及工事，务请组织西南两路参谋团着手设计，一面准备部队，一面先征集就地民工构筑碉堡为第一线；其次郴州、宜章、乐昌、曲江乃至英德为第二线；先待第一线工作完成，再修第二线，总期于此两个月内，第一线碉堡设法赶成，以为一劳永逸之计”。\footnote{《蒋中正电陈济棠等赣南残共西窜组织西南两路参谋团着手设计构筑（1934年5月18日）》，蒋中正文物档案02020002000036。}6月中旬，陈诚还判断：“第三国际绝不轻许匪军一旦放弃数年来经营成功，自命已成为伪中央苏区根据地之赣巢，而另谋新匪区之盘据。”7月初，他已基本改变看法，认为：“月来我军分路推进，外线合围，愈趋有效；惟匪绝不让我长驱，坐以待亡，其最后策略，似必集中全力，或凭险邀击，或乘虚冲逃。”\footnote{《电呈依据匪情拟进剿方略乞钧裁（1934年6月14日）》、《电呈针对匪军策略贡陈进剿意见（1934年7月1日）》，《陈诚先生书信集·与蒋中正先生往来函电》（上），第136～137页。}双方均已考虑下一步的战略计划。

\subsection{头陂、白水之役}


广昌战役失败后，中革军委的战争指导发生重大变化。第五次反“围剿”前期，中革军委曾明确表示：“我们战术的基本原则，是要求以敏活的机动来实行进攻的战斗。对于占领的支撑点和阵地实行任何的防御，都是不适宜的。”\footnote{项英、王稼祥、彭德怀：《关于进攻坚固阵地战斗的指示（1933年11月7日）》，《项英军事文选》，第261页。}广昌战役后，中革军委虽然继续坚持“短促突击”的运动防御，但此时“短促突击”已越来越变成短距离的战术对抗，战术的机动性大大削减。相应地，中革军委更加重视堡垒的修筑，欲以堡垒对堡垒，与国民党军展开寸土必争的保卫战，阵地防御成为红军基本的战斗方式。之所以如此，当然首先与中共已把大规模的战略转移列入议程，尽力抵御国民党军对苏区的深入，为战略转移赢得准备时间成为其重要考虑之一；其次，此前的屡屡失利使之对继续坚持运动防御的可能性发生怀疑也不能不说是重要原因。国民党方面观察：“赣匪自广昌建宁连城头陂白水相继失陷迭遭惨败后，变更策略，采用堡垒战术，企图节节顽抗，阻我前进。”\footnote{《陈诚罗卓英电蒋中正顾祝同赣匪自广昌建宁连城头陂白水相继失陷后采用堡垒战术企图节节顽抗阻我前进及伪第三十五军团在驿前以北地带利用复杂地形构筑线状阵地设置地雷等（1934-08-30）》，蒋中正文物档案002090300095074。}他们的判断是：

\begin{quote}

土匪向来主张游击战，运动战，但自我军碉堡封锁政策成功以来，亦渐趋重工事，尤其在广昌战役以后，更为重视。盖我军碉堡步步进逼，彼于运动战无机可觅，于是不得不改变方针，遍筑强固工事，与我作阵地战，阻我前进。\footnote{《第三路军在贯桥附近攻击匪强固阵地之所见》，陈诚文物档案702-00038-001，台北“国史馆”藏。}
\end{quote}


撤离广昌后，红军退守头陂、白水一带，构筑工事，准备抗击国民党军的进一步进攻。国民党军在攻击取得阶段性成果基础上，也不得不停步消化，构筑碉堡和封锁线，准备下一步的进攻。同时，重新调整部署：以北路军第八纵队6个师由泰和向兴国方向推进；第六路军第七纵队4个师由龙岗向古龙岗方向推进；第三路军第五纵队4个师首先进占头陂等地，尔后集中第三、第五纵队和东路军第十纵队共9个师，向宁都、白水、驿前、小松市、石城方向推进；东路军第二路军6个师由朋口、连城向汀州方向推进；南路军3个师由筠门岭向会昌、宁都方向推进；另以3个师集结于南丰、广昌地区为总预备队。此时，红军控制区域不断缩小，到1934年6月只剩下瑞金、于都、兴国、赣县、会昌、石城、宁化、长汀等中心区域的寥寥数县。为抵御国民党军的攻击，中革军委决定“六路分兵全线抵御”。7月23日，中革军委通报了国民党军六路进攻的情况和红军六路防御的部署：红三军团六师和红二十一师在兴国西北的沙村地区，抵抗敌周浑元部进攻；红二十三师和江西军区地方部队在兴国东北的古龙岗地区，抵抗敌薛岳部进攻；红五军团十三师在广昌南部的头陂地区，抵抗敌罗卓英部进攻；红三军团四、五师和五军团三十四师、十五师在石城北部的贯桥、驿前地区，抵抗敌樊松甫、汤恩伯部进攻；红一、九军团和红二十四师在闽西的连城、朋口地区，抵抗敌李延年部进攻，红二十二师在会昌筠门岭以北地区，抵抗敌陈济棠部进击。


7月1日，国民党军第三路军完成在广昌的集结，准备继续向南深入，兵锋直指石城。3日，国民党军进占头陂。红军在头陂东部天府山地区集结，一方面瞰制头陂方向国民党军，使其不敢贸然深入；另方面牵制通向石城的另一重要通道白水方向国民党军。


国民党军为打通进攻石城的通道，决定两路会进，第五纵队由头陂向东，攻下天府山，再和由广昌向南的第十纵队会攻白水。9日，国民党军第十四师向天府山发动进攻，占据天府山，随后，红军集中主力展开反攻。当日夜，双方展开激战，“短兵相接，往复肉搏，凡数十次，双方在黑暗中混战，达五六小时”。\footnote{《赣粤闽湘鄂北路剿匪军第三路军五次进剿战史》（上），第十章，第7页。}半夜时分，红军已占领除最高点外的大部分阵地，国民党军第十四师被压缩到山顶最后一线，面临着被歼灭的命运。10日凌晨，闻讯增援的国民党军第十一师赶到战场。国民党军在援军支援下发起反攻，红军功亏一篑，未能达到歼灭敌人的目标，于黎明前退出战斗。此役，红军给对方以较大打击，国民党军伤亡达400余人。


为配合天府山地区第五纵队的进攻，国民党军第十纵队也于9日开始向赤水地区进犯。10日，在空中力量掩护下，攻占赤水，至此，国民党军已经打开通向石城的第一道门户。红军主力退往中司、驿前一带，前锋则位于白水西南的大寨脑一线，在此构筑工事，准备抗击国民党军的进一步进攻。


7月中旬，国民党军第十纵队开始向大寨脑地区进兵。红三军团在此以堡垒战术与国民党军周旋，虽然延缓了国民党军的进攻步伐，但由于国民党军的炮火优势，损失颇大，红军堡垒作用有限：“一切条件支点全不坚固，受不起一个迫击炮弹，打不出枪，手榴弹投不出去，没有副防御，不能抑制敌于我手榴弹火力下”。\footnote{《彭杨关于三军团向半桥良田集中问题及行进中与战斗情况的报告（1934年7月23日）》，《中央苏区第五次反“围剿”》（下），第86页。}22日，国民党军付出近300人伤亡的代价，攻占大寨脑、鸭子岭，兵锋直指贯桥、驿前一带红军主力。


这一时期，国民党军虽然凭借绝对优势的兵力、装备，攻击屡获成功，但继续进兵其实也承受着很大压力。随着其向苏区内部的不断深入，后勤给养成为亟待解决的问题。同时，因为连续作战，水土不服，传染病对其形成很大杀伤，国民党军将领回忆：

\begin{quote}

广昌、石城之间，匪我相持缠斗，达三个月之久，大军云集于狭窄地区，时值酷暑，部队皆无蚊帐，又以蔬菜稀少，饮水不洁，及在阵地露宿等因素，致痢疾及疟疾患者极多，死于道旁者累累。师野战医院病患达千余，人满为患；我团两千余人，只有七人未患疟疾，我亦染患此病。疾病对于战力消耗，数倍于作战受伤，军队健康之维护，实为一极其重要之事。\footnote{《石觉先生访问记录》，第71页。}
\end{quote}


该回忆道出了国民党军的实况，让我们看到中共方面处于困境时，国民党方面其实也并不轻松。稍后，陈诚在给妻子的家书中写道：“此次至各处视察，情状极惨，沿途死病士兵夫无人处理者，不知其数。自广昌至石城，每日死亡计二百以上，各师患病者占三分之二，行营及总司令部无人过问，而各部又无法处理。”\footnote{《至各处视察沿途死病士兵夫不知其数各部军风纪亦不甚好（1934年10月13日）》，《陈诚先生书信集·家书》（上），第283页。}战局进行到此时，双方都各有苦衷，国民党军也有骑虎难下之势。蒋介石可资依赖的主要是人多势众、实力雄厚，抗击打能力相对较强。然而，当时蒋介石又不能不顾虑地处赣南红军之后虎视眈眈的粤桂势力，如果与红军作战消耗太大，必将影响其应对粤桂挑战的能力。所以蒋介石此时表面看气势汹汹，内心则谨慎有加，苦心焦虑谋划的，是代价最小的结局。

\subsection{红七军团北上与红六军团西征}


广昌战役后，战场上不利的形势使红军的战略转移更加迫在眉睫。6月下旬，中共中央政治局召开会议，落实共产国际的有关指示，决定派出红六军团、红七军团分别往西、往北。\footnote{参见金冲及主编《毛泽东传》（上），中央文献出版社，1996，第327页。}北上和西进计划在共产国际同意中央苏区突围的电文中有明确交待：“（1）为防备不得不离开，要规定加强在赣江西岸的基地，同这些地区建立固定的作战联系，成立运粮队和为红军建立粮食储备等；（2）现在就用自己的一部分部队经福建向东北方向发起战役，以期最后这些部队成为将来闽浙皖赣边区苏区的骨干力量。”\footnote{《共产国际执行委员会政治书记处政治委员会给埃韦特和中共中央的电报（1934年6月16日）》，《共产国际、联共（布）与中国革命档案资料丛书》第14册，第143页。}


根据共产国际的指示，中革军委决定抽调主力红军一部，组建抗日先遣队北上。6月下旬，红七军团军团部和所属第十九师受命从福建连城回到瑞金，待命北上。7月初，中革军委命令红七军团改编为北上抗日先遣队，北上深入到闽、浙、皖、赣地区“发展游击战争，创造游击区域”，并力争“建立新的苏维埃的根据地”，以牵制国民党“围剿”中央苏区力量，促其“调回一部到其后方去”。\footnote{《中央政治书记处、中央政府人民委员会、中革军委会关于派七军团以抗日先遣队名义向闽浙挺进的作战训令》，《中共中央文件选集》第10册，第332页。}红七军团由寻淮洲任军团长、乐少华任政治委员、粟裕任参谋长、刘英任政治部主任，突击补充了2000多名新战士，全军团共计6000余人。


7月6日晚，红七军团从瑞金出发，经福建长汀、连城、永安，于下旬进入闽中。随后进到尤溪地区。29日，攻占闽江南岸的尤溪口。此时，中革军委突然改变原定计划，命令红七军团东进，占领水口，威胁并相机袭取福州，企图调动和吸引更多的国民党军回援。遵照中革军委的命令，红七军团主力改向水口、福州方向前进，一部兵力继续北上。8月1日，攻占福州近郊水口，威胁福州。同时正式向部队宣布，红七军团对外以“中国工农红军北上抗日先遣队”的名义活动。2日，红七军团从水口绕道大湖向福州进发，7日进抵福州西北郊，当晚对福州发起试探攻击。对此，国民党军预有准备，第八十七师主力回防福州，第四十九师也经上海驰援，红军缺乏攻坚装备，兵力有限，被迫撤出战斗。红军进迫福州，中革军委意在围魏救赵，以此缓解闽西一带国民党军对苏区的压力，对此，蒋介石一度因难以摸清红军的意图：“闽北匪情是否另辟匪区，抑仅系牵制东路而窜扰耶”，而为之“心颇不安”。\footnote{《蒋介石日记》，1934年8月2、3日。}不过，红军北上兵力毕竟有限，国民党方面很快看清了这一点：“东路军总部探知其诡谋，密调劲旅，兼程进剿，闽西国军仍向石城长汀一带迈进，并末因之移动”；而先遣队的实力却因之“完全暴露”：“赤匪的诡计和实情，已经给我方看破了，所以我方对前方剿赤的军队，不调一兵一卒，以致松懈前方清剿的工作，只由他方面调来有余的兵力，已经可制这些残赤的死命了。”\footnote{《福州严重之一夜，匪分两路进袭谋占省垣》，1934年8月15日福建《民报》；陈肇英：《剿灭残赤的必胜观》，《党务周报》第1卷第17期，1934年8月2日。}


攻打福州失败后，红七军团向闽东地区转移，8月下旬，进入浙西南地区。9月初，转进至闽北苏区休整、补充。9月4日，中革军委来电，对原定行动计划作出补充：要求：“继续彻底的破坏进攻我红十军及闽北苏区敌人的后方”；“在闽浙赣皖边境创造广泛的游击运动及苏区根据地”，\footnote{《朱德、周恩来、王稼祥关于七军团作战计划补充指示致寻淮洲、乐少华并曾洪易电（1934年9月4日）》，中共福建省委党史研究室等编《中国工农红军北上抗日先遣队》，中共党史出版社，1990，第85页。}在执行上述两项中心任务时，首先在浙西一带活动，破坏交通和铁路线，支援中央苏区的反“围剿”战争。9月7日，中革军委又电批评：“在近几天来，七军团是在同闽北部队作盲目的儿戏，以致使较弱而孤立的敌人进到渔梁地域。七军团未完成自己的战斗任务。”\footnote{《朱德关于七军团行动问题致寻淮洲、乐少华电（1934年9月7日）》，《中国工农红军北上抗日先遣队》，第89页。}从当时实际情况看，中革军委出动红七军团北上，旨在使其于浙西一带牵制国民党军，尤其红军战略转移进入最后准备阶段时，这一要求更加迫切。但是，红七军团深入国民党军占领区域后，所遇困难确实相当严重，保证部队的生存成为执行任务的必要前提。


面对不利局面，红七军团仍然尽力完成中革军委赋予的任务。根据中革军委指示，9月9日，红七军团转到浙西一带活动，18日抵遂安白马，接中革军委电令：“主力应即向遂安前进，以袭击的方法占领该城而确实保持之于我军手中。”\footnote{《朱德关于七军团袭占遂安县城问题致寻淮洲、乐少华电（1934年9月18日）》，《中国工农红军北上抗日先遣队》，第102页。}中革军委还要求，攻占遂安后，“于安徽边境淳安、寿昌、衢县、开化的范围内发展游击战争和苏维埃运动”，并在“浙皖边境，约在徽州、建德、兰溪、江山、屯溪地域建立新的苏维埃根据地”。\footnote{《朱德关于七军团在浙西的政治工作和作战任务致寻淮洲、乐少华并刘畴西、聂洪钧电（1934年9月18日）》，《中国工农红军北上抗日先遣队》，第103页。}根据这一命令，先遣队派出侦察分队，至遂安附近了解敌情，准备袭取遂安。但国民党军在此严阵以待，国民党方面报纸报道：“省方开往部队昨已安全到达遂安县城，防地并已布置就绪，龙游、寿昌、淳安、衢县均已派定相当部队向遂安西南北三方面严密防堵，东西则有陈调元派队开入遂境，负责进剿，包围之势已成”\footnote{《窜遂残匪被包围国省军今日总攻，三日内即可将股匪全部消灭》，1934年9月21日《东南日报》。}，红军军事上处不利态势。另外，从政治上看，浙西一带群众基础薄弱，“常、遂交通一带之群众很差，几次逃跑一空，连向导均无人找到……因此，在我们的行动上均难以秘密”。\footnote{《曾洪易关于七军团对打击敌左右纵队之意见致朱德电（1934年9月19日）》，《中国工农红军北上抗日先遣队》，第107页。}七军团指挥员报告：

\begin{quote}

根据我们最近十天在浙西行动艰苦的经验，由于浙西电话、交通事业的发达，我们的行动敌人很快就知道，飞机每天都能跟追随，敌人队伍运动与转移均极迅速。并党与群众工作基础完全没有，伤病员安置极度困难。我们最近在浙皖边境，寄存群众家中一批伤员，据闽浙赣谍息全被敌搜出，影响士兵战斗情绪、战斗力极大。\footnote{《曾、寻、乐关于七军团与皖赣、皖南地方武装配合，创造浙皖新苏区致党中央、中革军委的请示电（1934年9月30日）》，《中国工农红军北上抗日先遣队》，第120页。}
\end{quote}


9月下旬，军团领导决定放弃在浙西建立根据地的计划，迅速摆脱敌人，改向皖赣边前进。10月初，到达皖赣苏区休整，此时，“部队已不足两千人”。\footnote{《红军抗日先遣队北上经过的报告》，《闽浙皖赣边区史料》，第61页。}


由于国民党军加紧向皖赣苏区进攻，10月21日，中革军委电令红七军团“逐步转回闽浙赣苏区”。\footnote{《朱德关于七军团取石门后转回闽浙赣苏区致寻淮洲、乐少华电（1934年10月21日）》，《中国工农红军北上抗日先遣队》，第135页。}遵照中革军委命令，红七军团在浮梁、德兴间通过两道封锁线，进入闽浙赣苏区，并于11月初与红十军会合。11月4日，中革军委决定红七军团与红十军合编为红军第十军团，下辖两个师，红七军团改编为第十九师，红十军改编为第二十师。由于红七军团出发后就出现领导层的种种问题，曾洪易消极避战，寻淮洲与乐少华则“不断的闹私人意见”，领导层其他成员间也矛盾重重，“不是你争我吵，便是你走我溜，甚至闹成打架”，\footnote{刘英：《北上抗日与坚持浙闽边三年斗争的回忆》，《闽浙皖赣边区史料》，第75、85页。}因此，军团领导层根据中共中央意见进行调整，刘畴西任军团长兼第二十师师长，乐少华任军团政治委员兼第二十师政治委员，寻淮洲任第十九师师长；并决定：闽浙赣省苏维埃主席方志敏兼任闽浙赣军区司令员，曾洪易任中共闽浙赣省委书记兼军区政治委员，粟裕任军区参谋长。


11月18日，中共中央决定成立以方志敏为主席的军政委员会，指示其向皖南挺进，“创造皖浙边苏区”。\footnote{《中央军区关于十军团行动任务及军政委员会组成的决定致方、曾、刘、乐、寻、聂电（1934年11月18日）》，《中国工农红军北上抗日先遣队》，第145页。}此后，红十军团在中央苏区主力红军转移进行长征的情况下，勉力支撑，往返于闽南与皖浙赣之间，作战十余次，减员达1/3以上，坚持奋斗至次年1月，“由于军政首长决心的不够”，\footnote{刘英：《北上抗日与坚持浙闽边三年斗争的回忆》，《闽浙皖赣边区史料》，第83页。}“没有下断然决心”\footnote{方志敏：《我从事革命斗争的略述》，《方志敏文集》，江西人民出版社，1999，第92页。}冲破国民党军的封锁线，在婺源怀玉山陷入重围，主力分割成数段。1月下旬，大部溃散或牺牲。方志敏、刘畴西被俘就义，只有800余人突围到闽浙赣边继续坚持游击战争。


红七军团北上，如朱德所指出的：“是准备退却，派先遣队去做个引子，不是要北上，而是要南下”，\footnote{《粟裕战争回忆录》，第134页。}是为配合主力红军即将进行的战略转移服务的。在派出北上先遣队稍后，中共中央和中革军委又决定湘赣苏区的红六军团进行西征，作为中央红军转移和西征的先遣队。两路部队，“一路是探路，一路是调敌”，\footnote{《周恩来在中共中央政治局会议上的发言提纲（1943年11月15日）》。}意图明确。


第五次反“围剿”开始后，湘赣苏区和湘鄂赣苏区的红六军团遭到国民党军围攻。虽然红军在梅花山、沙市等战斗中给进犯之敌以严重打击，但仍然难以抵挡国民党军的步步进逼和蚕食，根据地日渐缩小。1934年6月底、7月初，东华山、松山战斗失利后，湘赣苏区中心区域被占领，红六军团被分割、压缩在遂川、万安、泰和三县交界方圆仅数十里的狭小区域内，处境艰难。7月23日，中共中央书记处和中革军委命令，鉴于“敌人正在加紧对于湘赣苏区的封锁与包围”，“六军团继续留在现地区，将有被敌人层层封锁和紧缩包围之危险，而且粮食及物质的供给将成为尖锐的困难，红军及苏区之扩大受着很大的限制”，要求红六军团离开湘赣边，转移到湖南中部开展游击战争并创立新的苏区。命令强调：“六军团以自己在湘中的积极的行动，消灭敌人的单个部队，最大的发展当地的游击战争与土地革命。直至创立新的苏区，给湘敌以致命的威胁，迫使他不得不进行作战上及战略上的重新部置。”\footnote{《党中央书记处、中革军委关于红六军团向湖南中部转移给六军团及湘赣军区的训令》，《中共中央文件选集》第10册，第355～356页。}中共中央下令红六军团转移，除保存湘赣苏区现有力量的考虑外，更希望红六军团西进后，通过与湘鄂赣苏区红二军团建立可靠联系，在湖南中部开辟根据地，形成江西、四川两苏区联结的前提，这样将极大地便利主力红军的西征。


8月7日，任弼时等率领红六军团主力由江西遂川突围西征。据任弼时等报告，当时红六军团第十七、十八师共有6800余人，长、短枪及机枪3200余支。\footnote{《任弼时、王稼祥致中革军委关于红六军团状况的报告（1934年7月31日）》，《湘赣革命根据地史料选编》（下），江西人民出版社，1984，第729页。}根据中革军委指示，红六军团出发后，11日到达湖南桂东，次日正式宣布军政委员会和红六军团成立。军政委员会为西征最高领导机关，任弼时任主席，萧克、王震、李达、张子意分任军团军团长、政治委员、参谋长、政治部主任。军团辖红十七、十八两个师。


9月初，红六军团进入湘西南地区。中革军委指示其“协同二军团于湘西及湘西北地域发展苏维埃及游击运动，并于凤凰、松桃、干城、永绥地域建立巩固的根据地，其后方则背靠贵州，以吸引更多湘敌于湘西北方面”。\footnote{《中革军委致电六军团关于今后活动地区及其任务的指示（1934年9月8日）》，《湘鄂川黔革命根据地历史文献汇集》，湖南省档案馆，1984，第142页。}随后，该军团在城步一带活动。由于国民党军集中湘、桂、黔三地兵力施以压迫难以立足，被迫向贵州国民党军相对较弱地区进发，到达黔中瓮安一带。10月4日，中革军委命令红六军团向北至黔北印江与红二军团会合。北进途中，7日在石阡遭遇国民党军堵截，损失惨重，此后一直处于国民党军的围追堵截中，部队不断减员，“子弹极缺乏”。\footnote{《何键报告书》，《湘赣革命根据地》（下），第1293页。}此时，红三军领导人贺龙、关向应等得知红六军团准备北上会合，率红三军主力南下接应。24日，红六军团主力突破国民党军封锁线，在印江与红三军会师。随后，红三军奉命恢复红二军团番号，贺龙任军团长，任弼时任政治委员，辖两个师；红六军团由萧克任军团长，王震任政治委员，暂时编为3个团。两军成立以任弼时、贺龙、关向应为首的红二、六军团总指挥部。

\subsection{高虎脑、驿前之役}


1934年8月，国民党军经过短暂休整后，开始执行其进军石城计划。国民党军第三路军确定的推进总方针是：“以进占石城与东路军会取长汀之目的，拟即由白水附近向石城逐次筑碉前进，并准备随时与匪主力决战，以达成任务。推进计划分为三期，第一期进占驿前，第二期进占小松市，第三期进占石城。”\footnote{《赣粤闽湘鄂北路剿匪军第三路军五次进剿战史》（上），第十章，第17页。}8月5日，国民党军开始向贯桥一带红军阵地发动攻击。为阻止国民党军的快速进展，红三军团奉命撤至贯桥地区依托高虎脑有利地形，阻击南进石城之敌。红五军团于广昌、宁都之间布防，阻截西进宁都之敌。具体部署是：红三军团第五师为正面，在鹅形、蛤蟆寨、高虎脑设防；第四师为第五师的右翼，在画眉山、东家边、老寨、宝峰山、腊烛形设防；红五军团第三十四师为第五师的左翼，在船形、高脚岭、香炉寨设防。红军在高虎脑、万年亭到驿前20余里区域内构筑了以5个支撑点为骨干的立体防御工事，工事“外壕深广，各堡垒前均有三铁鹿砦，鹿砦之外，遍置竹钉，竹钉之外，埋设地雷”，\footnote{《赣粤闽湘鄂北路剿匪军第三路军五次进剿战史》（上），第十章，第29页。}防御相当严密。其中，高虎脑位于驿前北面，背靠两座大山，地势险峻复杂，紧紧卡住绕镇而过的广（昌）石（城）公路，是红军纵深防御体系的重点。红军在阵地前沿大筑工事：

\begin{quote}

挖了三道一米多深的堑壕。从山上砍来竹子，削成二、三寸的竹钉，用火锅炒一下后钉在木板上，放在堑壕前面，上面加以伪装，还挖了许多陷阱，设置了路障，并在阵地前沿埋设了地雷群。为防敌炮火和飞机轰炸，我们锯倒粗大的松树，作为工事的顶盖，再铺上二、三尺厚的泥土封顶，将各个防御工串构筑得像碉堡一样坚固。各工事之间还修挖了出击的通道。\footnote{温宗学：《回忆高虎脑阻击战》，《石城文史资料》第3辑，政协石城文史资料委员会，1990，第40页。}
\end{quote}


防御准备相当充分。


6日，当国民党军第八十九师对高虎脑展开全面进攻时，红军利用有利地形和坚固工事对其实施沉重打击：“我炮兵及机枪的短距离射击，手榴弹与地雷的爆烈在距离我约百米远的反斜面上，给敌以群集的牺牲。”\footnote{《三军团半桥附近战斗详报（1934年8月5日至7日）》，《中央苏区第五次反“围剿”》（下），第97页。}国民党军“团长以下伤亡甚重，且时近黄昏陆空不易协同”，红军“控置后方之部队得以从容活动”。\footnote{《国民党东路军第十纵队汤恩伯部与中央苏区红军在赣南一带战斗详报》，《中华民国史档案资料汇编》第5辑第1编军事（4），第348页。}次日，国民党军在火炮支援下，首先向鹅形发动强攻，红军“守备支点被炮火毁坏，机枪子弹用尽”，\footnote{《三军团半桥附近战斗详报（1934年8月5日至7日）》，《中央苏区第五次反“围剿”》（下），第100页。}被迫退出鹅形。接着国民党军又于下午以猛烈炮火掩护，不惜代价猛攻高虎脑，在付出重大牺牲后终于攻占高虎脑主阵地，红军退向驿前，国民党军进占贯桥。高虎脑战斗，倾泻在红军阵地上的各种炮弹和炸弹达2000多发，手榴弹达5000余枚。红军通过顽强抗击，使国民党军攻占不足3000米水平距离的山地，却付出了伤亡3000人的惨重代价，\footnote{国民党方面战史的统计数字是，整个贯桥高虎脑之役，伤亡1500人，其中军官100余人，阵亡团长1人、营长3人。}精锐部队第八十九师丧失了战斗力，不得不退出战斗。不过，此役红军伤亡达1373人，另有40人失踪。其中第五师“原来两千多人的大队伍，只剩下不到一千人，有的连队只剩二十多人”，\footnote{谢振华：《第五次反“围剿”中的高虎脑战斗断忆》，《军史资料》1985年第7期。}红五师政委陈阿金及众多指挥员牺牲。以当时国共双方兵力对比，红军仍然难以承受这样的消耗。


经过这次战役，国民党军对红军的防御能力有了新的认识，顾祝同向蒋介石报告，由于红军防御方式的变更，使“飞机侦炸及炮兵射击均不易生效力”。\footnote{《顾祝同电蒋中正此次贯桥之役我军伤亡奇重此次匪军方略变更因此我军飞机侦炸及炮兵射击均不易生效力（1934年8月10日）》，蒋中正文物档案002090300087290。}陈诚判断：“此次贯桥之役，匪构成数线强固阵地，并以比较精锐之向称突击部队顽强固守，似有因运动战（无）机可寻，转而采取阵地战之势。”为此他调整部署，一面完成碉堡修筑，“再行进展，并研究对于攻击匪强固阵地之对策，以利作战”；一面紧急向南京方面提出：“山迫炮效力极小，昨日飞机轰炸命中非常精确亦不能予以破坏。”\footnote{《陈诚1934年8月8日致顾祝同电》，《中华民国史档案资料汇编》第5辑第1编军事（4），第104页；《陈诚罗卓英电蒋中正等据报共军大岭格一带碉堡异常坚固请抽调卜福斯山炮俾供施行平射（1934年）》，蒋中正文物档案002090300083086。}要求增调新型山炮，加强火力，对付红军更加顽强的碉堡战术。


面对国民党军的厉兵秣马，红军也在调整兵力，修筑堡垒，彭德怀等前线指挥员致电朱德强调，面对国民党军即将发动的进攻，“我皆不应与之决战，以四、五两师三十四师防御动作，依靠支点高度消灭其有生力量，并以短促的反突击消灭敌人部分”。\footnote{《彭杨关于不与敌罗汤樊纵队决战之部署报告（1934年8月9日）》，《中央苏区第五次反“围剿”》（上），第203页。}次日，中革军委复电指示：“三、五军团抗击陈路军的任务仍以防御及部分的反突击来阻击及削弱敌人，但不必进入决战以节约我军兵力。”\footnote{《三军团驿前及其以北战斗详报（1934年8月10日至30日）》，《中央苏区第五次反“围剿”》（下），第110页。}8月13日，经过数天整顿，国民党军开始向驿前发动攻击。14日，两军在万年亭一带展开血战，“始则手榴弹互相投掷，密如骤雨，继则白刃搏斗，往反冲突，血肉横飞，极其惨烈”。\footnote{《赣粤闽湘鄂北路剿匪军第三路军五次进剿战史》（上），第十章，第35页。}国民党军第八十八师一部突入红军阵内后，被红军切断予以全歼，余部在飞机掩护下退回原阵地。此后，“因两次攻坚伤亡过大”，\footnote{《三军团驿前及其以北战斗详报（1934年8月10日至30日）》，《中央苏区第五次反“围剿”》（下），第117页。}国民党军不得不暂时停止进攻，直到其炮兵第一旅一营炮兵携带新式山炮到达前线，空军主力也向南集结，完成攻击部署后，才开始新一轮进攻。


8月28日，国民党军汤恩伯、罗卓英、樊崧甫三部，由罗卓英统一指挥，分左中右三路开始新一轮攻击。凌晨4时，国民党军集中三个炮兵营、飞机数十架，首先对红军阵地展开猛烈轰击，压制红军突击部队，同时“敌步兵预先利用暗夜接近出击地，此时乘机、炮轰炸及烟雾弥漫之际，猝然奇袭我阵地”。\footnote{《陈伯钧日记（1933～1937）》，第285页。}“这一来，我们的短促突击就根本用不上了”，\footnote{萧华：《第五次反“围剿”战斗中的大寨脑、驿前战斗》，《回忆中央苏区》，第476页。}“绵亘廿余里之匪阵地带，一弹指顷，已深深埋葬于烟幕之下”。\footnote{《赣粤闽湘鄂北路剿匪军第三路军五次进剿战史》（上），第十章，第39页。}在国民党军优势火力的强大攻击下，红军修筑的防御工事效果大打折扣，受对方“步炮空火力摧损，死伤极大”，保护山守军更“因被炮弹爆烟遮蔽视界机关枪突然发生故障”，致使敌军侵入阵地后守军都未有反应。同时，由于“忽视了敌人十五天的严重准备，所以在敌人开始突击时便暴露了我们精神松懈的极大弱点”。\footnote{《国民党北路军顾祝同部与中央苏区红军作战情形报告书》，《中华民国史档案资料汇编》第5辑第1编军事（4），第126页；《三军团驿前及其以北战斗详报（1934年8月10日至30日）》，《中央苏区第五次反“围剿”》（下），第121、129页。}许多部队在国民党军突然发起攻击时，正在熟睡、出操或在阵地下面休息，第十三师三十九团阵地内“无人守备，只有一个哨”，一遇进攻，措手不及，阵地纷纷失守。29日，国民党军已逼近驿前，凌晨3时许，中革军委致电彭德怀、杨尚昆，指示：“关于驿前之最后扼守或放弃由彭杨作战况决定，但如撤退必须有计划有掩护的进行。”\footnote{《三军团驿前及其以北战斗详报（1934年8月10日至30日）》，《中央苏区第五次反“围剿”》（下），第120、126页。}30日，红军撤离驿前，退向小松市方向。是役，红军损失重大，伤亡两千余人。


高虎脑、驿前之役，国共双方消耗均大，国民党军由于处于仰攻状态，伤亡更重。三军团指挥员彭德怀和杨尚昆在总结高虎脑战斗的经验时谈道：“假如我们在几百里距离的赤色版图上，一开始就使敌人遭受这样的抵抗，而给敌人消耗量当是不可计算的，要记着广昌战斗我们有生力量的消耗是数倍于敌的。”\footnote{《三军团半桥附近战斗详报（1934年8月5日至7日）》，《中央苏区第五次反“围剿”》（下），第103页。}他们的总结也在国民党方面将领那里得到印证，汤恩伯当时在电文中汇报：“匪占各高地之工事极度强固，飞机炸弹及炮弹多未破坏，工事前面埋有地雷，俟国军接近拉线即爆发。匪兵守工事者甚少，皆在周侧隐伏，如国军接近工事则从两侧包围反攻。”\footnote{《汤恩伯电蒋中正》，蒋中正文物档案002090300081286。}时任国民党军第四师团长的石觉后来谈到此役时说：“我军总共耗时七十二天，只进展廿九公里，表面上我们成功，事实上敌人利用这段时间向西南方面突围逃逸。这在历史上留下一个大疑问，何种作法为对，实难断定。”\footnote{《石觉先生访问记录》，第68页。}


在进攻驿前同时，龙岗方向的国民党军第六路军也向古龙岗展开进攻。8月，蒋介石就该线“进剿”计划下达手令，令其9月份“到达石城、古龙岗与兴国”。\footnote{《蒋中正总统档案·事略稿本》第27册，第380页。}8月上旬，国民党军第三纵队向石城进攻；第六纵队向兴国古龙岗进攻，第八纵队则向兴国的老营盘迅速推进。第八纵队出动后，先后占领沙村、杨公山等地，红六师在沙村阻击，从7月26日到8月25日，不到一个月时间里第八纵队伤亡、失踪1720人。\footnote{《曹德卿、徐策报告（1934年8月27日）》。}9月1日，国民党军控制由泰和往兴国的必经要冲老营盘。9月4日，第八纵队向新田、蓝田圩一线发动进攻。红五、红一军团先后到达该地区展开顽强防御，努力迟滞国民党军的迅速南下，保住红军往西的出口。经过半个多月的激战，红军虽有效控制了国民党军的南进步伐，但损失较大，不得不步步后撤，9月22日，第八纵队进据新田、蓝田圩一线，30日，进占距兴国县城仅30公里的高兴圩，直逼兴国城。


进攻古龙岗的国民党军第六纵队于9月11日开始发起攻击，12日双方在雄岭下“激战终日”，\footnote{1934年9月14日《江西民国日报》。}打成相持。21日，中革军委决定将在雄岭下防御的红二十一、二十三两师编为红八军团。22日，国民党军出动4个师部队向红八军团大举进攻，23日，基本控制雄岭下地区。红八军团先退兴国南坑，再退天子嵊一线。


江西红军连遭失利时，9月初，红一、红九军团与红二十四师在福建温坊取得局部战斗的胜利，给苏区军民稍许安慰。温坊位于长汀东南部，8月下旬，国民党军李延年纵队进占朋口一带后，计划经温坊推进长汀。中革军委对国民党军这一动向高度注意，决定在此方向集中兵力对敌实施打击。8月26日，朱德签发《我军目前的作战任务及行动部署》，命令集结在宁化曹坊的红一军团以打击李纵队为目标，即速赶至长汀童坊、河田地区，与红九军团和红二十四师会合，待命出击。31日，朱德再次命令红一、红九军团与红二十四师主力集结“温坊中屋村间进行突击李纵队的任务”。\footnote{《朱关于一、九军团突击李纵队行动部署的指示（1934年8月31日）》，《中央苏区第五次反“围剿”》（上），第203页。}9月1日，国民党军第三师第八旅大意轻敌，离开朋口单兵向其西十余里的温坊突进，准备在此构筑碉堡。该部到达温坊后，对近在眼前的红军主力仍没有感觉，报告：“当面之匪，系伪廿四师，并无其它匪情。”\footnote{《国民党军第3师李玉堂部与中央苏区红军在闽西赣南一带战斗详报》，《中华民国史档案资料汇编》第5辑第1编军事（4），第381页。}国民党军这一违背常规的举动，连红军将领也感困惑，红九军团报告：

\begin{quote}

当时我们对敌情判断确有怀疑之处，认为一军团情报不确，因为敌人从来没有用两团兵力而远离其堡垒前进，李敌纵然胆大也不敢以两团军力挺出其封锁线十余里之温坊。纵或有之，仅系游击性质，绝不敢在该地构筑工事，今既构筑阵地，必有较大的兵力在后尾跟进或已隐蔽集结于温坊附近，必系我军的侦察不确，一时为其所欺骗。\footnote{《九军团温坊附近两次战斗详报》，《中央苏区第五次反“围剿”》（下），第138页。}
\end{quote}


这既说明国民党军在温坊行动的轻率，也显示屡遭失利后，红军将领自信心已大受影响。


确认国民党军实为单兵突进后，红军决定抓住机会，以一部迂回到敌军背后，截断其后路，红一军团、红九军团与红二十四师分路出击，将进驻温坊的国民党军紧紧包围。战斗至次日凌晨2时许，国民党军阵地被全线突破，两个团遭歼，第三师第八旅旅长许永相仅以身免，部队“逃脱的仅约百余”。\footnote{《二十四师温坊两次战斗详报》，《中央苏区第五次反“围剿”》（下），第134页。}该师事后总结，战斗失败的主要原因是：“各级官长轻敌之心太盛，阵地不能慎密布置，工事不能迅速完成，以致猝遇匪袭，噬脐莫及。”\footnote{《国民党军第3师李玉堂部与中央苏区红军在闽西赣南一带战斗详报》，《中华民国史档案资料汇编》第5辑第1编军事（4），第384页。}


在各部均取得进展时，福建却传败讯，令早已对福建方面进展缓慢的蒋介石深为恼火，他特电蒋鼎文令其牢记稳扎稳打战术，声称背离这一战术“虽胜必罚”。\footnote{《蒋中正电蒋鼎文东路向河田长汀前进须稳重师长等不照剿共战术虽胜必罚（1934年9月4日）》，蒋中正文物档案02020002000052。}同时派顾祝同到龙岩，在给蒋鼎文的电报中说是协助，而给顾的电文中则称可令蒋“向前指挥”\footnote{《蒋中正电顾祝同东路失利往龙岩指挥可使蒋鼎文向前负责或调其来赣（1934年9月9日）》，蒋中正文物档案02020002000055。}或调赣，实际是以顾代蒋，贬蒋鼎文为前敌总指挥。随后，蒋介石以“欺上陷下，又复临阵脱逃”\footnote{《蒋中正电蒋鼎文第八旅长许永相应就地枪决（1934年9月20日）》，蒋中正文物档案002010200119019。}之罪将第八旅旅长许永相枪决，第三师师长李玉堂由中将降为上校。一家欢乐一家愁，在蒋介石大感痛心时，林彪、聂荣臻对此战结果自是颇为欣慰：“苦战一年，此役颇可补充。”\footnote{《温坊战斗之战果与一军团守备布置的报告》，《中央苏区第五次反“围剿”》（下），第137页。}


9月3日，国民党军根据东路军总指挥蒋鼎文的部署，再次出动第三、第九师3个团进攻温坊，显然，蒋鼎文对红军在温坊地区的兵力没有足够的估计，在两个团遭消灭后，仅仅增加一个团兵力，犯了逐次增兵的兵家大忌。对此，红军指挥部有充分估计，9月3日凌晨，朱德命令“我一、九军团及二十四师主力在有利条件下应在温坊阵地前给敌以短促突击以消灭其先头部队”。\footnote{《朱德致林彪、聂荣臻急电（1934年9月3日）》。}3日上午，国民党军先头部队第九师第五十团进至温坊东北，红一军团截其归路后，实施包围并予以全歼，同时，第三师后续部队也遭到红军的痛击。两次战斗，红军将领回忆共歼敌2000余人，俘虏2400余人，仅红二师四团就俘虏1600人。\footnote{《耿飚回忆录》（上），第155页。}国民党方面的数据是：自身伤811人，阵亡395人。\footnote{《蒋鼎文电蒋中正（1934年9月16日）》，蒋中正文物档案002090300081193。}红军伤亡仅900余人，其中一军团伤亡600余人。\footnote{《林聂关于温坊附近战斗情况及四日我在中屋村等地的部署情况报告》，《中央苏区第五次反“围剿”》（下），第148页。}


温坊战斗的局部胜利不足以扭转战局。由于江西战事紧急，红一军团在战斗胜利后奉命开往江西，仅红九军团与红二十四师继续留守松毛岭、白衣洋岭一线。9月27日，国民党军第九、第三十六师开始向松毛岭、白衣洋岭发动进攻，29日，红军阵地全线失守。随后，红军根据转移命令，主力逐渐向江西集中，福建方面战事大体结束。

\section{红军实施转移}

\subsection{红军长征前的准备}


国民党军8月底攻占驿前后，照例修建碉堡、构筑公路，并将碉堡线逐渐向小松市红军阵地伸展，其前锋部队则试探性地对红军阵地展开攻击，蚕食红军防线。在屡遭失败，国民党军前锋步步向红军中心区进逼时，红军实施战略转移已势在必行。


为集中策划、调度即将到来的战略大转移，1934年夏，中共中央书记处会议决定由博古、李德、周恩来组成“三人团”，负责筹划红军战略转移的有关事宜，一系列有关战略转移的计划、行动陆续展开。遵义会议决议提到，中革军委在“八九十三个月战略计划”中，已经“提出了这一问题，而且开始了退出苏区的直接准备”。\footnote{《中央关于反对敌人五次“围剿”的总结的决议》，《中共中央文件选集》第10册，第466页。决议同时批评中革军委在“五六七三个月战略计划”中没有提出这一问题，这种批评显得牵强，因为“五六七三个月战略计划”制定根据李德的回忆在4月，此时转移计划还未确定，中革军委不可能未卜先知。}最初预定的突围时间是10月底11月初，因为“根据我们获得的情报，蒋介石企图在这期间集中力量发动新的进攻，突围的日期选择在这时，必然会使敌人扑个空”；“从华南地区的地理气候上来考虑，这也是行军和作战的最有利的时间”。\footnote{〔德〕奥托·布劳恩：《中国纪事1932～1939》，第112页。}中共叛将杨岳彬在给国民党方面的呈文中透露了中共的这一意图：“今年四五月，赤党伪中央，已从匪军中，调出湘南籍之干部多名，潜回湘南各县，布置交通路线，并闻有从湘赣边匪区，抽调一部匪军，编为湘南游击队，窜扰湘南之说。”\footnote{《匪情实录》，赣粤闽湘鄂“剿匪”军东路总司令部《东路月刊》第3、4期合刊。}


根据向西转移的方针，从1934年7、8月份开始，红军已开始部署重要物资、资材的西运事宜。7月下旬，红十五师奉命开往福建方面，“到福建搬运胜利品”，\footnote{《彭绍辉日记》，1934年7月24日，第19页。}所谓胜利品，实际就是准备搬运的重要资材。7、8月份，红一军团主力开赴闽西掩护红九军团安全转运资材到赣南。8月中旬，朱德电令一、九军团，要求其在闽西苏区“确实掩护资材”\footnote{《朱关于掩护资材西运的行动部署给一、九军团的指示（1934年8月12日）》，《中央苏区第五次反“围剿”》（上），第204页。}西运。同时，中共中央还秘密派出国家政治保卫局保卫大队到于都、登贤等红军预定集结地域，侦察路线。中革军委则由副总参谋长张云逸率领一个分队，潜往赣粤湘边界，侦察敌情和道路交通。


9月8日，中革军委电示第三军团指挥员彭德怀、杨尚昆，要求其在9月底前“阻止敌人于石城以北”，“在执行这一任务时应最高度的节用有生兵力及物质资材。在战斗的间隙中除三分之一的值班部队外，主力应集结补充、整理训练，并加强部队的政治团结”。同时强调在石城地区的防御战应进行“运动防御”，“不要准备石城的防御战斗，而应准备全部的撤退”。\footnote{《朱关于三军团应阻石城以北敌人以保卫瑞金部署的指示（1934年9月8日）》，《中央苏区第五次反“围剿”》（上），第216～217页。}14日，周恩来向林彪、聂荣臻等传达了中共中央和中革军委关于红军准备战略转移的决定。24日，中革军委再电彭、杨，重申：“为避免过多的损失及确实突击，第一地区阵地，应作为掩护地带，而第二地区阵地应作为主要抵抗地带。”\footnote{《中革军委致彭德怀、杨尚昆电（1934年9月24日）》。}显然，中共中央此时已将保存有生力量随时准备撤退作为主要考虑。中革军委发出命令，要求各军团在10月1日前组织好后方机关，加强运输队的建设，把敌人占领县区的军事部立刻改为县区游击队司令部和政治部，县区军事部长为游击队司令员、队长，县区委书记兼游击队政治委员；并规定“如在边区和中心区域有被敌人侵犯之可能时，则将军事部作上述改组”，为红军主力突围后苏区继续坚持作出组织上的初步安排。


与此同时，为减轻红军突围西进的阻力，中共中央决定与“围剿”中央苏区的国民党南路军总司令、广东地方实力派陈济棠接触，展开停火谈判。陈济棠出于自身生存需要的考虑，把红军在江西的存在作为其与南京中央间的一道屏障，因此，对“剿共”军事阳奉阴违，和中共间一直有信使往还。早在1933年11月，共产国际代表就报告：“在中央苏区，广州政府代表已开始进行停战谈判。”\footnote{《埃韦特给皮亚特尼茨基的报告（1933年11月22日）》，《共产国际、联共（布）与中国革命档案资料丛书》第13册，第625页。}此后，谈判断续进行。当红军开始准备撤离时，谈判进一步加紧。潘汉年回忆：“当我们的红军向中国西部推进时，广州的军阀们认为，如果蒋介石能消灭红军，那就会给他们造成很大的威胁。在他们看来，红军至今是南京军队和广州军队之间的屏障。我们曾派代表去进行谈判。第一次我们未能达成什么协议。第二次我们终究争取到广州同意进行谈判。”\footnote{《潘汉年在共产国际执行委员会书记处非常会议上的发言（1935年10月15日）》，《共产国际、联共（布）与中国革命档案资料丛书》第15册，第54页。}1934年“八一”节前，双方通过谈判达成停战协议，并设立联络电台。9月，朱德致信陈济棠，声明“红军粉碎五期进攻之决战，已决于10月间行之”，表示愿就停止双方作战、恢复贸易、政治开放、军事反蒋、代购军火等问题与粤军举行秘密谈判。\footnote{朱德：《关于抗日反蒋问题给陈济棠的信》，《朱德选集》，人民出版社，1983，第18页。}14日，博古向共产国际报告了双方的接触。共产国际对与粤方接触高度重视，指示中共在谈判中主要应提出代购军火和取消封锁、恢复贸易，不应附加其他过高条件，以免“丧失利用广州人和南京人之间矛盾的机会”。\footnote{《秦邦宪给共产国际执行委员会的电报（1934年9月14日）》、《共产国际执行委员会政治书记处政治委员会给中共中央的电报（1934年9月23日）》，《共产国际、联共（布）与中国革命档案资料丛书》第14册，第236、253页。}


中共中央的这一表态迅速得到粤方回应，10月6日，中共中央、中革军委派潘健行（潘汉年）、何长工为代表，同陈济棠的代表杨幼敏、黄质文等在寻乌进行会谈。双方经过数日反复协商，达成了就地停战、互通情报、解除封锁、互相通商和必要时可互相借道等5项协议。其中借道一条，言明红军有行动时事先将经过要点告诉陈济棠，陈部即后撤20公里让红军通过，红军保证不进入广东腹地。中共和粤方成立的这一协议，为红军的顺利突围转移准备了极为有利的条件。张闻天在《红色中华》发表的文章明确指出：“近来国民党内部的某些军阀，愿意同我们在反蒋方面进行某些条件的妥协，我们显然是不会拒绝利用这种机会的。”\footnote{张闻天：《一切为了保卫苏维埃》，《红色中华》第239期，1934年9月29日。}长征开始后，中革军委给红军各军团指挥官发出指示，告之：“现我方正与广东谈判，让出我军西进道路，敌方已有某种允诺。故当粤军自愿的撤退时，我军应勿追击及俘其官兵；但这仅限于当其自愿撤退时，并绝不能因此而消弱警觉性及经常的战斗准备。”\footnote{《军委关于我方正与广东谈判让出西进道路，如粤军自愿撤退我军应勿追击的指示（1934年10月26日）》，《红军长征档案史料选编》，学习出版社，1996，第18页。}事实上，长征初期，红军之所以能顺利实现转移，和粤方放开道路直接相关，蒋介石曾在日记中指称“粤陈纵匪祸国，何以见后世与天下”。\footnote{《蒋介石日记》，1934年10月30日。}


突围方针确定后，8、9月间，中共中央在宣传、组织上做了一系列的准备工作。8月18日，中央红军机关刊物《红星》发表周恩来撰写的社论，提出：“我们更要在远殖的行动中增加我们的兵力，采取更积极的行动，求得在运动战中消灭更多的白军，我们要坚决挺进到敌人的后方去，利用敌人的空虚，大大的开展游击运动……创造新的苏区，创造新的红军，更多的吸引敌人的部队调回后方，求得整个的战略部署的变动。我们要反对对敌人后方的恐慌观念……要在抗日先遣队胜利的开展之形势下，时刻准备着全部出动去与日本帝国主义作战！”\footnote{周恩来：《新的形势与新的胜利》，《红星》第60期，1934年8月20日。}这隐约透露出红军战略转移的信息。负责组织工作的李维汉回忆：“1934年7、8月间，博古把我找去，指着地图对我说：现在中央红军要转移了，到湘西洪江建立新的根据地。你到江西省委、粤赣省委去传达这个精神，让省委作好转移的准备，提出带走和留下的干部名单，报中央组织局……根据博古的嘱咐，我分别到江西省委、粤赣省委去传达。”\footnote{李维汉：《回忆长征》（上），中共党史资料出版社，1986，第343页。}刘建华回忆，1934年8月底9月初，“按照中央局的部署，党省委和团省委派我带了两个人，到茶梓和乱石做党、团二线工作。所谓二线工作，即一线是公开的，随时准备跟随部队行动；二线是秘密的，准备在中央红军转移后，留在当地坚持斗争”。\footnote{刘建华：《风雷激荡二十年》，中央文献出版社，1999，第66～67页。}9月4日，中革军委发出号令，要求：“发展更多的苏区于敌人背后，瓦解敌军，改变敌人的战略部署，把中央苏区革命先进的光荣事业扩大到全中国去，发展民族革命战争，开始与帝国主义直接作战，这是我们彻底粉碎敌人五次‘围剿’的基本方针。”\footnote{《中革军委为扩大红军的紧急动员的号令》，《红色中华》第228期，1934年9月4日。}


9月19日，中华苏维埃共和国人民执行委员会主席张闻天发出指示信，决定调整苏维埃机构，取消国民经济部、财政部、粮食部，成立财政经济委员会，“保卫局与裁判部可合并为肃反委员会”，并规定“在战争特别紧张的区域甚至苏维埃所有的各部都可以取消，而由个别同志直接负责去解决当前特别重要的战争问题”；“所有苏维埃机关中各种无用文件都应销毁”，“机关工作人员中的家属，应该另行安顿”，下级机关在“同上级领导机关脱离交通关系时，依然能够去进行工作”。这实际是在布置苏区失陷后的工作。29日，张闻天又发表《一切为了保卫苏维埃》一文，指出：“为了保卫苏区，粉碎五次‘围剿’，我们在苏区内部求得同敌人的主力决战，然而为了同样的目的，我们分出我们主力的一部分深入敌人的远后方，在那里发动广大的群众斗争，开展游击战争，解除敌人的武装，创造新的红军主力与新的苏区……我们有时在敌人优势兵力的压迫之下，不能不暂时的放弃某些苏区与城市，缩短战线，集中力量，求得战术上的优势，以争取决战的胜利。”\footnote{张闻天：《一切为了保卫苏维埃》，《红色中华》第239期，1934年9月29日。}文章发表于《红色中华》，是中央红军准备实行战略转移的第一个公开信号。


在中共中央积极准备撤离时，9月17日，博古致电共产国际，提出：“中央和革命军事委员会根据我们的总计划决定从10月初集中主要力量在江西的西南部对广东的力量实施进攻战役。最终目的是向湖南南部和湘桂两省的边境地区撤退。全部准备工作将于10月1日前完成，我们的力量将在这之前转移并部署在计划实施战役的地方。”\footnote{《秦邦宪给共产国际执行委员会的电报（1934年9月17日）》，《共产国际、联共（布）与中国革命档案资料丛书》第14册，第251页。}9月30日，共产国际致电中共中央，表示：“考虑到这样一个情况，即今后只在江西进行防御战是不可能取得对南京军队的决定性胜利的，我们同意你们将主力调往湖南的计划。”\footnote{《共产国际执行委员会政治书记处政治委员会给中共中央的电报（1934年9月30日）》，《共产国际、联共（布）与中国革命档案资料丛书》第14册，第256页。}这意味着中共的转移计划已得到共产国际的完全同意。10月8日，中共中央向新成立的苏区中央分局发出训令，提出在国民党军不断深入苏区的形势下，如果红军主力继续在缩小的苏区内部作战，“因为地域上的狭窄，使红军行动与供给补充上感觉困难，而损失我们最宝贵的有生力量，并且这也不是保卫苏区的有效的方法。因此，正确的反对敌人的战斗与澈底粉碎敌人五次‘围剿’，必须使红军主力突破敌人的封锁，深入敌人的后面去进攻敌人。这种战斗的方式似乎是退却的，但是却正相反，这才是进攻敌人，克服敌人堡垒主义，以取得胜利的重要方式。因为这样的行动，将在离开堡垒的地区中得到许多消灭敌人的战斗机会，解除敌人的武装壮大红军，在广大的新的区域中，散布苏维埃影响，创立新的苏区”。\footnote{《中国共产党中央委员会给中央分局的训令》，《红军长征档案史料选编》，第3～4页。}


在中央红军加紧突围准备之时，国民党“围剿”军主力尚在休整。此时，国民党军久攻之下，师劳兵疲，且因长宿野外，患病者甚多；另外，蒋介石本人对战局估计十分乐观，9月6日，蒋在日记中预定计划：“一、进剿至石城宁都与长汀之线，当可告一段落，以后即用少数部队迫近，与飞机轰炸当可了事。二、用政治方法招降收编，无妨乎。”\footnote{《蒋介石日记》，1934年9月6日。}正是在这种情绪影响下，国民党军在9月底才开始新一轮进攻，9月26日其第三、第十、第五纵队共6个师向石城攻击，第八纵队向兴国攻击，第七纵队向古龙岗攻击，第四纵队向长汀攻击，南路军向会昌进攻，并拟于10月14日总攻瑞金、宁都。


9月25日，国民党军主攻部队第三路军下达攻击命令，以小松市为主攻方向。26日，国民党军6个师全线展开攻击，先后攻占中华台、陈古岭、分水坳等高地，30日占领小松市。第十纵队指挥官汤恩伯亲至第一线观察，发现石城城北石榴花、鹅项坳一带工事密布，是红军主力集中地区。10月3日，国民党军第八十八师在专门调集的飞机、大炮掩护下，向石榴花、鹅项坳高地猛烈攻击。两高地是石城最后屏障，红军在此进行了顽强固守，但难以抵挡对方的猛烈炮火，被迫撤出战斗。7日，国民党军第十一师进占石城。红军顽强奋战，在石城阻挡了国民党军的进攻步伐，为准备战略转移争取了时间，自己也付出很大伤亡，战役结束后，“红三军团老的连长完全死伤”。\footnote{陈云：《遵义政治局扩大会议传达提纲》，《遵义会议文献》，第39页。}肖华回忆：“这一仗打得很壮烈，损失也很大。连续战斗的伤亡，原一万多人的‘少共国际师’到这时只剩下五千来人了。”\footnote{肖华：《中国青年的灿烂花朵——回忆“少共国际师”》，《忆少先队和少共国际师》，江西人民出版社，1979，第6～7页。}


与石城激战同时，向古龙岗进攻的国民党军第六路军遭到红军顽强阻击，前进一度受阻；后由于红军主力部队的撤出，行动加速，先后占领天子嵊、风车坳等高地，10月10日完成对古龙岗的占领。兴国方面，红一军团和红五军团在庙背以北和高兴圩西南地区顽强抗击周浑元纵队的进攻。战至9月30日，庙背、高兴圩等地先后失守，红一、五军团被迫“向兴国退去”，\footnote{《顾祝同电蒋中正（1934年9月30日）》，蒋中正文物档案002090300081101。}撤至新圩、文陂地区，继续抗击国民党军对兴国的攻击。


在前方进行阻击战时，红军转移已箭在弦上，主力部队纷纷撤往后方集中，只有地方部队继续留在前方阻击敌人。9月21日，中革军委决定将红二十一、红二十三师合编为第八军团，周昆任军团长，红八军团司令部由第二十一师司令部代理。此外，9月15日还成立了教导师，张经武任师长，何长工任政治委员。10月2日，博古亲至红三军团，召集团以上干部会，彭绍辉日记记载，彭德怀在会上指出：“我军在北线迟滞敌人，争取时间的任务已完成。我军要向敌人反攻，主力须转移。”\footnote{《彭绍辉日记》，1934年10月2日，第36页。}这已经非常明确地传达出红军即将进行战略转移的信息，一些回忆录所说红军许多中、高级指挥员一直到长征开始仍不知将要进行战略转移的说法是站不住脚的。7日，中共中央和中革军委命令红二十四师和地方武装接替中央红军主力的防御任务，主力集中瑞金、于都地区，准备执行新的任务。9日，红军总政治部发布《关于准备长途行军与战斗的政治指令》，要求“加强部队的政治军事训练，发扬部队的攻击精神，准备突破敌人的封锁线，进行长途行军与战斗”。\footnote{《总政治部关于准备长途行军与战斗的政治指令》，《红军长征档案史料选编》，第9～11页。}


应该说，长征出发前，中共中央为长征进行的物质、舆论、组织等各方面的准备是较为充分的。在武器弹药、粮款筹集、兵员发展上都做了大量的工作。有关长征初期的回忆文章写道：“我们的帽子、衣服、布草鞋、绑带、皮带，从头到脚，都是崭新的新东西。”\footnote{艾平：《第六个夜晚》，《中国工农红军第一方面军长征记》，人民出版社，1958，第61页。}耿飚回忆，出发前，“种种迹象表明，红军要有大的行动。师部不断通知我们去领棉衣，领银元，领弹药，住院的轻伤员都提前归队，而重伤员和病号，则被安排到群众家里。地图也换了新的，我一看，不是往常的作战区域，这说明，部队要向新的地域开进”。\footnote{《耿飚回忆录》（上），第169页。}李一氓回忆，他长征前几天赶回瑞金时，“看到别人都有了充分的物质准备，因为他们早已得到一路出发的正式通知”，“有些人有新的胶底帆布鞋，有些人不知在什么地方搞来的很不坏的雨衣，有的人还有很好的水壶，很好的饭盒，很新的油纸雨伞，五节的大电筒”。\footnote{《李一氓回忆录》，第161页。}可见，长征前夕对于要进行战略转移这一点，已经传布到相当范围。就此而言，人们没有理由忽视陈云当时的说法：“此次红军抛弃数年经营之闽赣区域而走入四川，显系有计划之行动。当去年退出江西以前，以我之目光观之，则红军已进行了充分准备。”\footnote{廉臣（陈云）：《随军西行见闻录》，《中国工农红军第一方面军长征记》，第5页。}1939年，当李德在共产国际遭到批评并接受了大部分指责时，仍然特别就长征问题作出辩解，强调：“在技术方面，我认为，远征的准备工作是好的，突破四道防线的计划也一样，比较容易的克服这些防线就证明了这一点。”\footnote{《布劳恩给共产国际执行委员会的书面报告（1939年9月22日）》，《共产国际、联共（布）与中国革命档案资料丛书》第15册，第349页。}


一些论者及回忆录提到红军开始长征时，没有对进军方向及进军计划作出交代，对这些质疑，董必武当年的回答应有借鉴意义：

\begin{quote}

主力转移自然是由西向北前进，这是毫无疑问的。至于转移到什么地方，经过什么路线，走多少时候等问题，系军事上的秘密，不应猜测，而且有些问题要临时才能决定。如行军走哪条路，什么时候到达什么地方，有时定下了，还没有照着做，或做了一部分，忽因情况变了又有更改，这是在行军中经常遇到的，只要大的方向知道了，其余的也就可以不问。\footnote{董必武：《长征纪事》，《董必武选集》，人民出版社，1985，第16页。}
\end{quote}


对于一场军事行动而言，必要的保密应该属于常识，将近一个月的准备动员事实上也是当时可能有的时间极限。


为便于随军行动，中共中央、中央政府、中革军委机关和直属部队编为两个纵队。第一野战纵队由红军总部和干部团组成，叶剑英任纵队司令员兼政治委员，下辖4个梯队，博古、李德、周恩来、朱德等随该纵队行动。第二野战纵队由中共中央机关、苏维埃中央政府机关、后勤部队、卫生部门、总工会、青年团等组成，罗迈（李维汉）任司令员兼政治委员，毛泽东、张闻天、王稼祥等随该纵队行动。10月8日，中革军委公布各军团组成情况：一军团17280人，三军团15205人，五军团10868人，八军团9022人，九军团10238人，第一野战纵队4893人，第二野战纵队9853人，共7.7万余人。同时，计划给红一、三、五、八、九5个军团补充9700人。13日，中革军委下令将补充团人员拨付给各野战军团，\footnote{《朱德关于各补充团正式拨给各兵团管辖并负责训练的指令（1934年10月13日）》，《红军长征·文献》，解放军出版社，1995，第66页。}包括中央两纵队在内的野战军人数达到8.6万余人。


中央红军主力转移前，中共中央决定由项英、瞿秋白、陈毅、陈潭秋、贺昌等组成中共苏区中央分局、中央军区和中华苏维埃共和国中央政府办事处，项英为中央分局书记、中央军区司令员兼政治委员，陈毅为办事处主任，统一领导中央苏区和闽浙赣苏区的红军和地方武装。


10月10日晚，中央红军开始实行战略转移。中共中央、中革军委率领第一、第二野战纵队，分别由瑞金县的田心、梅境地区出发，向集结地域开进。16日，中央红军各部队在于都河以北地区集结完毕。中央红军的战略大转移由此开始。关于这一次战略大转移，李德有一个说法是：“蒋介石的第四和第五次‘围剿’的经验说明，在三十年代中期新的国际环境和民族状况中，比较小的和互相隔绝的苏区是不能长期坚持下去的。”\footnote{〔德〕奥托·布劳恩：《中国纪事1932～1939》，第113页。}

\subsection{红军转移中蒋介石的对策}


随着红军的转移，中共历史上的一个时代实际即宣告结束。作为尾声，值得一提的是红军开始战略转移前后蒋介石的态度，事实上，这和整个中央苏区的发展进程仍然有逻辑上的一致性。


红军转移开始之初，进展十分迅速。湘江战役前的两个月多一点时间内，红军从赣南西部转移到广西境内，行军3500里，且几乎没有遭遇大的战斗。以红一军团为例，10月第一个月日行军里程基本在60～90里之间，共行军11天，计860里，平均每日78.2里。11月行军24天，计1530里，平均每日63.75里。\footnote{《红军第一军团长征中经过的地点及里程一览表》，《中国工农红军第一方面军长征记》，第423～425页。遵义会议后的1935年2月红军由于需要摆脱国民党军追兵，28天内有26天在行军，计1620里，平均每日62.3里，3月行军24天，计1535里，平均64里。可以看出，这时的行军速度和1934年11月相仿，慢于1934年10月。}考虑到转移人员多达10余万人，进行的又是超远距离连续行军，此种速度应称快捷。共产国际对此的评论是：“运动的目的——使自己的部队与四川红军兵团部队会合——几乎在没有来自敌人方面任何干扰的情况下，完全实现了。”\footnote{《共产国际执行委员会东方书记处关于中国军事形势的通报材料（1935年2月11日）》，《共产国际、联共（布）与中国革命档案资料丛书》第14册，第373页。}在1935年初，这一判断应该说稍显乐观，但其对红军转移初期状况的描述并不夸张。无独有偶，当时供职于国民党中央党部的王子壮也在日记中写道：“赤匪主力并未消灭，以不堪中央军之压迫相率西去，经赣州、大余等地而趋湘南，盖欲入川以会合徐向前之股匪也。复以何键为追剿总司令，然以何之兵不善战，恐难收阻止之效，据报告已骎骎西进。”\footnote{《王子壮日记》第2册，第165页。}“骎骎西进”，的确是红军初期快速转移的最好注脚。


不过，应该特别强调的是，红军初期转移的顺利，并不能如李德在前述答辩中用来反证中共中央准备的成功，虽然中共中央的准备工作不像曾被指责的那样仓促、零乱，尚属中规中矩，但也并无出奇之举。要实现初期堪称快捷的转移速度，远非红军单方面所可决定，当红军实际处于被动撤退这样一种境地，向着一个没有群众基础的地域挺进时，更能够在实质上决定红军命运的，还是其对手方的动向。红军和粤方的谈判成功，使其在长征之初事实上为中共开放了西进道路自是重要原因，同时，蒋介石的态度其实也十分复杂。李宗仁就曾谈道：“就战略的原则来说，中央自应四方筑碉，重重围困，庶几使共军逃窜无路，整个就地消灭……但此次中央的战略部署却将缺口开向西南，压迫共军西窜。”\footnote{《李宗仁回忆录》（下），广西政协文史资料委员会编印，1981，第625～626页。}这样的说法虽属一家之言，但并非无稽之谈，事实上，正如毛泽东在论述红色政权为什么能够存在时强调的政治割据因素一样，长征初期的进程和国内政治力量间之诡谲互动有着无法忽视的关系。


对于蒋介石在“围剿”最后阶段的作战可能的态度，共产国际代表曾有精当的预测：“蒋介石需要这样来同红军作战，使他在消灭红军之后不是太被削弱地出现在福建和广东的边界上。如果做不到这一点，那么同红军作战就不能达到他在军事上和政治上得到加强的目的。”在共产国际代表看来，蒋的碉堡战术既是其对付红军运动战的手段，还包含着他观测形势发展的战略考虑：“蒋介石比德国人更了解江西和在它之后的省份（福建、广东、湖南），因此他知道，突破不会给他带来太多的东西。他认为，根据各种预测，突破可以取得成功，但为此他要付出昂贵的代价。此外，红军最终是会被打散的，那时斗争就永不会结束。因此蒋介石想把红军置于包围圈内，看他们在碉堡包围情况下怎么办。”\footnote{《施特恩关于中国军事政治形势的报告（1933年11月8日）》，《共产国际、联共（布）与中国革命档案资料丛书》第13册，第600、602页。}


第五次“围剿”最后阶段，南京方面握有充分的主动权，正因此，素喜在内斗中纵横捭阖的蒋介石，对“围剿”结局没有放弃开放各种可能性。抗战前夕，蒋与中共谈判期间，曾在日记中写道：“政治胜利至七八分为止，必须让敌人一条出路，此围攻必缺之道乎。”\footnote{《蒋介石日记》，1937年6月3日。}这段话，相当程度上传达出蒋进行政争的内在理路和处事方式，对观察蒋的种种作为不无参考意义，其对苏区封锁及进攻的部署，隐约可以窥测到他的如上复杂思虑。“围剿”期间，其在东南北三面展开对红军的围困，西面赣州方向力量始终空虚，没有布置有力部队，这固与赣州一带临近广东、湖南，牵涉到政治地理的复杂形态相关，但蒋的无所作为仍让人印象深刻。随着“围剿”的深入，西线赣州方向地位日形重要，1934年8月，国民党“围剿”军西路军总司令何键明确谈道，在合围完成状况下，红军一定会寻找出路，“如窜粤窜浙，均为事实上所不可能，一定只有下述两条路，一为窜闽，一为从粤湘接壤的地方窜入湖南，再图窜别的地方”。\footnote{《何总司令在扩大纪念周中报告最近剿匪情况与修建飞机场的必要》，《西路军公报》第14期，1934年8月15日。}9月底，何键更进一步向蒋报告，据其得到的情报，中共准备西进与湘黔部队合股，并建议“抽调三师驻扎赣州，原驻部队及南路军集合赣州、信丰、安远间，以防匪窜越赣水”。\footnote{《贺国光电蒋中正转何键报上海伪中央会议决议放弃闽南老巢绕窜湘桂边境盘据川滇为其大本营现判断赣匪企图西窜其与川黔之匪合股现建议抽调三师驻扎赣州而将原驻部队及南路军集合于赣州信丰安远间期防匪窜越赣水等语（1934年9月29日）》，蒋中正文物档案002090300086050。}10月初，白崇禧也提醒蒋加强西路兵力，“宜驻良口、赣州、南康”，\footnote{《白崇禧电蒋中正何键据廖磊报称王家烈遭共匪击退沿乌江退往思南且贺龙意图与萧克合股现拟各路于赣南闽西会剿其中西路军宜驻良口赣州南康而南路军驻守筠门岭安远信丰之线等以阻匪西窜（1934年10月5日）》，蒋中正文物档案002090300096129。}以防止中共主力西进。但是，何、白这些急如星火的建议在蒋这里并未得到认真回应，除象征性要求加强构筑碉堡工事外，缺乏有效的实际动作。以致红军开始西进时，赣州一带防御相当空虚，“南康县周百里无国军”，而赣州商会则向蒋介石报告：“沿江两岸数百里均无兵驻”，形成“西路军之部署难以防堵”\footnote{《南康县商会等电蒋中正熊式辉南康县百里无国军请即派军驰镇（1934年10月26日）》、《赣县党部暨商会电蒋中正股匪西窜驻军集剿城内与沿江两岸空虚请速令罗师兼程前进并盼示覆（1934年10月25日）》、《黄克鼎电蒋中正现已完成砻市经大龙大舸大坳背之碉堡又赣匪西窜西路军之部署难以防堵请派兵控制遂上以备截击而以罗霖师等部队向桂汝增移（1934年10月28日）》，蒋中正文物档案002090300082173、002090300082174、002090300084042。}之局。


不仅是何键等，早在1934年6月，陈诚曾针对红军可能的突围提出建议：“我军部署应仍以有力一部深入匪巢，诱导其它各路前进，以达合围之目的。另以一部准备匪如西窜随时均可截击。故拟（一）以我三、五两路军主力之罗（卓英）、樊（崧甫）、汤（恩伯）三纵队，会取石城，威胁瑞金，匪必东援伪都，而忽视西面。此时即令周（浑元）纵队乘机袭取兴国，诱导南路军北取雩都。（二）第六路军就现地停止集结主力，策应各方，待机前进。（三）兴国附近地势险恶，道路崎岖，匪化最深，仅以现有之周纵队兵力进取兴国似嫌单薄，故拟请增调九十九师归周指挥，以厚实力。”\footnote{《电呈依据匪情拟进剿方略乞钧裁（1934年6月14日）》，《陈诚先生书信集·与蒋中正先生往来函电》（上），第136～137页。}单从战场形势看，陈的建议切中要害，应属高明。对此，蒋介石的反应同样出人意料，他不仅没有按陈诚所说在西线增加兵力，反而在红军行将开始转移时指示：“欲促进战局之从早结束，则东路应增加兵力，如能将第四与八十九两师由汤（恩伯）带领东移，则东路即可单独向长汀、瑞金进展，一面北路军占领宁都与薛（岳）路会合后，即可由宁都与东路军由长汀会占瑞金，可免石城与长汀线之兵力与时间也。”\footnote{《蒋中正电陈诚战局欲早结束东路应增加兵力（1934年10月9日）》，蒋中正文物档案02020002000064。}难明蒋介石机心的陈诚回电表达出自己的疑惑：“为求歼匪于赣南计，我军重点应偏于西翼地区，使东路军不必急进，免迫匪西窜，如为减少部队调动之疲劳及免费时间计，汤纵队似可不必东调。”\footnote{《陈诚电蒋中正江西赤匪欲牵制我主力以西窜及我东南西北各路部队进剿部署情形（1934年10月15日）》，蒋中正文物档案002090300090032。}陈诚的质疑，恰恰反证出蒋介石西兵东调的用心，即将红军往西线逼迫，其战略目的极有可能指向的是逼走而不是围歼红军。晏道刚回忆，在得知红军离开赣南的消息时，蒋介石曾表白：“不问共军是南下或西行、北进，只要他们离开江西，就除去我心腹之患”；“红军不论走那一条路，久困之师经不起长途消耗，只要我们迫堵及时，将士用命，政治配合得好，消灭共军的时机已到，大家要好好策划”。\footnote{晏道刚：《蒋介石追堵长征红军的部署及其失败》，《红军长征在贵州·史料选辑》，《贵州社会科学》编辑部、贵州省博物馆编印，1983，第347页。}


“围剿”最后阶段蒋的真实态度，前引9月6日蒋日记表明其此时尚少对红军穷追猛打的腹案，倒是有一厢情愿的招安设想。蒋日记还显示，此时他对粤方和中共的接触已有掌握，10月17日，日记“注意”栏第一条就是：“粤陈联匪已成乎，应制止之。”\footnote{《蒋介石日记》，1934年10月17日。}从语气上看，他对粤、共接触应早有了解，情报掌握尚称迅捷。不过，蒋所谓制止并无下文，倒是此前几天，10月14日，他还在催促粤方向北开动，要求其“进取于都，以为我封锁线前进部队之据点”。如果粤方果真向北开动，充当堵截红军西进的主力，蒋自然乐于见到，但衡诸宁、粤、共三方微妙关系，粤方出动概率几乎为零，而蒋也不会天真到对此真的寄予期待，所以对于粤方提出的索取津贴要求，他一改往常的慷慨态度，以“待其确实占领于都后照办”，将皮球踢了回去。\footnote{《蒋中正总统档案·事略稿本》第28册，台北，“国史馆”，2007，第275页。与此相对照，蒋在军情紧急时，擅打银弹攻势，不等部队开拔就发放开拔费并不鲜见。参见金以林《从反叛到瓦解——石友三1931年反蒋失败的个案考察》，《近代中国》第156期，2004年3月。}


蒋介石的所谓“放水”\footnote{蒋纬国晚年在口述自传中谈及这段历史时称：“从整体来看，当时与其说是没有包围成功而被中共突围，不如说是我们放水。”见刘凤翰整理《蒋纬国口述自传》，中国大百科全书出版社，2008，第6页。}举动，当年窥破奥妙者或已有人在。11月底，在政坛翻滚有年的何其巩上书蒋介石，提出：“赣匪之可虑，不在其窜逃，而在其守险负隅，旷日持久……赣匪倘能在赣川以东，合围而聚歼之，固为上策，否则有计划的网开一面，迫其出窜，然后在追剿中予以节节之击灭，似亦不失为上策中之中策也。”这或许是何揣测蒋介石心态的一种进言。何并针对西南三省谈道：“川滇黔三省，拥有七千万以上之人口，形险而地腴，煤盐油矿以及各种金属皆不缺乏，足为国防之最后支撑点。宜乘徐匪猖獗之时，或在赣匪西窜之时，力加经营。即钧座不能亲往，亦宜派遣忠义大员，统率重兵入川。”对何的上述看法，蒋批曰：“卓见甚是，当存参考。”\footnote{《何其巩呈蒋中正抗日救国根本大计在于安定北方巩固中部经营西南之意见（1934年11月25日）》，蒋中正文物档案002080200194043。}以此为基础，1934年底除旧迎新之际，蒋瞻前顾后，在日记中将“追剿”红军、抗日准备与控制西南三者巧妙结合：“若为对倭计，以剿匪为掩护抗日之原则言之，避免内战，使倭无隙可乘，并可得众同情，乃仍以亲剿川、黔残匪以为经营西南根据地之张本，亦未始非策也。当再熟筹之！”\footnote{《蒋介石日记》，1934年12月29日。}这段话，和前述种种结合看，的确意味深长。


从国民党方面作战计划看，在红军积极准备撤离的关键阶段，其最高指挥机关并没有实施围歼的严密计划。1934年9月29日，时任南昌行营参谋长的贺国光向蒋报告：“综合各方情报判断，赣匪似有西窜企图，如能在沿赣江至信丰、安远之第一线及宁冈、桂东、汝城、仁北、曲江之第二线以东地区歼灭之固善，倘被窜逸，则以在沿湘江桂江之第三线封锁较为确实。李白何迭次来电同此主张。诚以此线利用河流天然形势南北连贯不易绕越防守极易。除湘江段已颁发西路补助金赶筑外，广西段据李白来电亦在进行中。”\footnote{《贺国光电蒋中正据各方情报判断江西共军似有西窜企图请准拨款修筑湘赣桂边区碉堡封锁线并派马吉第或聂洸前前往督察工程并与桂方接洽（1934年9月29日）》，蒋中正文物档案002090300083152。}可见国民党军一开始就将围堵红军重点放在了湘江一线。\footnote{1934年10月31日，当中央红军还在湘赣边境时，蒋介石给前方陈济棠、何键、顾祝同的电令中就要求将其“聚歼于湘江以东地区”，也可作为正文的佐证。参见《蒋中正电各高级将领赣南共军已由赣州大庾间逃窜嘱各将士截追（1934年10月31日）》，蒋中正文物档案002020200029001。}更有意思的是，在红军主力已离开赣南但尚未远飏时，蒋甚至告诉部下：“匪主力未窜过汝城桂东以西之前，我三路不必急取瑞金，而可先取雩都也。”\footnote{《蒋中正电熊式辉贺国光共匪主力未过汝城桂东以西前因东路军可由长汀相机进取瑞金故三路不必急进而可先取雩都（1934年11月3日）》，蒋中正文物档案002090300096148。}蒋此举显然不能单从军事角度加以解释，其中含有很深的用心：既与其当时策动前方将领以“剿共”尚未完成呼吁推迟召开国民党五全大会相关，也不无担心占领瑞金刺激红军回返而暂时控兵不进的意图。


看到这些，我们不能不钦佩第五次反“围剿”开始之初共产国际驻华军事总顾问弗雷德对蒋动向的分析：“要么他应该率领强大的部队由福建和赣江上游（向赣州方向）推进，出现在红军和闽赣两省军队之间，要么指望红军向赣西和湖南突破，停止在原地没有希望的斗争。这也会把他拖到湖南，他会在那里面临准备新战役的任务。”\footnote{《施特恩关于中国军事政治形势的报告（1933年11月8日）》，《共产国际、联共（布）与中国革命档案资料丛书》第13册，第600～601页。}有理由相信，蒋当时选择了后者。正由于此，10月中旬，当红军正紧锣密鼓地进行着长征的最后准备时，蒋介石却从江西前线撒手西去，开始其被当时报章称为“万里长征”\footnote{《蒋委员长出巡纪要》（下），《军政旬刊》第39期，1934年11月10日，第118页。}的一个多月西北、华北之行。在蒋看来，红军的西去，也许就像围棋里对手逃大龙一样，只要自身立于主动，大可顺水推舟、就势而为，所谓“赣匪一旦窜遁，则无论跟踪追剿之师，因地留戍之师，回防中部北部之师，控制西南一带之师，皆能左右逢源，不虞粘滞，从此大局可期永安”。\footnote{《何其巩呈蒋中正抗日救国根本大计在于安定北方巩固中部经营西南之意见（1934年11月25日）》，蒋中正文物档案002080200194043。}而此时蒋的西北、华北之行，当然不会和对日问题没有关系，也不会不落入背靠西北、正和日本相互提防的苏俄眼中。考虑到当时蒋介石正和苏俄寻求更紧密关系，西北之行传达的意思，对中苏、中日乃至国共关系，都透着历史深处一言难尽的微妙。


1927年后的苏维埃运动，中共凭借其真诚的信仰、严密的组织、强大的动员、坚固的武力，最大限度地发挥出共产革命的威力。有了中共本身的实力，加上风云际会的国内外政治态势，革命的烈火才会以其可能有的最猛烈态势熊熊燃烧。在此过程中，中共革命不断挑战着能与不能的边际，创造着属于自己的辉煌。


不过，革命的张力不可能无限制地伸展，夺取政权是革命的既定目标，但当年这样的目标事实上还难以企及。因此，当国民党军对苏区集中全力展开进攻时，无论是鄂豫皖、湘鄂西、湘鄂赣、闽浙赣，最终都无可挽回地走向了失败。当年国共之间的对垒，并不完全在同一个数量级内进行，中共的发展，更多的是利用国民党统治的内部冲突，当这种冲突趋于平稳、南京政府力量不断上升时，中共受到的压力将空前增大。在一个丛林法则的世界中，无论是历史还是现实，实力终究是进退成败的关键。第五次“围剿”和反“围剿”的发端、进行乃至最终结局，如果放到这一大背景下衡量，将可以有一个更为持平的理解。


中国革命是一个持续推进的过程，在这期间，每一代革命者都付出了牺牲，也有他们的困难和局限。无论是个体还是群体，革命者们在创造历史的同时，不可避免地受着时势的制约：国内的、国外的、社会的、政治的、经济的。中央苏区后期，诸多的不利因素共同挤迫着中国革命，这些当年为了生存和理想揭竿而起、意气风发中还略带青涩的革命者们，终究敌不过占据着天时地利又正处盛年的南京政权，革命，将在撤离中央苏区的路途上，重新出发。

